% Lesson 21 — Vocabulary: Words and expressions related to helping with fund-raising activities/support groups in the acquisition of community resources/services — Anglais Tle A — Cours signé OnBuch+
% Programme officiel MINESEC. Compiler avec : tectonic EN21-vocabulary-words-and-expressions-related-to-helping-wit.tex
\newcommand{\DOCMATIERE}{Anglais}
\newcommand{\DOCNIVEAU}{Tle A}
\newcommand{\DOCMODULE}{Chapter 2 : Economic Life and Occupations}
\newcommand{\DOCLECON}{EN21}
\newcommand{\DOCTITRE}{Lesson 21 — Vocabulary: Words and expressions related to helping with fund-raising activities/support groups in the acquisition of community resources/services}
% OnBuch+ — préambule « cours signés » ENRICHI (compilé avec tectonic / XeLaTeX)
% Système visuel « Pop » d'OnBuch+ : fond crème (jamais blanc), bordures encre
% épaisses, ombres « dures » décalées, titres Archivo Black en capitales.
% Version MATHÉMATIQUES : graphes (pgfplots/tikz), boîtes pédagogiques
% (définition, propriété, méthode, attention, le savais-tu, exercice résolu,
% exercice, TP, bilan), page de garde et signature OnBuch+ ancrée en bas.
%
% Macros attendues dans main.tex AVANT \input{preamble} :
%   \DOCMATIERE  \DOCNIVEAU  \DOCMODULE  \DOCLECON (numéro)  \DOCTITRE
\documentclass[11pt]{article}

\usepackage{fontspec}
\usepackage[english,french]{babel} % français = langue principale ; l'anglais via \eng{..} / textebox
\usepackage[a4paper,margin=2cm,top=3.4cm,bottom=2.8cm,headheight=1.4cm,footskip=1.3cm]{geometry}
\usepackage{amsmath,amssymb,amsfonts}
\usepackage{mathtools} % \xmapsto, \coloneqq... souvent utilisés par les modèles
\usepackage{cancel} % simplifications barrées (bilans de forces)
% \abs{z} (module, valeur absolue) et \norm{u} (norme), utilisés par les modèles
\providecommand{\abs}{}\let\abs\relax\DeclarePairedDelimiter{\abs}{\lvert}{\rvert}
\providecommand{\norm}{}\let\norm\relax\DeclarePairedDelimiter{\norm}{\lVert}{\rVert}
\usepackage[table]{xcolor} % [table] : fournit aussi \rowcolors (alternance de lignes)
\usepackage{pagecolor}
\usepackage{tikz}
\usetikzlibrary{arrows.meta,positioning,calc,shapes.geometric,shapes.misc,shapes.arrows,decorations.pathmorphing,decorations.pathreplacing,decorations.markings,patterns,shadows,mindmap,trees,fit,backgrounds,matrix,intersections,angles,quotes}
\usepackage{pgfplots}
\usepackage[version=4]{mhchem} % équations nucléaires \ce{^{235}_{92}U}
\usepackage[european,siunitx]{circuitikz} % schémas électriques
\pgfplotsset{compat=1.18}
\usepgfplotslibrary{fillbetween,groupplots} % aires entre courbes (calcul intégral)
\usepackage{enumitem}
\usepackage{fancyhdr}
% [most] charge toutes les bibliothèques tcolorbox (listings, minted, raster,
% documentation...) et gonfle inutilement la mémoire TeX sur un document long
% (~300 boîtes) : on ne charge que ce qui est réellement utilisé.
\usepackage[skins,breakable]{tcolorbox}
\usepackage{graphicx}
\usepackage{colortbl}
\usepackage{booktabs}
\usepackage{tabularx}
\usepackage{diagbox} % case d'en-tête diagonale des tableaux à double entrée
% Colonnes extensibles centrée / gauche / droite (C, L, R), souvent utilisées par les modèles
\newcolumntype{C}{>{\centering\arraybackslash}X}
\newcolumntype{L}{>{\raggedright\arraybackslash}X}
\newcolumntype{R}{>{\raggedleft\arraybackslash}X}
\usepackage{array}
\usepackage{multicol}
\usepackage{multirow}
\usepackage{siunitx}
% parse-numbers=false : accepte aussi \SI{2,5 \times 10^{-3}}{...}
\sisetup{locale=FR,output-decimal-marker={,},group-minimum-digits=4,parse-numbers=false}
\DeclareSIUnit\ohm{\ensuremath{\symup{\Omega}}} % U+2126 absent de la police
\DeclareSIUnit\division{div} % réglages d'oscilloscope (ms/div, V/div)
\DeclareSIUnit\tr{tr} % tours (vitesse de rotation, ex. \SI{1500}{\tr\per\minute})
\DeclareSIUnit\molar{M} % concentration molaire (ex. \SI{3}{\molar}, \SI{1}{\micro\molar})
\DeclareSIUnit\an{an} % année (le modèle confond parfois avec \year, un compteur TeX) — ex. \SI{5}{\kilo\meter\per\an}
\DeclareSIUnit\FCFA{FCFA} % franc CFA (ex. \SI{500}{\FCFA\per\kilo\gram})
\DeclareSIUnit\degreeBrix{\textdegree Brix} % teneur en sucre d'un sirop/jus (réfractométrie)
\DeclareSIUnit\month{mois} % le modèle utilise parfois \month, un compteur TeX, comme unité de durée
\usepackage{qrcode}
% \degree : commande fréquemment utilisée par les modèles pour un angle ou une
% température en dehors de \SI{}{} (non fournie par les paquets chargés).
\providecommand{\degree}{^{\circ}}
% \qed (fin de démonstration), utilisé par les modèles sans amsthm
\providecommand{\qed}{\hfill$\square$}
% \barre : double barre des valeurs interdites dans un tableau de variations (tkz-tab)
\providecommand{\barre}{\ensuremath{\|}}
% environnement proof (amsthm non chargé) utilisé par les modèles
\ifdefined\proof\else\newenvironment{proof}[1][Démonstration]{\par\noindent\textit{#1.}\ }{\qed\par}\fi
\usepackage{lastpage}
\usepackage{hyperref}
\hypersetup{hidelinks}

% ============================================================
% Palette « Pop » OnBuch+ — mêmes hex que la classe P (plus_ui.dart)
% ============================================================
\definecolor{poporange}{HTML}{F59321}
\definecolor{popdark}{HTML}{E07A0C}
\definecolor{popcream}{HTML}{FFF6E8}
\definecolor{popink}{HTML}{1C1714}
\definecolor{poppurple}{HTML}{7A5AE0}
\definecolor{popgreen}{HTML}{2E9E6A}
\definecolor{poppink}{HTML}{E0457B}
\definecolor{popblue}{HTML}{2F7DE1}
\definecolor{popgold}{HTML}{C98A12}
\definecolor{popmuted}{HTML}{8A7F76}
\definecolor{popline}{HTML}{EADFD2}
\definecolor{popgray}{HTML}{8A7F76}
\colorlet{popgrey}{popgray}
% Alias tolérants (le modèle nomme parfois les couleurs différemment)
\colorlet{popwhite}{white}
\colorlet{popblack}{popink}
\colorlet{popred}{poppink}
\colorlet{popyellow}{popgold}
% Teintes claires (fonds de tableaux/encadrés) — le modèle nomme parfois
% les couleurs avec un suffixe L (ex. popblueL) sans les avoir définies.
\colorlet{poporangeL}{poporange!20}
\colorlet{popdarkL}{popdark!20}
\colorlet{poppurpleL}{poppurple!20}
\colorlet{popgreenL}{popgreen!20}
\colorlet{poppinkL}{poppink!20}
\colorlet{popblueL}{popblue!20}
\colorlet{popgoldL}{popgold!20}
\colorlet{popredL}{popred!20}
\colorlet{popgrayL}{popgray!20}
% Teintes claires pour fonds de tableaux / zones
\colorlet{popblueL}{popblue!10!white}
\colorlet{poporangeL}{poporange!14!white}
\colorlet{popgreenL}{popgreen!12!white}
\colorlet{poppurpleL}{poppurple!12!white}
\colorlet{poppinkL}{poppink!12!white}
\colorlet{popgoldL}{popgold!14!white}

\pagecolor{popcream}

% ============================================================
% Typographie
% ============================================================
\defaultfontfeatures{Ligatures=TeX}
\setmainfont{PlusJakartaSans-Medium.ttf}[Path=../../../fonts/,BoldFont=PlusJakartaSans-Bold.ttf]
% Maths en sans-serif (Fira Math) avec chiffres et lettres droites de Plus
% Jakarta, pour que les formules se fondent dans le texte.
\usepackage[math-style=ISO,bold-style=ISO]{unicode-math}
\setmathfont{FiraMath-Regular.otf}[Path=../../../fonts/]
\setmathfont{PlusJakartaSans-Medium.ttf}[Path=../../../fonts/,range={up/num,up/{Latin,latin},\mathup}]
\usepackage{icomma} % virgule décimale sans espace en mode math
% Caractères mathématiques Unicode absents des polices de texte (Plus Jakarta,
% Archivo Black) : rendus par leur équivalent mathématique, en texte comme en
% formule — sinon ils s'affichent en carré vide (ex. « Arithmétique dans ℤ »).
\usepackage{newunicodechar}
\newunicodechar{ℝ}{\ensuremath{\mathbb{R}}}
\newunicodechar{ℤ}{\ensuremath{\mathbb{Z}}}
\newunicodechar{ℂ}{\ensuremath{\mathbb{C}}}
\newunicodechar{ℕ}{\ensuremath{\mathbb{N}}}
\newunicodechar{ℚ}{\ensuremath{\mathbb{Q}}}
\newunicodechar{𝔻}{\ensuremath{\mathbb{D}}}
\newunicodechar{α}{\ensuremath{\alpha}}
\newunicodechar{β}{\ensuremath{\beta}}
\newunicodechar{γ}{\ensuremath{\gamma}}
\newunicodechar{δ}{\ensuremath{\delta}}
\newunicodechar{ε}{\ensuremath{\varepsilon}}
\newunicodechar{θ}{\ensuremath{\theta}}
\newunicodechar{λ}{\ensuremath{\lambda}}
\newunicodechar{μ}{\ensuremath{\mu}}
\newunicodechar{σ}{\ensuremath{\sigma}}
\newunicodechar{τ}{\ensuremath{\tau}}
\newunicodechar{φ}{\ensuremath{\varphi}}
\newunicodechar{ψ}{\ensuremath{\psi}}
\newunicodechar{ω}{\ensuremath{\omega}}
\newunicodechar{Σ}{\ensuremath{\Sigma}}
% majuscules produites par \MakeUppercase dans les titres (α-aminés -> Α)
\newunicodechar{Α}{\ensuremath{\alpha}}
\newunicodechar{Β}{\ensuremath{\beta}}
\newunicodechar{Γ}{\ensuremath{\Gamma}}
\newunicodechar{Δ}{\ensuremath{\Delta}}
\newunicodechar{Ε}{\ensuremath{\varepsilon}}
\newunicodechar{Θ}{\ensuremath{\Theta}}
\newunicodechar{Λ}{\ensuremath{\Lambda}}
\newunicodechar{Μ}{\ensuremath{\mu}}
\newunicodechar{Τ}{\ensuremath{\tau}}
\newunicodechar{Φ}{\ensuremath{\Phi}}
\newunicodechar{Ψ}{\ensuremath{\Psi}}
\newunicodechar{Ω}{\ensuremath{\Omega}}
\newunicodechar{↦}{\ensuremath{\mapsto}}
\newunicodechar{∈}{\ensuremath{\in}}
\newunicodechar{∉}{\ensuremath{\notin}}
\newunicodechar{∧}{\ensuremath{\wedge}}
\newunicodechar{⇔}{\ensuremath{\Leftrightarrow}}
\newunicodechar{⟺}{\ensuremath{\Longleftrightarrow}}
\newunicodechar{⇒}{\ensuremath{\Rightarrow}}
\newunicodechar{⟹}{\ensuremath{\Longrightarrow}}
\newunicodechar{∘}{\ensuremath{\circ}}
\newunicodechar{∪}{\ensuremath{\cup}}
\newunicodechar{∩}{\ensuremath{\cap}}
\newunicodechar{≡}{\ensuremath{\equiv}}
\newunicodechar{⊂}{\ensuremath{\subset}}
\newunicodechar{⊥}{\ensuremath{\perp}}
\newunicodechar{∃}{\ensuremath{\exists}}
\newunicodechar{∀}{\ensuremath{\forall}}
\newunicodechar{∛}{\ensuremath{\sqrt[3]{\,}}}
\newunicodechar{∜}{\ensuremath{\sqrt[4]{\,}}}
\newunicodechar{⃗}{\ensuremath{{}^{\to}}}
\newunicodechar{ⁿ}{\ifmmode^{n}\else\textsuperscript{\ensuremath{n}}\fi}
\newunicodechar{ˣ}{\ifmmode^{x}\else\textsuperscript{\ensuremath{x}}\fi}
\newunicodechar{ᵇ}{\ifmmode^{b}\else\textsuperscript{\ensuremath{b}}\fi}
\newunicodechar{ʳ}{\ifmmode^{r}\else\textsuperscript{\ensuremath{r}}\fi}
\newunicodechar{ᵘ}{\ifmmode^{u}\else\textsuperscript{\ensuremath{u}}\fi}
\newunicodechar{ʸ}{\ifmmode^{y}\else\textsuperscript{\ensuremath{y}}\fi}
\newunicodechar{ᵃ}{\ifmmode^{a}\else\textsuperscript{\ensuremath{a}}\fi}
\newunicodechar{ᵉ}{\ifmmode^{e}\else\textsuperscript{\ensuremath{e}}\fi}
\newunicodechar{ᵏ}{\ifmmode^{k}\else\textsuperscript{\ensuremath{k}}\fi}
\newunicodechar{ᵖ}{\ifmmode^{p}\else\textsuperscript{\ensuremath{p}}\fi}
\newunicodechar{ᴺ}{\ifmmode^{N}\else\textsuperscript{\ensuremath{N}}\fi}
\newunicodechar{ᶜ}{\ifmmode^{c}\else\textsuperscript{\ensuremath{c}}\fi}
\newunicodechar{ʰ}{\ifmmode^{h}\else\textsuperscript{\ensuremath{h}}\fi}
\newunicodechar{ᵠ}{\ifmmode^{q}\else\textsuperscript{\ensuremath{q}}\fi}
\newunicodechar{ₙ}{\ifmmode_{n}\else\textsubscript{\ensuremath{n}}\fi}
\newunicodechar{ₐ}{\ifmmode_{a}\else\textsubscript{\ensuremath{a}}\fi}
\newunicodechar{ᵢ}{\ifmmode_{i}\else\textsubscript{\ensuremath{i}}\fi}
\newunicodechar{ᵣ}{\ifmmode_{r}\else\textsubscript{\ensuremath{r}}\fi}
\newunicodechar{ₖ}{\ifmmode_{k}\else\textsubscript{\ensuremath{k}}\fi}
\newunicodechar{ᵦ}{\ifmmode_{\beta}\else\textsubscript{\ensuremath{\beta}}\fi}
\newunicodechar{ₓ}{\ifmmode_{x}\else\textsubscript{\ensuremath{x}}\fi}
\newfontfamily\popdisplay{ArchivoBlack-Regular.ttf}[Path=../../../fonts/]
\newfontfamily\poplogo{SpaceGrotesk-Bold.ttf}[Path=../../../fonts/]
\newfontfamily\popsemi{PlusJakartaSans-SemiBold.ttf}[Path=../../../fonts/]
\newfontfamily\popxbold{PlusJakartaSans-ExtraBold.ttf}[Path=../../../fonts/]
\newfontfamily\popxboldls{PlusJakartaSans-ExtraBold.ttf}[Path=../../../fonts/,LetterSpace=3.0]

\setlist[enumerate]{leftmargin=1.7em,itemsep=5pt,topsep=4pt}
\setlist[itemize]{leftmargin=1.7em,itemsep=4pt,topsep=4pt,label={\color{poporange}\textbullet}}
\setlength{\parindent}{0pt}
\setlength{\parskip}{5pt}
\renewcommand{\arraystretch}{1.25}

% pgfplots : style Pop par défaut
\pgfplotsset{
  every axis/.append style={axis line style={popink,line width=1pt},tick style={popink},
    grid style={popline,line width=0.5pt},label style={font=\small},tick label style={font=\footnotesize}},
  popcurve/.style={poporange,line width=1.8pt,no markers,smooth},
}
% Même style disponible dans un \draw TikZ simple (hors pgfplots)
\tikzset{popcurve/.style={poporange,line width=1.8pt}}
\tikzset{popgrid/.style={popline,very thin}}
\colorlet{popcurve}{poporange} % le nom du style est parfois utilisé comme couleur

% ============================================================
% Pill / squiggle
% ============================================================
\newcommand{\poppill}[3]{%
  \tikz[baseline=-0.55ex]\node[draw=popink,line width=1.2pt,rounded corners=2.6pt,%
    fill=#2,inner xsep=6.5pt,inner ysep=3pt]{%
    \popxboldls\fontsize{7}{7}\selectfont\color{#3}\MakeUppercase{#1}};%
}
\newcommand{\popsquiggle}[1]{%
  \tikz[baseline=-0.4ex]{
    \draw[#1,line width=2.6pt,line cap=round]
      plot [smooth] coordinates {(0,0.09) (0.34,0.01) (0.68,0.21) (1.02,0.01) (1.36,0.10)};
  }%
}
% Mot-clé surligné (marqueur orange sous le texte)
\newcommand{\cle}[1]{{\popxbold\color{popdark}#1}}

% ============================================================
% En-tête / pied
% ============================================================
\pagestyle{fancy}
\fancyhf{}
\renewcommand{\headrulewidth}{0pt}
\renewcommand{\footrulewidth}{0pt}
\lhead{\raisebox{-0.15ex}{\poplogo\fontsize{13}{13}\selectfont\color{popink}OnBuch\color{poporange}+}}
\rhead{\poppill{\DOCMATIERE\ \textbullet\ \DOCNIVEAU\ \textbullet\ Leçon \DOCLECON}{white}{popink}}
\lfoot{\popxbold\fontsize{7.6}{7.6}\selectfont\color{popmuted}\MakeUppercase{Cours signé OnBuch+}\ \textbullet\ {\popsemi\DOCTITRE}}
\rfoot{\tikz[baseline=-0.6ex]\node[rounded corners=8.5pt,fill=popink,minimum height=17pt,minimum width=17pt,inner xsep=5pt,inner sep=0]{\popxbold\fontsize{8}{8}\selectfont\color{popcream}\thepage\,/\,\pageref*{LastPage}};}

% ============================================================
% Titres
% ============================================================
\newcommand{\chaptitre}[2]{%
  \begin{tcolorbox}[enhanced,colback=poporange,colframe=popink,boxrule=2.6pt,arc=11pt,
      drop shadow=popink,left=15pt,right=15pt,top=12pt,bottom=11pt,before skip=3pt,after skip=18pt]
    \poppill{#2}{popink}{popcream}\par\vspace{9pt}
    {\raggedright\hyphenpenalty=10000\exhyphenpenalty=10000\popdisplay\fontsize{22}{24}\selectfont\color{popcream}\MakeUppercase{#1}\par}
  \end{tcolorbox}
}
% Section numérotée : pastille orange avec le numéro + titre Archivo
\newcounter{popsec}
\newcommand{\coursec}[1]{%
  \stepcounter{popsec}\setcounter{popsub}{0}%
  \par\vspace{16pt}\noindent
  \begin{minipage}{\linewidth}
  \tikz[baseline=-0.7ex]\node[draw=popink,line width=1.6pt,fill=poporange,rounded corners=4pt,minimum size=22pt,inner sep=0]{\popdisplay\fontsize{12}{12}\selectfont\color{popcream}\thepopsec};%
  \hspace{8pt}{\popdisplay\fontsize{15}{17}\selectfont\color{popink}\MakeUppercase{#1}}\par
  \vspace{4pt}\noindent{\color{poporange}\rule{36pt}{3.6pt}}
  \end{minipage}\par\nopagebreak\vspace{8pt}
}
\newcounter{popsub}
\newcommand{\courssub}[1]{%
  \stepcounter{popsub}%
  \par\vspace{9pt}\noindent{\popxbold\fontsize{11.5}{13}\selectfont\color{popink}{\color{poporange}\thepopsec.\thepopsub}\hspace{6pt}#1}\par\nopagebreak\vspace{4pt}
}

% ============================================================
% Boîtes pédagogiques — même recette : fond blanc, bordure épaisse colorée,
% ombre dure encre, badge plein. Titre optionnel [#1].
% ============================================================
\tcbset{popbase/.style={breakable,enhanced,colback=white,boxrule=2.3pt,arc=9pt,
  shadow={3pt}{-3pt}{0pt}{popink},left=14pt,right=14pt,top=11pt,bottom=11pt,before skip=11pt,after skip=15pt,
  fontupper=\popsemi\fontsize{10.6}{14.5}\selectfont\color{popink},
  fontlower=\popsemi\fontsize{10.6}{14.5}\selectfont\color{popink}}}
% Coche dessinée (le glyphe ✓ n'existe pas dans les polices du document)
\newcommand{\popcheck}[1]{\tikz[baseline=-0.5ex]\draw[#1,line width=1.6pt,line cap=round,line join=round](-0.13,0.0)--(-0.04,-0.09)--(0.13,0.1);}
\providecommand{\coche}{\popcheck{popgreen}}
\providecommand{\resultat}[1]{\par\smallskip\noindent\fcolorbox{popgreen}{popgreenL}{\parbox{\dimexpr\linewidth-2\fboxsep-2\fboxrule\relax}{\textbf{Résultat.} #1}}\par\smallskip}
\newcommand{\popboxhead}[3]{\poppill{#1}{#2}{white}\ifx\relax#3\relax\else\hspace{6pt}{\popxbold\fontsize{10.6}{12}\selectfont\color{popink}#3}\fi\par\vspace{7pt}}

\newtcolorbox{aretenir}[1][]{popbase,colframe=popgreen,before upper={\popboxhead{À retenir}{popgreen}{#1}}}
% Texte en anglais (document de lecture, dialogue, extrait) : cadre
% d'étude. Un titre optionnel (source, auteur) s'affiche dans l'en-tête.
\newtcolorbox{textebox}[1][]{popbase,colframe=popblue,before upper={\popboxhead{Text}{popblue}{#1}\begin{otherlanguage}{english}},after upper={\end{otherlanguage}}}
% Marques ✓ / ✗ (forme correcte / fautive) souvent utilisées par les modèles
\providecommand{\cross}{\textcolor{poppink}{\ensuremath{\times}}}
\providecommand{\xmark}{\cross}\providecommand{\crossmark}{\cross}\providecommand{\cmark}{\ensuremath{\checkmark}}
% Ligne de réponse / justification à compléter (inventée par les modèles)
\providecommand{\justification}[1]{\par\noindent\textit{Justification~:} \dotfill\par}
% \textipa (package tipa non chargé) : la police ne couvre pas l'API -> simple passage
\providecommand{\textipa}[1]{#1}
% Soulignement (ulem non chargé) : \uline, \ul
\providecommand{\uline}[1]{\underline{#1}}\providecommand{\ul}[1]{\underline{#1}}
% \glossaire{\item ...} inventée par les modèles : liste de glossaire
\providecommand{\glossaire}[1]{\begin{itemize}#1\end{itemize}}
% \alert (beamer) inventée par les modèles : texte mis en évidence
\providecommand{\alert}[1]{\textbf{\textcolor{poppink}{#1}}}
% variantes ASCII d'ordinaux écrites en macro
\providecommand{\iere}{\textsuperscript{re}}\providecommand{\ieme}{\textsuperscript{e}}\providecommand{\ere}{\textsuperscript{re}}\providecommand{\eme}{\textsuperscript{e}}\providecommand{\er}{\textsuperscript{er}}\providecommand{\ie}{\textsuperscript{e}}
% Ordinaux français écrits en macro par les modèles : 1\ière, 3\ième
\newcommand{\ière}{\textsuperscript{re}}\newcommand{\ième}{\textsuperscript{e}}
% \exemple (inventée par les modèles) : étiquette « Exemple : »
\providecommand{\exemple}{\textbf{Exemple~:}\ }
% \square utilisable aussi hors mode math (cases à cocher)
\AtBeginDocument{\let\squareMath\square\renewcommand{\square}{\ensuremath{\squareMath}}}
% Mot ou phrase en anglais à l'intérieur d'un paragraphe français
\newcommand{\eng}[1]{\ifmmode\text{\textit{#1}}\else\begingroup\selectlanguage{english}\itshape #1\endgroup\fi}
\let\esp\eng % alias toléré
\newtcolorbox{exemplebox}[1][]{popbase,colframe=poporange,before upper={\popboxhead{Exemple}{poporange}{#1}}}
\newtcolorbox{definition}[1][]{popbase,colframe=popblue,before upper={\popboxhead{Définition}{popblue}{#1}}}
\newtcolorbox{propriete}[1][]{popbase,colframe=poppurple,before upper={\popboxhead{Propriété}{poppurple}{#1}}}
\newtcolorbox{methode}[1][]{popbase,colframe=popgold,before upper={\popboxhead{Méthode}{popgold}{#1}}}
\newtcolorbox{attention}[1][]{popbase,colframe=poppink,before upper={\popboxhead{Attention}{poppink}{#1}}}
\newtcolorbox{experience}[1][]{popbase,colframe=popblue,colback=popblueL,before upper={\popboxhead{Expérience}{popblue}{#1}}}
% « Le savais-tu ? » : carte violette pleine (contexte, culture, Cameroun)
\newtcolorbox{savaistu}[1][]{popbase,colback=poppurple,colframe=popink,
  fontupper=\popsemi\fontsize{10.4}{14.2}\selectfont\color{white},
  before upper={\poppill{Le savais-tu\,?}{white}{poppurple}\ifx\relax#1\relax\else\hspace{6pt}{\popxbold\color{white}#1}\fi\par\vspace{7pt}}}
% Exercice résolu : énoncé en haut, \tcblower puis la solution sur fond crème
\newcounter{popexo}\newcounter{popexr}
\newtcolorbox{exoresolu}[1][]{popbase,colframe=poporange,colbacklower=popcream,
  segmentation style={popink,line width=1.2pt,dashed},
  before upper={\stepcounter{popexr}\popboxhead{Exercice résolu \thepopexr}{poporange}{#1}},
  before lower={\poppill{Solution}{popink}{popcream}\par\vspace{6pt}}}
% Exercice (non résolu en place) — la correction est dans la partie Corrigés
\newtcolorbox{exercice}[1][]{popbase,colframe=popink,
  before upper={\stepcounter{popexo}\popboxhead{Exercice \thepopexo}{popink}{#1}}}
% Corrigé
\newtcolorbox{corrige}[1][]{popbase,colframe=popgreen,colback=popgreenL,
  before upper={\popboxhead{Corrigé}{popgreen}{#1}}}
% Objectifs / prérequis / bilan
\newtcolorbox{objectifs}{popbase,colframe=popink,colback=poporangeL,
  before upper={\popboxhead{Objectifs de la leçon}{popink}{}}}
\newtcolorbox{prerequis}{popbase,colframe=popink,colback=popblueL,
  before upper={\popboxhead{Prérequis}{popblue}{}}}
\newtcolorbox{bilan}{popbase,colframe=popink,colback=white,boxrule=2.8pt,
  before upper={\popboxhead{Fiche bilan}{poporange}{L'essentiel à connaître pour l'examen}}}
% Figure encadrée avec légende
\newtcolorbox{popfigure}[1][]{enhanced,breakable=false,colback=white,colframe=popline,boxrule=1.4pt,arc=8pt,
  left=10pt,right=10pt,top=10pt,bottom=8pt,before skip=10pt,after skip=12pt,halign=center,
  lower separated=false,halign lower=center,
  fontlower=\popsemi\fontsize{9}{11.5}\selectfont\color{popmuted},
  #1}
\newcommand{\legende}[1]{\par\vspace{6pt}{\popsemi\fontsize{9}{11.5}\selectfont\color{popmuted}#1}}

% ============================================================
% Signature OnBuch+
% ============================================================
\newcommand{\poplogoinline}[1]{{\poplogo\fontsize{#1}{#1}\selectfont\color{popink}OnBuch\color{poporange}+}}
\newcommand{\signatureonbuch}{%
  \begin{tcolorbox}[enhanced,colback=popink,colframe=popink,boxrule=2.2pt,arc=10pt,
      drop shadow=poporange,left=14pt,right=14pt,top=10pt,bottom=10pt,before skip=0pt,after skip=0pt]
    \tikz[baseline=-0.7ex]\node[circle,fill=popgreen,draw=popcream,line width=1.3pt,minimum size=22pt,inner sep=0]{\popcheck{white}};%
    \hspace{9pt}\begin{minipage}[c]{0.62\linewidth}
      {\popxbold\fontsize{12}{13}\selectfont\color{popcream}Signé\ \textbullet\ {\poplogo OnBuch\color{poporange}+}}\par\vspace{2pt}
      {\raggedright\popsemi\fontsize{8}{10.5}\selectfont\color{popline}\MakeUppercase{Équipe pédagogique OnBuch+}\\\MakeUppercase{Conforme au programme officiel MINESEC}\par}
    \end{minipage}\hfill
    {\poplogo\fontsize{22}{22}\selectfont\color{popcream}OnBuch\color{poporange}+}
  \end{tcolorbox}%
}

% ============================================================
% Page de garde
% #1 titre · #2 matière · #3 niveau · #4 module · #5 n° leçon · #6 durée ·
% #7 description / objectifs (texte ou itemize)
% ============================================================
\newcommand{\pagegarde}[7]{%
  \thispagestyle{empty}
  \newgeometry{a4paper,margin=1.7cm,top=1.6cm,bottom=1.4cm}%
  \begin{tikzpicture}[remember picture,overlay]
    \fill[poporange!18] ([xshift=-3.2cm,yshift=-3.6cm]current page.north east) circle (4.6cm);
    \fill[poppurple!14] ([xshift=2.4cm,yshift=7.6cm]current page.south west) circle (3.8cm);
  \end{tikzpicture}%
  {\poplogo\fontsize{30}{30}\selectfont\color{popink}OnBuch\color{poporange}+}\hfill
  \raisebox{0.6ex}{\poppill{Cours signé \textbullet\ Premium}{popink}{popcream}}\par
  \vspace{1pt}\popsquiggle{poppurple}\par
  \vspace{8pt}
  \begin{tcolorbox}[enhanced,colback=poporange,colframe=popink,boxrule=2.8pt,arc=13pt,
      drop shadow=popink,left=18pt,right=18pt,top=12pt,bottom=13pt,after skip=10pt]
    \poppill{#2 \textbullet\ #3}{popink}{popcream}\hspace{5pt}\poppill{#4}{white}{popink}\par\vspace{8pt}
    {\popdisplay\fontsize{14}{14}\selectfont\color{popink}LEÇON #5}\par\vspace{4pt}
    {\raggedright\hyphenpenalty=10000\exhyphenpenalty=10000\popdisplay\fontsize{25}{27}\selectfont\color{popcream}\MakeUppercase{#1}\par}
    \vspace{8pt}
    {\popxbold\fontsize{9.5}{9.5}\selectfont\color{popink}\MakeUppercase{Durée conseillée : #6}}
  \end{tcolorbox}
  \begin{tcolorbox}[enhanced,colback=white,colframe=popink,boxrule=2.2pt,arc=10pt,
      drop shadow=popink,left=16pt,right=16pt,top=10pt,bottom=10pt,after skip=8pt]
    \poppill{Dans cette leçon}{poppurple}{white}\par\vspace{5pt}
    \popsemi\fontsize{9.3}{12.6}\selectfont\color{popink}#7
  \end{tcolorbox}
  \noindent\begin{minipage}[t]{0.487\linewidth}
    \begin{tcolorbox}[enhanced,colback=popgreen,colframe=popink,boxrule=2.2pt,arc=10pt,
        drop shadow=popink,left=11pt,right=11pt,top=10pt,bottom=10pt,height=3.1cm,valign=center]
      \tcbox[colback=white,colframe=popink,boxrule=1.3pt,arc=5pt,boxsep=0pt,nobeforeafter,
        left=3pt,right=3pt,top=3pt,bottom=3pt,valign=center]{\qrcode[height=1.9cm]{https://play.google.com/store/apps/details?id=com.luvvix.onbuch&hl=fr}}%
      \hspace{8pt}\begin{minipage}[c]{0.5\linewidth}
        {\popdisplay\fontsize{11}{12.5}\selectfont\color{white}\MakeUppercase{Télécharge OnBuch}}\par\vspace{4pt}
        {\raggedright\popsemi\fontsize{8.4}{11}\selectfont\color{white}Gratuit \textbullet\ Google Play\\Tuteur IA, annales, concours, résultats\par}
      \end{minipage}
    \end{tcolorbox}
  \end{minipage}\hfill
  \begin{minipage}[t]{0.487\linewidth}
    \begin{tcolorbox}[enhanced,colback=poppurple,colframe=popink,boxrule=2.2pt,arc=10pt,
        drop shadow=popink,left=11pt,right=11pt,top=10pt,bottom=10pt,height=3.1cm,valign=center]
      \tikz[baseline=-0.7ex]\node[circle,fill=white,draw=popink,line width=1.3pt,minimum size=1.6cm,inner sep=0]{\popdisplay\fontsize{24}{24}\selectfont\color{poppurple}+};%
      \hspace{8pt}\begin{minipage}[c]{0.6\linewidth}
        {\popdisplay\fontsize{11}{12.5}\selectfont\color{white}\MakeUppercase{Passe à OnBuch+}}\par\vspace{4pt}
        {\raggedright\popsemi\fontsize{8.4}{11}\selectfont\color{white}Tous les cours signés, Tuteur IA illimité, fiches PDF — onglet OnBuch+ de l'app\par}
      \end{minipage}
    \end{tcolorbox}
  \end{minipage}
  \vspace{8pt}
  \popcontenu
  \vfill
  \signatureonbuch
  \restoregeometry
  \newpage
}

% Rangée « Ce cours contient » (page de garde)
\newcommand{\poptile}[3]{% #1 couleur #2 gros texte #3 légende
  \begin{tcolorbox}[enhanced,colback=#1,colframe=popink,boxrule=2pt,arc=9pt,shadow={2.5pt}{-2.5pt}{0pt}{popink},
      left=6pt,right=6pt,top=9pt,bottom=9pt,halign=center,valign=center,height=1.75cm,width=\linewidth,nobeforeafter]
    {\popdisplay\fontsize{12.5}{13}\selectfont\color{white}\MakeUppercase{#2}}\par\vspace{3pt}
    {\centering\popsemi\fontsize{7.8}{9.5}\selectfont\color{white}#3\par}
  \end{tcolorbox}}
\newcommand{\popcontenu}{%
  \par\noindent{\popxboldls\fontsize{7.5}{7.5}\selectfont\color{popmuted}CE COURS SIGNÉ CONTIENT}\par\vspace{7pt}
  \noindent\begin{minipage}[t]{0.235\linewidth}\poptile{popblue}{Cours}{complet, illustré, pas à pas}\end{minipage}\hfill
  \begin{minipage}[t]{0.235\linewidth}\poptile{poporange}{Méthodes}{et exercices résolus}\end{minipage}\hfill
  \begin{minipage}[t]{0.235\linewidth}\poptile{poppink}{Exercices}{type examen + corrigés}\end{minipage}\hfill
  \begin{minipage}[t]{0.235\linewidth}\poptile{popgreen}{Fiche bilan}{l'essentiel à retenir}\end{minipage}\par}

% Sommaire
\newtcolorbox{sommaire}{enhanced,colback=white,colframe=popink,boxrule=2.2pt,arc=10pt,
  drop shadow=popink,left=18pt,right=18pt,top=13pt,bottom=13pt,before skip=6pt,after skip=20pt,
  before upper={\poppill{Sommaire}{poppurple}{white}\par\vspace{9pt}\popsemi\fontsize{10.6}{14.5}\selectfont\color{popink}}}

% Signature de fin de cours — poussée tout en bas de la dernière page
\newcommand{\findecours}{%
  \par\vfill\nopagebreak
  \begin{center}\popsquiggle{poporange}\end{center}\vspace{4pt}
  \signatureonbuch
}
\AtBeginDocument{\def\llbracket{\mathopen{[\mkern-3.5mu[}}\def\rrbracket{\mathclose{]\mkern-3.5mu]}}} % intervalles d'entiers (glyphe ⟦ absent de la police)
% Glyphes absents de Fira Math : redéfinis à partir de glyphes disponibles
\AtBeginDocument{%
  \def\setminus{\mathbin{\backslash}}\let\smallsetminus\setminus
  \def\vdots{\mathord{\vbox{\baselineskip3.2pt\lineskiplimit0pt\kern4pt\hbox{.}\hbox{.}\hbox{.}}}}%
  \def\ddots{\mathinner{\mkern1mu\raise6.5pt\hbox{.}\mkern2mu\raise3.6pt\hbox{.}\mkern2mu\raise0.7pt\hbox{.}\mkern1mu}}%
  \def\triangle{\mathord{\scalebox{0.9}{$\symup{\Delta}$}}}%
  \def\nabla{\mathord{\raisebox{\depth}{\scalebox{1}[-1]{$\symup{\Delta}$}}}}%
  \def\checkmark{\popcheck{popdark}}%
  \def\star{\mathbin{\ast}}\def\bigstar{\mathord{\ast}}%
  \def\diamond{\mathbin{\scalebox{0.6}{$\Diamond$}}}%
  \def\bigcup{\mathop{\mathchoice{\raisebox{-0.3em}{\scalebox{1.7}{$\cup$}}}{\scalebox{1.25}{$\cup$}}{\cup}{\cup}}\displaylimits}%
  \def\bigcap{\mathop{\mathchoice{\raisebox{-0.3em}{\scalebox{1.7}{$\cap$}}}{\scalebox{1.25}{$\cap$}}{\cap}{\cap}}\displaylimits}%
  \def\lesseqqgtr{\mathrel{\lessgtr}}\def\bigoplus{\mathop{\oplus}\displaylimits}%
}
\providecommand{\ssi}{si et seulement si} % abréviation fréquente
\AtBeginDocument{\providecommand{\wideparen}{\overparen}} % arc de cercle

%% FIGFIX-BEGIN v5 — correctifs globaux des figures (docs/onbuch-generation/tools/figfix)
\usepackage{adjustbox}\usepackage{etoolbox}\tcbuselibrary{raster}
% toute figure TikZ/pgfplots plus large que la ligne est réduite à la largeur disponible
\BeforeBeginEnvironment{tikzpicture}{\begin{adjustbox}{max width=\linewidth}}
\AfterEndEnvironment{tikzpicture}{\end{adjustbox}}
% légendes pgfplots lisibles par-dessus les courbes
\pgfplotsset{every axis/.append style={legend style={fill=white,fill opacity=0.92,text opacity=1,draw=black!15,inner sep=3pt}}}
% étiquettes d'arêtes (syntaxe edge["..."]) avec fond blanc
\tikzset{every edge quotes/.append style={fill=white,inner sep=1.2pt}}
%% FIGFIX-END
\begin{document}

\pagegarde{Lesson 21 — Vocabulary: Words and expressions related to helping with fund-raising activities/support groups in the acquisition of community resources/services}{Anglais}{Tle A}{Chapter 2 : Economic Life and Occupations}{EN21}{2 h}{Cette leçon de vocabulaire présente les mots et expressions liés aux activités de collecte de fonds et aux groupes de soutien pour l'acquisition de ressources et services communautaires. Elle couvre les champs lexicaux, les collocations, les expressions idiomatiques, les faux amis et l'emploi du lexique en contexte, en vue du Bac.
\vspace{4pt}
\begin{multicols}{2}\small
\begin{itemize}
\item À la fin de cette leçon, je sais identifier et définir le vocabulaire des collectes de fonds (fund-raising) et des groupes de soutien (support groups)
\item employer correctement les collocations essentielles : raise funds, make a donation, launch an appeal, volunteer for...
\item distinguer les mots de la même famille : donate / donation / donor, contribute / contribution / contributor
\item éviter les faux amis et calques du français (to assist ≠ assister à, a donation ≠ une dotation...)
\item prononcer et orthographier correctement les mots clés du thème
\item réemployer le lexique dans des phrases, dialogues et courts textes sur la vie communautaire
\end{itemize}\end{multicols}}

\begin{sommaire}
\begin{multicols}{2}
\begin{enumerate}
\item Mise en situation et découverte du thème
\item Champ lexical 1 : Fund-raising activities
\item Champ lexical 2 : Support groups and community resources/services
\item Collocations, expressions idiomatiques et mots de la même famille
\item Faux amis et pièges des francophones
\item Emploi du lexique en contexte et production guidée
\item Activité d'intégration
\item Méthodes et exercices résolus
\item Exercices
\item Corrigés des exercices
\item Fiche bilan
\end{enumerate}
\end{multicols}
\end{sommaire}

\begin{objectifs}
\begin{itemize}
\item À la fin de cette leçon, je sais identifier et définir le vocabulaire des collectes de fonds (fund-raising) et des groupes de soutien (support groups)
\item employer correctement les collocations essentielles : raise funds, make a donation, launch an appeal, volunteer for...
\item distinguer les mots de la même famille : donate / donation / donor, contribute / contribution / contributor
\item éviter les faux amis et calques du français (to assist ≠ assister à, a donation ≠ une dotation...)
\item prononcer et orthographier correctement les mots clés du thème
\item réemployer le lexique dans des phrases, dialogues et courts textes sur la vie communautaire
\item rédiger un court texte (annonce, appel à contributions, lettre de remerciement) avec le vocabulaire approprié
\item comprendre un texte ou une conversation portant sur une action caritative ou communautaire
\end{itemize}
\end{objectifs}

\begin{prerequis}
\begin{itemize}
\item Vocabulaire général de la vie communautaire vu en Première (community, development, project)
\item Les verbes suivis de l'infinitif ou du gérondif (to help + BV, to contribute to + V-ing)
\item Le présent simple et le present perfect pour décrire des actions et leurs résultats
\item Les articles et les noms dénombrables/indénombrables (money, equipment, information)
\end{itemize}
\end{prerequis}

\newpage

\coursec{Mise en situation et découverte du thème}

\begin{experience}[Situation d'accroche : two appeals for help]
Lisez attentivement la situation suivante, puis réagissez à l'oral avec votre voisin.

\eng{In Bamenda, the Parent-Teacher Association of GBHS is organizing a fund-raising gala to buy computers for the school library. Volunteers are needed to mobilize donors, sell raffle tickets and manage the budget. At the same time, a support group for flood victims in Limbe is appealing for contributions to acquire blankets and medicines. How do we talk about all these actions in correct English? This lesson gives you the exact words and expressions to describe fund-raising activities and support groups working for community resources and services.}

« À Bamenda, l'association parents-enseignants du GBHS organise un gala de collecte de fonds pour acheter des ordinateurs pour la bibliothèque de l'école. On a besoin de volontaires pour mobiliser les donateurs, vendre des billets de tombola et gérer le budget. Au même moment, un groupe de soutien aux victimes des inondations de Limbé lance un appel à contributions pour acquérir des couvertures et des médicaments. Comment parler de toutes ces actions dans un anglais correct ? Cette leçon vous donne les mots et expressions exacts pour décrire les activités de collecte de fonds et les groupes de soutien qui œuvrent pour les ressources et services de la communauté. »

\textbf{Problématique :} quels mots anglais permettent de nommer ces actions de solidarité, ces acteurs et ces ressources, sans calquer le français ?
\end{experience}

\courssub{Lecture découverte : un texte support}

\begin{textebox}[The PTA Fund-raising Day at GBHS Bamenda]
Last Saturday, the PTA of GBHS Bamenda organized a fund-raising gala in the school hall. Parents, alumni and well-wishers made generous donations to buy computers for the school library. Volunteers sold raffle tickets at the entrance, while others collected contributions in baskets. The chairperson of the organizing committee announced that a support group of former students had pledged to maintain the equipment for three years. A local businessman offered to sponsor the event and promised to match every franc raised. During the evening, a gospel choir performed free of charge, and a famous artist donated two paintings for the auction. Some guests gave cash; others made in-kind donations such as books, chairs and a photocopier. The treasurer kept a careful record of every amount received. By the end of the day, the organizers had raised over two million francs CFA. The headmaster thanked all the benefactors and said the new computer room would open in January. He also appealed to the community to set up a permanent fund for poor students who cannot pay their fees. Everybody agreed that when a community pulls together, no project is too big.
\end{textebox}

\textbf{Glossaire (mots difficiles) :}

\begin{center}
\begin{tabularx}{\linewidth}{|l|l|X|}\hline
\rowcolor{popblueL} \textbf{Mot anglais} & \textbf{Nature} & \textbf{Traduction française} \\ \hline
fund-raising gala & nom & gala de collecte de fonds \\ \hline
alumni & nom pluriel & anciens élèves \\ \hline
well-wisher & nom & sympathisant, personne bienveillante \\ \hline
donation & nom & don \\ \hline
raffle ticket & nom & billet de tombola \\ \hline
contribution & nom & contribution, cotisation \\ \hline
to pledge & verbe & promettre (formellement), s'engager à \\ \hline
to sponsor & verbe & parrainer, financer \\ \hline
auction & nom & vente aux enchères \\ \hline
in-kind donation & nom & don en nature \\ \hline
treasurer & nom & trésorier(-ière) \\ \hline
benefactor & nom & bienfaiteur(-trice) \\ \hline
to appeal to & verbe & lancer un appel à \\ \hline
to pull together & verbe & se serrer les coudes, unir ses efforts \\ \hline
\end{tabularx}
\end{center}

\courssub{Repérage des mots nouveaux}

Relisez le texte et surlignez tous les mots liés à deux grandes familles :

\begin{itemize}
\item \textbf{La collecte de fonds} (\eng{fund-raising}) : \eng{donation, raffle ticket, contribution, to pledge, to sponsor, auction, to raise money, benefactor}.
\item \textbf{Les groupes de soutien et les ressources} (\eng{support groups and resources}) : \eng{support group, volunteers, committee, treasurer, equipment, fees, fund}.
\end{itemize}

\begin{definition}[Fund-raising]
\eng{Fund-raising} (ou \eng{fundraising}) : \emph{the activity of collecting money for a charity, a school, a church or a community project}. Traduction : « collecte de fonds ». Le verbe correspondant est \eng{to raise funds / to raise money} (« collecter des fonds / de l'argent ») et la personne qui collecte est \eng{a fund-raiser} (« un collecteur de fonds »).
\end{definition}

\begin{methode}[Deviner le sens d'un mot inconnu avant de vérifier]
Face à un mot nouveau dans un texte, ne courez pas tout de suite au dictionnaire. Suivez ces quatre étapes :
\begin{enumerate}
\item \textbf{Identifiez sa nature} : est-ce un nom, un verbe, un adjectif ? Observez sa place dans la phrase et ses marques (\eng{-tion} = souvent un nom ; \eng{-ed} = souvent un verbe au passé).
\item \textbf{Observez les mots qui l'entourent} : le contexte donne des indices (ex. \eng{the treasurer kept a record of every amount received} → on devine que \eng{treasurer} est une personne qui gère l'argent).
\item \textbf{Proposez un équivalent français} plausible, même approximatif.
\item \textbf{Vérifiez} ensuite dans le glossaire ou le dictionnaire pour confirmer ou corriger votre hypothèse.
\end{enumerate}
Cette stratégie, appelée \eng{guessing from context}, est essentielle à l'épreuve de compréhension de l'écrit.
\end{methode}

\begin{exoresolu}[Guessing from context]
Énoncé : dans la phrase \eng{A local businessman offered to sponsor the event and promised to match every franc raised}, devinez le sens de \eng{to sponsor} et de \eng{to match} sans dictionnaire.
\tcblower
Solution détaillée :
\begin{enumerate}
\item \eng{to sponsor} est un verbe (précédé de \eng{offered to}). Le sujet est \eng{a local businessman} et l'objet \eng{the event} : un homme d'affaires fait quelque chose pour un événement de collecte de fonds. Hypothèse : « financer, parrainer ».
\item \eng{to match} est un verbe ; son objet est \eng{every franc raised} (« chaque franc collecté »). Hypothèse : « égaler, donner l'équivalent » — si la foule donne 100 000 FCFA, lui donne aussi 100 000 FCFA.
\item Vérification au glossaire/dictionnaire : \eng{to sponsor} = parrainer ; \eng{to match} = égaler. Hypothèses confirmées.
\end{enumerate}
Traduction complète : « Un homme d'affaires local a proposé de parrainer l'événement et a promis de verser l'équivalent de chaque franc collecté. »
\end{exoresolu}

\courssub{Questions de compréhension}

\begin{exercice}[Comprehension check]
\textbf{A. True or False?} Justify with a line from the text.
\begin{enumerate}
\item The gala was organized to buy books for the library.
\item Former students promised to maintain the computers.
\item The artist sold his paintings and kept the money.
\item More than two million francs were collected.
\end{enumerate}
\textbf{B. Multiple choice.} Choose the correct answer.
\begin{enumerate}
\item The event took place (a) in a church (b) in the school hall (c) in a hotel.
\item The businessman promised to (a) sing with the choir (b) match every franc raised (c) buy the paintings.
\item The headmaster wants a permanent fund for (a) poor students (b) new teachers (c) the football team.
\end{enumerate}
\textbf{C. Short answers.}
\begin{enumerate}
\item Who kept a record of the money received?
\item What did the guests give besides cash?
\item When will the computer room open?
\end{enumerate}
\end{exercice}

\begin{corrige}[Comprehension check]
\textbf{A.} 1. False — the gala was organized to buy \emph{computers}. 2. True — \eng{a support group of former students had pledged to maintain the equipment}. 3. False — he \emph{donated} two paintings for the auction. 4. True — \eng{the organizers had raised over two million francs CFA}.
\textbf{B.} 1. (b) 2. (b) 3. (a).
\textbf{C.} 1. The treasurer. 2. In-kind donations: books, chairs and a photocopier. 3. In January.
\end{corrige}

\courssub{Carte mentale du thème}

\begin{popfigure}
\begin{tikzpicture}[mindmap, grow cyclic, every node/.style={concept, font=\small}, concept color=popblue!60, level 1/.append style={level distance=3.2cm, sibling angle=72}, level 2/.append style={level distance=2.2cm, sibling angle=45}]
\node[concept color=poporange!70, font=\small\bfseries] {Helping the community}
  child[concept color=popblue!50] { node {fund-raising} }
  child[concept color=popgreen!50] { node {support groups} }
  child[concept color=poppurple!50] { node {resources} }
  child[concept color=poppink!50] { node {services} }
  child[concept color=popgold!60] { node {people involved} };
\end{tikzpicture}
\legende{Carte mentale du thème : cinq branches lexicales à explorer dans la leçon.}
\end{popfigure}

Chaque branche sera développée dans les sections suivantes : les activités de collecte (\eng{gala, raffle, auction}), les groupes (\eng{support group, committee, volunteers}), les ressources (\eng{equipment, computers, blankets}), les services (\eng{maintenance, library, health care}) et les personnes (\eng{donor, sponsor, benefactor, treasurer}).

\courssub{Brainstorming : la solidarité autour de nous}

\begin{experience}[Class brainstorming]
En groupes de quatre, listez en cinq minutes toutes les actions de solidarité qui existent dans votre quartier ou votre village, puis essayez de les nommer en anglais :
\begin{itemize}
\item la \textbf{njangi} ou tontine ;
\item la \textbf{harvest thanksgiving} (action de grâce de la moisson) à l'église ;
\item le \textbf{community work} (travail communautaire : nettoyage du marché, réparation d'un pont) ;
\item les collectes pour les funérailles, les mariages, les frais scolaires d'un orphelin.
\end{itemize}
Chaque groupe présente ensuite deux phrases à la classe, par exemple : \eng{In our quarter, women organize a njangi every month to help members pay school fees.}
\end{experience}

\begin{savaistu}[Njangi, tontine... en anglais ?]
Au Cameroun anglophone, les \emph{njangi} ou tontines sont une forme traditionnelle de solidarité financière : chaque membre cotise régulièrement et le total revient à tour de rôle à l'un d'eux. En anglais, on dit \eng{a rotating savings and credit association (ROSCA)} ou, plus simplement, \eng{a contribution club} ou \eng{a savings club}. Exemple : \eng{My mother belongs to a njangi that meets every Sunday.} Attention : le mot \eng{njangi} est accepté dans l'anglais camerounais, mais dans un examen international, expliquez-le : \eng{a njangi, that is, a rotating savings group}.
\end{savaistu}

\begin{attention}[Premiers pièges des francophones]
\begin{itemize}
\item Ne dites pas \eng{to collect funds} pour « collecter des fonds » au sens noble : le verbe naturel est \eng{to raise funds / money}. \eng{They raised two million francs} (« Ils ont collecté deux millions de francs »).
\item \eng{A donation} est un « don », pas une « dotation » administrative.
\item \eng{To support} signifie « soutenir » (aider) : \eng{We support the project} = « Nous soutenons le projet ». Ne traduisez pas « supporter une équipe » par \eng{support} seul dans tous les contextes ; pour le football, on dit aussi \eng{to support a team}, c'est correct, mais « supporter la douleur » se dit \eng{to bear / to endure}.
\item Prononciation : \eng{fund} se prononce « feunde » (voyelle courte), \eng{raise} se prononce « réize » (comme \eng{days}), et \eng{donation} « do-né-cheune » avec l'accent sur la deuxième syllabe.
\end{itemize}
\end{attention}

\begin{aretenir}[Bilan de la mise en situation]
\begin{itemize}
\item Le texte support nous a permis de découvrir en contexte le vocabulaire de la solidarité : \eng{fund-raising, donation, to pledge, to sponsor, support group, volunteers, benefactor}.
\item La stratégie \eng{guessing from context} (nature du mot → contexte → hypothèse → vérification) permet de comprendre un texte sans traduire chaque mot.
\item Le thème se déploie en cinq branches : \eng{fund-raising, support groups, resources, services, people involved}.
\end{itemize}
\end{aretenir}

\coursec{Champ lexical 1 : Fund-raising activities}

Dans la section précédente, nous avons découvert le thème général de la leçon à travers la situation de l'école de Buea qui cherche à financer un nouveau forage (\eng{borehole}) et une bibliothèque. Nous avons entendu des mots comme \eng{funds}, \eng{donation}, \eng{gala}, \eng{raffle}. Il est temps maintenant d'organiser ce vocabulaire de manière systématique. Nous allons explorer le premier grand champ lexical : celui des \cle{activités de collecte de fonds} (\eng{fund-raising activities}), en classant les mots en quatre familles logiques : les \textbf{acteurs} (people), les \textbf{actions} (actions), les \textbf{événements} (events) et les \textbf{résultats} (results).

\courssub{Observation : un texte pour découvrir le vocabulaire en contexte}

\begin{textebox}[A fund-raising campaign in Buea — reading text]
The Parent-Teacher Association of Government High School Buea has launched an appeal to raise funds for a new borehole and a school library. The organisers have planned several events. First, a fund-raising gala will take place in the school hall in December. Local businesses have agreed to act as sponsors: a brewery will provide drinks, and a bank has pledged 500,000 FCFA. Volunteers will sell raffle tickets at the market every Saturday, and the proceeds from the raffle will buy books for the library.

A charity concert is also scheduled, with famous artists performing for free. In addition, the school is organising a sponsored walk from Buea Town to Mile 17: each participant will collect money from well-wishers for every kilometre walked. Some benefactors prefer to make a donation in kind: a trader from Douala has already donated bags of cement, and a retired teacher has contributed fifty books.

The target is five million FCFA. So far, the amount raised is two million FCFA in cash, plus donations in kind. “Every contribution counts,” says the chairperson of the organising committee. “We thank all our donors and contributors. Together, we will give our children clean water and books.”
\end{textebox}

\textbf{Glossaire (mots difficiles)}

\begin{center}
\begin{tabularx}{\linewidth}{|l|l|X|}
\hline
\rowcolor{popblueL} \eng{English} & Nature & Français \\
\hline
\eng{borehole} & nom & forage, puits foré \\
\hline
\eng{to pledge} & verbe & promettre (un don), s'engager à donner \\
\hline
\eng{proceeds} & nom (pluriel) & recettes, produit (d'une vente) \\
\hline
\eng{well-wisher} & nom & sympathisant, personne bienveillante \\
\hline
\eng{donation in kind} & expression & don en nature (biens ou services) \\
\hline
\eng{target} & nom & objectif (à atteindre) \\
\hline
\eng{chairperson} & nom & président(e) (d'un comité) \\
\hline
\end{tabularx}
\end{center}

\textbf{Questions de compréhension (orales ou écrites)}

\begin{enumerate}
\item What has the Parent-Teacher Association launched an appeal for?
\item Name three fund-raising events mentioned in the text.
\item Who pledged 500,000 FCFA?
\item What will the proceeds from the raffle be used for?
\item What is the difference between a donation in cash and a donation in kind? Give one example of each from the text.
\end{enumerate}

\courssub{Les acteurs : people involved in fund-raising}

\begin{definition}[Les acteurs de la collecte de fonds]
\begin{itemize}
\item \cle{donor} (« don-neur ») : a person or organization that gives money or goods — un donateur.
\item \cle{fundraiser} (« fand-ré-zeur ») : a person whose job or role is to collect money for a cause — un collecteur de fonds, un organisateur de collecte. Attention : \eng{a fundraiser} peut aussi désigner l'événement lui-même (\eng{The gala was a successful fundraiser}).
\item \cle{organiser} (« or-ga-naï-zeur ») : the person who plans and arranges the event — l'organisateur. (Variante US : \eng{organizer}).
\item \cle{volunteer} (« vo-lon-tier ») : a person who works for free — un bénévole, un volontaire.
\item \cle{sponsor} (« spon-seur ») : a person or company that gives money to support an event, often in exchange for publicity — un sponsor, un mécène.
\item \cle{well-wisher} (« ouèl-ouï-cheur ») : a person who supports a cause with sympathy, good wishes or small gifts — un sympathisant.
\item \cle{benefactor} (« bé-ni-fak-teur ») : a person who gives generous help, especially money, to a good cause — un bienfaiteur.
\item \cle{contributor} (« kon-tri-biou-teur ») : anyone who gives something (money, goods, time) to a common fund — un contributeur.
\end{itemize}
\end{definition}

\begin{exemplebox}[Les acteurs en contexte]
\eng{Many donors responded to the school’s appeal.} « De nombreux donateurs ont répondu à l'appel de l'école. »

\eng{Volunteers sold raffle tickets at Mokolo market.} « Des bénévoles ont vendu des tickets de tombola au marché Mokolo. »

\eng{The brewery agreed to be the main sponsor of the gala.} « La brasserie a accepté d'être le sponsor principal du gala. »

\eng{An anonymous benefactor paid for the roof of the new classroom block.} « Un bienfaiteur anonyme a payé le toit du nouveau bloc de classes. »
\end{exemplebox}

\begin{attention}[Faux amis : sponsor, donateur, volontaire]
\begin{itemize}
\item \eng{Sponsor} se prononce « spon-seur » mais s'écrit avec un \textbf{-or} final, comme \eng{donor}, \eng{organiser} (anglais britannique) / \eng{organizer} (anglais américain). Ne pas écrire \emph{sponser}.
\item Un \eng{volunteer} n'est pas seulement un « volontaire » au sens militaire : c'est d'abord un \textbf{bénévole}, quelqu'un qui travaille gratuitement.
\item Ne confondez pas \eng{donor} (celui qui donne) et \eng{donee} (celui qui reçoit, terme juridique rare).
\end{itemize}
\end{attention}

\courssub{Les actions : what people do}

\begin{definition}[Les verbes et expressions de la collecte de fonds]
\begin{itemize}
\item \cle{to raise funds / to raise money} : collecter des fonds / de l'argent.
\item \cle{to make a donation} : faire un don.
\item \cle{to give / to contribute money} : donner / contribuer de l'argent.
\item \cle{to launch an appeal} : lancer un appel (aux dons).
\item \cle{to organise a gala / a charity concert / a sponsored walk} : organiser un gala / un concert de charité / une marche sponsorisée. (Variante US : \eng{organize}).
\item \cle{to sell raffle tickets} : vendre des tickets de tombola.
\item \cle{to collect money} : collecter de l'argent, faire la quête.
\item \cle{to pledge} : promettre un don, s'engager à verser une somme.
\item \cle{to donate} : faire don de, donner (à une œuvre).
\end{itemize}
\end{definition}

\begin{propriete}[Le verbe to raise : forme et emploi]
\eng{To raise} est un verbe \textbf{régulier} : \textbf{raise – raised – raised}. Il est toujours \textbf{transitif} : il prend obligatoirement un complément d'objet direct.

\begin{center}
\begin{tabular}{|l|l|l|}
\hline
\rowcolor{popblueL} Base verbale & Prétérit & Participe passé \\
\hline
raise & raised & raised \\
\hline
rise & rose & risen \\
\hline
\end{tabular}
\end{center}

Collocations fréquentes : \eng{to raise funds} (collecter des fonds), \eng{to raise money} (collecter de l'argent), \eng{to raise awareness} (sensibiliser), \eng{to raise a question} (soulever une question).

\eng{The students raised two million FCFA last term.} « Les élèves ont collecté deux millions de FCFA le trimestre dernier. »

\textbf{Ne pas confondre avec} \eng{to rise} (rise – rose – risen), qui est \textbf{intransitif} : il ne prend jamais de complément d'objet. \eng{Prices are rising.} « Les prix augmentent. » On ne peut pas dire \emph{rise funds} : on « lève » quelque chose (\eng{raise}), le soleil, lui, « se lève » (\eng{rise}).
\end{propriete}

\begin{attention}[Les prépositions obligatoires : ne pas calquer le français]
En anglais, la préposition est fixée par le verbe et ne se devine pas depuis le français :
\begin{itemize}
\item On dit \eng{to raise funds \textbf{for} a project} : \eng{They raised funds for a new borehole.} « Ils ont collecté des fonds \textbf{pour} un nouveau forage. »
\item On dit \eng{to donate money \textbf{to} an organization} : \eng{She donated money to the Red Cross.} « Elle a donné de l'argent \textbf{à} la Croix-Rouge. »
\item On dit \eng{to contribute \textbf{to} a fund} : \eng{Everyone contributed to the appeal fund.} « Chacun a contribué à la caisse de l'appel. »
\item On dit \eng{to participate \textbf{in} an event} : \eng{Fifty students participated in the sponsored walk.} « Cinquante élèves ont participé à la marche sponsorisée. »
\end{itemize}
Ne dites jamais \emph{raise funds to a project} ni \emph{donate money for an organization}.
\end{attention}

\courssub{Les événements : fund-raising events}

\begin{definition}[Les événements de collecte de fonds]
\begin{itemize}
\item \cle{fund-raising gala} : gala de bienfaisance (soirée de prestige avec vente de billets).
\item \cle{charity concert} : concert de charité (les artistes jouent gratuitement).
\item \cle{fun fair} : fête foraine, kermesse (jeux et stands dont les recettes vont à une œuvre).
\item \cle{sponsored walk / run} : marche / course sponsorisée (chaque participant se fait promettre une somme par kilomètre).
\item \cle{raffle / tombola} : tombola (vente de tickets avec tirage au sort de lots).
\item \cle{auction} (« ok-cheune ») : vente aux enchères.
\item \cle{harvest thanksgiving} : action de grâce de la moisson (cérémonie religieuse avec offrandes, très courante dans les églises du Cameroun anglophone).
\item \cle{appeal fund} : fonds d'appel, caisse spéciale ouverte pour une cause précise.
\end{itemize}
\end{definition}

\begin{exemplebox}[Les événements en contexte]
\eng{The school launched an appeal to raise funds for a new borehole.} « L'école a lancé un appel pour collecter des fonds pour un nouveau forage. »

\eng{The charity concert attracted more than a thousand people.} « Le concert de charité a attiré plus de mille personnes. »

\eng{They organised an auction of paintings by local artists.} « Ils ont organisé une vente aux enchères de tableaux d'artistes locaux. »

\eng{During the harvest thanksgiving, the congregation offered food and money.} « Pendant l'action de grâce de la moisson, l'assemblée a offert des vivres et de l'argent. »
\end{exemplebox}

\courssub{Les résultats : money and outcomes}

\begin{definition}[Les mots du résultat]
\begin{itemize}
\item \cle{proceeds} (« pro-ssides », toujours au pluriel) : les recettes, le produit d'une vente ou d'un événement.
\item \cle{amount raised} : la somme collectée.
\item \cle{target} : l'objectif (financier) à atteindre.
\item \cle{budget} : le budget prévisionnel.
\item \cle{contribution} : la contribution (de chacun).
\item \cle{donation in cash} : don en espèces.
\item \cle{donation in kind} : don en nature.
\end{itemize}
\end{definition}

\begin{definition}[Donation in kind]
A \cle{donation in kind} is a gift of goods or services, not money — un don en nature : vivres, matériel, médicaments, livres, ciment, ou encore travail gratuit.

\eng{The trader made a donation in kind: twenty bags of rice.} « Le commerçant a fait un don en nature : vingt sacs de riz. »

À l'opposé : \eng{a donation in cash} (un don en espèces) : \eng{The bank made a cash donation of one million FCFA.} « La banque a fait un don en espèces d'un million de FCFA. »
\end{definition}

\begin{exemplebox}[Les résultats en contexte]
\eng{Proceeds from the raffle will buy books for the library.} « Les recettes de la tombola achèteront des livres pour la bibliothèque. »

\eng{Our target is five million FCFA, and we have already reached half of it.} « Notre objectif est de cinq millions de FCFA, et nous en avons déjà atteint la moitié. »

\eng{The total amount raised exceeded all expectations.} « La somme totale collectée a dépassé toutes les attentes. »
\end{exemplebox}

\begin{attention}[Proceeds : un pluriel piégeux]
\eng{Proceeds} s'emploie \textbf{toujours au pluriel} et avec un verbe au pluriel : \eng{The proceeds \textbf{were} given to the hospital.} « Les recettes ont été remises à l'hôpital. » Ne dites pas \emph{the proceed was}. De même, \eng{money} est indénombrable : \eng{The money \textbf{was} banked} (« L'argent a été mis à la banque »), jamais \emph{monies} au sens courant.
\end{attention}

\courssub{Carte mentale et tableau récapitulatif}

\begin{popfigure}
\begin{tikzpicture}[
mindmap,
grow cyclic,
every node/.style={concept, font=\small},
concept color=popblue,
level 1/.append style={level distance=3.2cm, sibling angle=90},
level 2/.append style={level distance=2.2cm, sibling angle=45, font=\scriptsize},
]
\node[concept, concept color=poporange, text=white] {Fund-raising}
child[concept color=popblue, text=white] { node {People}
  child { node {donor, benefactor} }
  child { node {volunteer, organiser} }
  child { node {sponsor, well-wisher} }
}
child[concept color=popgreen, text=white] { node {Actions}
  child { node {raise funds, pledge} }
  child { node {donate, contribute} }
  child { node {launch an appeal} }
}
child[concept color=poppurple, text=white] { node {Events}
  child { node {gala, concert} }
  child { node {raffle, auction} }
  child { node {sponsored walk} }
}
child[concept color=popgold, text=white] { node {Results}
  child { node {proceeds, target} }
  child { node {amount raised} }
  child { node {donation in kind} }
};
\end{tikzpicture}
\legende{Carte mentale : les quatre familles du champ lexical de la collecte de fonds.}
\end{popfigure}

\begin{center}
\begin{tabularx}{\linewidth}{|l|l|X|}
\hline
\rowcolor{popblueL} \eng{English} & Français & \eng{Exemple court} \\
\hline
\eng{donor} & donateur & \eng{Many donors answered the appeal.} \\
\hline
\eng{fundraiser} & collecteur de fonds & \eng{She is a professional fundraiser.} \\
\hline
\eng{organiser} & organisateur & \eng{The organisers booked the hall.} \\
\hline
\eng{volunteer} & bénévole & \eng{Volunteers collected money at the gate.} \\
\hline
\eng{sponsor} & sponsor, mécène & \eng{The bank is our main sponsor.} \\
\hline
\eng{well-wisher} & sympathisant & \eng{Well-wishers sent gifts and cards.} \\
\hline
\eng{benefactor} & bienfaiteur & \eng{A benefactor paid for the roof.} \\
\hline
\eng{contributor} & contributeur & \eng{Every contributor received a receipt.} \\
\hline
\eng{to raise funds} & collecter des fonds & \eng{They raised funds for a borehole.} \\
\hline
\eng{to make a donation} & faire un don & \eng{He made a donation of 10,000 FCFA.} \\
\hline
\eng{to contribute money} & contribuer de l'argent & \eng{Parents contributed money and food.} \\
\hline
\eng{to launch an appeal} & lancer un appel & \eng{The school launched an appeal.} \\
\hline
\eng{to sell raffle tickets} & vendre des tickets de tombola & \eng{Students sold raffle tickets at the market.} \\
\hline
\eng{to collect money} & collecter de l'argent & \eng{We collected money at the church door.} \\
\hline
\eng{to pledge} & promettre (un don) & \eng{The company pledged 500,000 FCFA.} \\
\hline
\eng{to donate} & faire don de & \eng{She donated fifty books to the library.} \\
\hline
\eng{fund-raising gala} & gala de bienfaisance & \eng{The gala will take place in December.} \\
\hline
\eng{charity concert} & concert de charité & \eng{The charity concert was a success.} \\
\hline
\eng{fun fair} & kermesse, fête foraine & \eng{The fun fair raised 200,000 FCFA.} \\
\hline
\eng{sponsored walk} & marche sponsorisée & \eng{He walked 10 km in the sponsored walk.} \\
\hline
\eng{raffle / tombola} & tombola & \eng{She won a goat in the raffle.} \\
\hline
\eng{auction} & vente aux enchères & \eng{They held an auction of paintings.} \\
\hline
\eng{harvest thanksgiving} & action de grâce de la moisson & \eng{The harvest thanksgiving filled the church.} \\
\hline
\eng{appeal fund} & fonds d'appel & \eng{The appeal fund is growing every week.} \\
\hline
\eng{proceeds} & recettes & \eng{Proceeds from the raffle will buy books.} \\
\hline
\eng{amount raised} & somme collectée & \eng{The amount raised reached two million FCFA.} \\
\hline
\eng{target} & objectif & \eng{Our target is five million FCFA.} \\
\hline
\eng{budget} & budget & \eng{The committee prepared a careful budget.} \\
\hline
\eng{donation in kind} & don en nature & \eng{He made a donation in kind: bags of cement.} \\
\hline
\eng{donation in cash} & don en espèces & \eng{The firm made a cash donation.} \\
\hline
\end{tabularx}
\end{center}

\courssub{Exercice résolu et entraînement}

\begin{exoresolu}[Complétez avec le mot qui convient]
Énoncé. Complétez chaque phrase avec le mot ou l'expression approprié(e) : \eng{proceeds}, \eng{pledged}, \eng{volunteers}, \eng{target}, \eng{sponsor}, \eng{donation in kind}.

\begin{enumerate}
\item The \ldots\ldots\ldots of the campaign is three million FCFA.
\item A soft-drink company agreed to be the official \ldots\ldots\ldots of the concert.
\item Thirty \ldots\ldots\ldots sold tickets at the entrance of the stadium.
\item The trader made a \ldots\ldots\ldots : he gave ten bags of rice, not money.
\item The \ldots\ldots\ldots from the auction will pay for new desks.
\item The minister \ldots\ldots\ldots one million FCFA but has not paid yet.
\end{enumerate}

\tcblower

Solution détaillée.
\begin{enumerate}
\item \eng{target} — il s'agit de l'objectif financier à atteindre.
\item \eng{sponsor} — une entreprise qui soutient financièrement un événement.
\item \eng{volunteers} — des bénévoles, qui travaillent gratuitement ; le verbe \eng{sold} confirme qu'il faut un nom de personnes au pluriel.
\item \eng{donation in kind} — la précision « not money » et l'exemple des sacs de riz indiquent un don en nature.
\item \eng{proceeds} — les recettes de la vente aux enchères ; attention au pluriel.
\item \eng{pledged} — le ministre a \emph{promis} la somme ; la fin de la phrase (« has not paid yet ») montre que l'argent n'est pas encore versé, donc \eng{donated} serait faux.
\end{enumerate}
\end{exoresolu}

\begin{exercice}[À vous de jouer]
\begin{enumerate}
\item Traduisez en anglais : « Les élèves ont organisé une marche sponsorisée pour collecter des fonds pour la bibliothèque. »
\item Traduisez en anglais : « Une bienfaitrice anonyme a fait don d'ordinateurs à l'école. »
\item Choisissez la bonne préposition : \eng{They raised money (for / to) a new borehole.} / \eng{She donated books (for / to) the school library.}
\item Corrigez l'erreur : \emph{The proceed of the gala was given to the hospital.}
\item Trouvez le mot de la même famille : \eng{to donate} (nom : \ldots), \eng{to contribute} (nom : \ldots), \eng{to organise} (nom de personne : \ldots).
\end{enumerate}
\end{exercice}

\begin{corrige}[Exercice « À vous de jouer »]
\begin{enumerate}
\item \eng{The students organised a sponsored walk to raise funds for the library.}
\item \eng{An anonymous benefactor donated computers to the school.} (ou \eng{made a donation of computers})
\item \eng{for} / \eng{to} : on collecte \textbf{pour} un projet, on donne \textbf{à} une organisation.
\item \eng{The \textbf{proceeds} of the gala \textbf{were} given to the hospital.} (\eng{proceeds} est toujours pluriel.)
\item \eng{donation}, \eng{contribution}, \eng{organiser}.
\end{enumerate}
\end{corrige}

\begin{aretenir}[L'essentiel du champ lexical 1]
\begin{itemize}
\item Quatre familles de mots : \textbf{people} (donor, volunteer, sponsor, benefactor\ldots), \textbf{actions} (to raise funds, to pledge, to donate\ldots), \textbf{events} (gala, raffle, sponsored walk, auction\ldots), \textbf{results} (proceeds, target, amount raised, donation in kind\ldots).
\item \eng{To raise} (raise – raised – raised) est transitif ; \eng{to rise} (rise – rose – risen) est intransitif.
\item Prépositions fixes : \eng{raise funds \textbf{for}}, \eng{donate money \textbf{to}}, \eng{contribute \textbf{to}}, \eng{participate \textbf{in}}.
\item \eng{Proceeds} est toujours au pluriel ; \eng{money} est indénombrable.
\end{itemize}
\end{aretenir}

\coursec{Champ lexical 2 : Support groups and community resources/services}

Dans la section précédente, nous avons découvert le vocabulaire des \emph{fund-raising activities} : les activités de collecte de fonds (\eng{a charity concert}, \eng{to raise money}, \eng{a donation}...). Mais une question se pose naturellement : \textbf{qui} collecte ces fonds, et \textbf{pour quoi faire} ? La réponse tient en deux expressions : les \cle{support groups} (groupes de soutien) et les \cle{community resources and services} (ressources et services communautaires). C'est tout l'objet de cette section : apprendre à nommer les groupes qui agissent, les ressources qu'ils acquièrent, les services qu'ils fournissent et les personnes qu'ils aident.

\courssub{Observation : un texte pour découvrir le lexique}

Lisons d'abord cet article de presse original. Soulignez mentalement tous les mots qui désignent des groupes, des ressources, des services et des bénéficiaires.

\begin{textebox}[A community pulls together — The Bamenda Herald]
When the floods destroyed several houses in the Mile 4 neighbourhood last September, the residents did not wait for outside help. The local \textbf{welfare committee} immediately set up a \textbf{support group} for the flood victims. Within two weeks, the group had collected blankets, food and cooking pots, and had found temporary shelter for twelve families.

But the committee wanted to do more. With the assistance of a \textbf{charitable organization} called Hope for All, it launched a campaign to \textbf{acquire} new \textbf{community resources}. The first project was a \textbf{borehole}: until then, people had to walk three kilometres to the nearest \textbf{water point}. The second project was the renovation of the \textbf{health centre}, which had no beds and very little \textbf{equipment}. An international \textbf{NGO} agreed to \textbf{equip} the centre and to \textbf{supply} medicines for one year.

Other groups joined the effort. The \textbf{PTA} of the government primary school raised money to repair a \textbf{classroom block} and buy new \textbf{furniture}. The \textbf{alumni association} of the secondary school offered to \textbf{maintain} the school \textbf{library}. A women's \textbf{cooperative} provided free meals to \textbf{orphans} and other \textbf{vulnerable people}, while a \textbf{self-help group} of young volunteers took charge of \textbf{sanitation} and cleaned the market square every Saturday.

“We cannot do everything,” admits Mrs Ngum, the chairperson of the welfare committee, “but we can \textbf{cater for the needy} in our own community. Our next step is to \textbf{advocate for} better \textbf{road maintenance} and a regular \textbf{water supply}. The whole \textbf{community at large} will benefit.”
\end{textebox}

\textbf{Glossaire (les mots difficiles) :}

\begin{center}
\begin{tabularx}{\linewidth}{|l|l|X|}
\hline
\rowcolor{popblueL} \textbf{Mot anglais} & \textbf{Nature} & \textbf{Traduction} \\
\hline
flood & noun & inondation \\
\hline
shelter & noun & abri, hébergement \\
\hline
borehole & noun & forage (puits équipé d'une pompe) \\
\hline
chairperson & noun & président(e) \\
\hline
the needy & plural noun & les nécessiteux, les démunis \\
\hline
to advocate for & verb & plaider en faveur de, militer pour \\
\hline
\end{tabularx}
\end{center}

\textbf{Questions de compréhension (répondez en anglais, par des phrases complètes) :}

\begin{enumerate}
\item Who set up a support group after the floods?
\item What were the two main projects of the welfare committee?
\item Which group raised money to repair the classroom block?
\item What does the women's cooperative do for orphans?
\item What will Mrs Ngum advocate for next?
\end{enumerate}

\begin{corrige}[Questions de compréhension]
\begin{enumerate}
\item \eng{The local welfare committee set up a support group after the floods.}
\item \eng{The two main projects were a borehole and the renovation of the health centre.}
\item \eng{The PTA of the government primary school raised money to repair the classroom block.}
\item \eng{The women's cooperative provides free meals to orphans and other vulnerable people.}
\item \eng{She will advocate for better road maintenance and a regular water supply.}
\end{enumerate}
\end{corrige}

\courssub{Types de groupes : qui agit ?}

\begin{definition}[Support group and related words]
Un \cle{support group} est un groupe de personnes qui s'unissent pour s'entraider ou pour aider les autres (moralement, matériellement ou financièrement). On distingue plusieurs types de groupes :
\begin{itemize}
\item \eng{a support group} : un groupe de soutien (ex. \eng{a support group for flood victims}) ;
\item \eng{a self-help group} : un groupe d'entraide, dont les membres s'aident mutuellement (ex. une \emph{tontine} de quartier) ;
\item \eng{a charity} / \eng{a charitable organization} : une œuvre de charité, une organisation caritative ;
\item \eng{an NGO} (\emph{non-governmental organization}) : une ONG ;
\item \eng{an association} : une association ;
\item \eng{a cooperative} : une coopérative (groupement de production ou d'achat en commun) ;
\item \eng{a welfare committee} : un comité de bien-être social (souvent au niveau d'un quartier ou d'un village) ;
\item \eng{a PTA} (\emph{Parent-Teacher Association}) : une association de parents d'élèves et d'enseignants ;
\item \eng{an alumni association} : une association d'anciens élèves.
\end{itemize}
\end{definition}

\begin{definition}[What is an NGO?]
An \cle{NGO} is a \emph{non-governmental organization}: an organization that is \textbf{independent from the government} and that works for \textbf{social or humanitarian aims} (health, education, the environment, human rights). En français : \emph{une organisation non gouvernementale (ONG)}. Notez la prononciation des lettres : « ène-dji-o ».
\end{definition}

\begin{attention}[Pluriel de NGO et majuscules]
Le pluriel de \eng{NGO} est \eng{NGOs} (et non \emph{NGO's} : l'apostrophe ne marque jamais le pluriel en anglais). Les sigles s'écrivent en majuscules : \eng{NGO}, \eng{PTA}, mais les mots développés restent en minuscules : \eng{a non-governmental organization}, \eng{a parent-teacher association}. Attention aussi à l'article : \eng{an NGO} (« an » devant le son voyelle « ène »), \eng{a PTA} (« a » devant le son consonne « pi »).
\end{attention}

\courssub{Community resources : les ressources communautaires}

Les \cle{community resources} sont les biens et infrastructures matériels dont dispose une communauté. Voici le vocabulaire essentiel :

\begin{center}
\begin{tabularx}{\linewidth}{|X|X|X|}
\hline
\rowcolor{popblueL} \textbf{English} & \textbf{Français} & \textbf{Exemple en contexte} \\
\hline
a borehole & un forage & \eng{The village dug a borehole.} \\
\hline
a water point & un point d'eau & \eng{The nearest water point is far.} \\
\hline
a health centre & un centre de santé & \eng{The health centre has no doctor.} \\
\hline
a classroom block & un bloc de salles de classe & \eng{They built a new classroom block.} \\
\hline
a library & une bibliothèque & \eng{The library needs more books.} \\
\hline
a community hall & une salle communautaire / des fêtes & \eng{Meetings are held in the community hall.} \\
\hline
electricity & l'électricité (indénombrable) & \eng{The school has no electricity.} \\
\hline
equipment & l'équipement, le matériel (indénombrable) & \eng{The equipment is modern.} \\
\hline
furniture & les meubles, le mobilier (indénombrable) & \eng{The furniture was donated.} \\
\hline
\end{tabularx}
\end{center}

\begin{attention}[Trois indénombrables piégeux]
\eng{Equipment}, \eng{furniture} et \eng{electricity} sont \textbf{indénombrables} en anglais : jamais de \emph{-s}, jamais de \emph{a/an}, verbe au singulier.
\begin{itemize}
\item \eng{The equipment IS new.} « L'équipement est neuf. » (et non \emph{equipments are})
\item \eng{The furniture WAS donated.} « Les meubles ont été donnés. » (et non \emph{furnitures})
\item Pour compter, on dit \eng{a piece of equipment} « un appareil / un équipement », \eng{a piece of furniture} « un meuble ».
\end{itemize}
Même logique pour \eng{information}, \eng{advice} et \eng{news}, également indénombrables : \eng{a piece of advice} « un conseil ».
\end{attention}

\courssub{Community services : les services communautaires}

À ne pas confondre avec les ressources (les \emph{choses}) : les \cle{community services} sont les \emph{prestations} dont bénéficie la population.

\begin{center}
\begin{tabularx}{\linewidth}{|X|X|X|}
\hline
\rowcolor{popgreenL} \textbf{English} & \textbf{Français} & \textbf{Exemple en contexte} \\
\hline
health care & les soins de santé & \eng{Health care is free for children.} \\
\hline
education & l'éducation, l'instruction & \eng{Every child deserves an education.} \\
\hline
counselling & le conseil, l'accompagnement psychologique & \eng{The centre offers counselling to victims.} \\
\hline
sanitation & l'assainissement, l'hygiène publique & \eng{Poor sanitation causes diseases.} \\
\hline
water supply & l'approvisionnement en eau & \eng{The water supply is irregular.} \\
\hline
security & la sécurité & \eng{The vigilante group ensures security.} \\
\hline
road maintenance & l'entretien des routes & \eng{Road maintenance is the council's duty.} \\
\hline
\end{tabularx}
\end{center}

\begin{exemplebox}[Ressource ou service ?]
Comparez :
\begin{itemize}
\item \eng{The NGO built a \textbf{health centre}.} (ressource : un bâtiment) « L'ONG a construit un centre de santé. »
\item \eng{The NGO provides \textbf{health care}.} (service : des soins) « L'ONG fournit des soins de santé. »
\item \eng{The PTA repaired a \textbf{classroom block}.} (ressource) « L'association de parents a réparé un bloc de classes. »
\item \eng{The state must guarantee \textbf{education}.} (service) « L'État doit garantir l'éducation. »
\end{itemize}
\end{exemplebox}

\courssub{Les actions des groupes : les verbes clés}

\begin{propriete}[Verbes d'aide et de fourniture]
\begin{itemize}
\item \eng{to support somebody} : soutenir quelqu'un (jamais de préposition après \emph{support}). \eng{They support the orphans.} « Ils soutiennent les orphelins. »
\item \eng{to assist somebody (in doing something)} : aider, assister (plus formel que \emph{help}). \eng{The NGO assisted the committee in organizing the campaign.} « L'ONG a aidé le comité à organiser la campagne. »
\item \eng{to help out} : donner un coup de main (familier ; phrasal verb séparable). \eng{Volunteers helped out at the centre.} / \eng{They helped the victims out.} « Des bénévoles ont donné un coup de main au centre. / Ils ont dépanné les victimes. »
\item \eng{to provide something} / \eng{to provide somebody WITH something} : fournir. \eng{The group provided food.} / \eng{The group provided the victims with food.}
\item \eng{to supply something (to somebody)} : approvisionner, fournir (souvent de façon régulière). \eng{The NGO supplies medicines to the clinic.}
\item \eng{to acquire} : acquérir, se procurer (formel). \eng{The school acquired new computers.} « L'école a acquis de nouveaux ordinateurs. »
\item \eng{to maintain} : entretenir. \eng{The alumni maintain the library.} « Les anciens élèves entretiennent la bibliothèque. »
\item \eng{to equip somebody/something WITH something} : équiper de. \eng{They equipped the centre with beds.}
\item \eng{to cater FOR the needy} : subvenir aux besoins des démunis (préposition \emph{for} obligatoire).
\item \eng{to advocate FOR something} : plaider pour, militer en faveur de. \eng{They advocate for clean water.}
\end{itemize}
\end{propriete}

\begin{exemplebox}[Exemples traduits]
\begin{itemize}
\item \eng{A support group for flood victims provided blankets and food.} « Un groupe de soutien aux victimes des inondations a fourni des couvertures et de la nourriture. »
\item \eng{The NGO equipped the health centre with new beds.} « L'ONG a équipé le centre de santé de lits neufs. »
\item \eng{The cooperative caters for the needy in the neighbourhood.} « La coopérative subvient aux besoins des démunis du quartier. »
\item \eng{The welfare committee advocates for better sanitation.} « Le comité de bien-être plaide pour un meilleur assainissement. »
\item \eng{The alumni association helps out during examinations.} « L'association des anciens élèves donne un coup de main pendant les examens. »
\end{itemize}
\end{exemplebox}

\begin{attention}[Prépositions : les pièges des francophones]
\begin{itemize}
\item \eng{to provide somebody \textbf{with} something} : ne dites jamais \emph{to provide somebody something} sur le modèle français « fournir quelqu'un quelque chose ». Mais \eng{to provide something \textbf{for} somebody} est correct : \eng{They provided food for the victims.}
\item \eng{to cater \textbf{for}} et \eng{to advocate \textbf{for}} : la préposition \emph{for} est obligatoire, même si le français dit « subvenir aux besoins » et « plaider pour ».
\item \eng{to listen \textbf{to}}, \eng{to wait \textbf{for}}, \eng{to depend \textbf{on}} : même logique. En revanche, \eng{to support}, \eng{to assist}, \eng{to help} se construisent \textbf{sans} préposition : ne dites pas \emph{to help to somebody} ni \emph{to support to the orphans}.
\end{itemize}
\end{attention}

\courssub{Les bénéficiaires : qui reçoit l'aide ?}

\begin{definition}[Beneficiaries]
Les \cle{beneficiaries} (bénéficiaires) sont les personnes qui reçoivent l'aide :
\begin{itemize}
\item \eng{the needy} : les nécessiteux, les démunis ;
\item \eng{the underprivileged} : les défavorisés ;
\item \eng{vulnerable people} : les personnes vulnérables ;
\item \eng{orphans} : les orphelins ;
\item \eng{flood victims} : les victimes d'inondations ;
\item \eng{the disabled} : les personnes handicapées ;
\item \eng{the community at large} : la communauté dans son ensemble, l'ensemble de la population.
\end{itemize}
\end{definition}

\begin{attention}[« The needy ARE... » : l'adjectif substantivé]
\eng{The needy}, \eng{the disabled}, \eng{the underprivileged}, \eng{the poor}, \eng{the rich} sont des \textbf{noms collectifs au pluriel} formés par \emph{the + adjectif} (on parle d'adjectif substantivé). Ils désignent toute une catégorie de personnes et prennent un \textbf{verbe au pluriel}. On ne leur ajoute \textbf{jamais} de \emph{-s}.
\begin{itemize}
\item \eng{The needy ARE our priority.} « Les démunis sont notre priorité. » (et non \emph{the needies are}, ni \emph{the needy is})
\item \eng{The disabled NEED special equipment.} « Les personnes handicapées ont besoin d'un équipement spécial. »
\item Pour parler d'\textbf{une seule} personne, il faut un nom : \eng{a disabled person}, \eng{a needy family}.
\end{itemize}
\end{attention}

\courssub{Schéma récapitulatif : le circuit de la solidarité}

\begin{popfigure}
\begin{center}
\begin{tikzpicture}[
  node distance=0.9cm,
  box/.style={draw=popblue, fill=popblueL, rounded corners, align=center, text width=3.6cm, minimum height=1.1cm, font=\small},
  box2/.style={draw=popgreen, fill=popgreenL, rounded corners, align=center, text width=4.2cm, minimum height=1.1cm, font=\small},
  box3/.style={draw=poporange, fill=poporangeL, rounded corners, align=center, text width=4.6cm, minimum height=1.1cm, font=\small},
  box4/.style={draw=poppurple, fill=poppurpleL, rounded corners, align=center, text width=4.6cm, minimum height=1.1cm, font=\small},
  fleche/.style={-{Stealth[length=3mm]}, thick, popdark}
]
\node[box, align=center] (groupe){\textbf{Support group}\\ NGO, PTA, cooperative, welfare committee...};
\node[box2, below=of groupe, align=center] (actions){\textbf{raises funds / provides services}\\ to raise, to equip, to maintain, to cater for...};
\node[box3, below=of actions, align=center] (ressources){\textbf{community resources}\\ borehole, health centre, classroom block, library...};
\node[box4, below=of ressources, align=center] (benef){\textbf{beneficiaries}\\ the needy, orphans, the disabled, the community at large};
\draw[fleche] (groupe) -- (actions);
\draw[fleche] (actions) -- (ressources);
\draw[fleche] (ressources) -- (benef);
\draw[fleche, dashed, popmuted] (benef.west) to[bend right=45] node[left, font=\scriptsize, align=center]{they may join\\ the group too!} (groupe.west);
\end{tikzpicture}
\end{center}
\legende{Le circuit de la solidarité communautaire : un groupe de soutien collecte des fonds et fournit des services, ce qui permet d'acquérir des ressources communautaires au profit des bénéficiaires.}
\end{popfigure}

\courssub{Exercice résolu}

\begin{exoresolu}[Complétez avec le mot qui convient]
Complétez chaque phrase avec un mot de la liste : \eng{borehole} — \eng{counselling} — \eng{the needy} — \eng{equipment} — \eng{advocate for} — \eng{alumni association} — \eng{sanitation}.

1. The \_\_\_\_\_\_ of our school donated fifty desks to the library.
2. The village needs a new \_\_\_\_\_\_ because the old water point is dry.
3. The support group offers free \_\_\_\_\_\_ to traumatized flood victims.
4. We must \_\_\_\_\_\_ better road maintenance in our subdivision.
5. The hospital received new medical \_\_\_\_\_\_ from a foreign NGO.
6. Volunteers clean the gutters every month to improve \_\_\_\_\_\_.
7. This charity caters for \_\_\_\_\_\_ in the whole region.

\tcblower
\textbf{Solution détaillée :}

1. \eng{alumni association} — seul ce groupe d'anciens élèves peut logiquement donner des tables à la bibliothèque de « notre » école. « L'association des anciens élèves de notre école a donné cinquante tables à la bibliothèque. »
2. \eng{borehole} — un forage remplace un point d'eau asséché. « Le village a besoin d'un nouveau forage car l'ancien point d'eau est à sec. »
3. \eng{counselling} — on offre un accompagnement psychologique à des victimes traumatisées. « Le groupe de soutien offre un accompagnement gratuit aux victimes d'inondations traumatisées. »
4. \eng{advocate for} — « plaider pour » un meilleur entretien des routes (verbe + préposition \emph{for}). « Nous devons plaider pour un meilleur entretien des routes dans notre arrondissement. »
5. \eng{equipment} — indénombrable : pas de \emph{-s}. « L'hôpital a reçu du matériel médical neuf d'une ONG étrangère. »
6. \eng{sanitation} — nettoyer les caniveaux améliore l'assainissement. « Des bénévoles nettoient les caniveaux chaque mois pour améliorer l'assainissement. »
7. \eng{the needy} — adjectif substantivé, sans \emph{-s} : « Cette œuvre de charité subvient aux besoins des démunis dans toute la région. »
\end{exoresolu}

\courssub{Pour aller plus loin : culture anglophone}

\begin{savaistu}[Charity shops and bake sales]
Au Royaume-Uni, les \eng{charity shops} sont des magasins qui vendent des objets d'occasion (vêtements, livres, vaisselle) donnés par le public : tous les bénéfices financent des œuvres caritatives. On en trouve dans presque chaque rue commerçante ! Aux États-Unis, les \eng{PTA} jouent un rôle central dans la vie scolaire : ils organisent notamment des \eng{bake sales}, des ventes de gâteaux faits maison par les parents, pour financer les sorties scolaires, les bibliothèques ou le matériel de sport. La tradition du bénévolat communautaire y est si forte que beaucoup d'universités américaines examinent les activités de volontariat (\eng{volunteer work}) des candidats à l'admission.
\end{savaistu}

\begin{aretenir}[L'essentiel de la section]
\begin{itemize}
\item Les groupes : \eng{support group}, \eng{self-help group}, \eng{charity}, \eng{NGO}, \eng{cooperative}, \eng{welfare committee}, \eng{PTA}, \eng{alumni association}.
\item Les ressources (choses) : \eng{borehole}, \eng{water point}, \eng{health centre}, \eng{classroom block}, \eng{community hall}, \eng{equipment}, \eng{furniture} — attention, ces deux derniers sont indénombrables.
\item Les services (prestations) : \eng{health care}, \eng{education}, \eng{counselling}, \eng{sanitation}, \eng{water supply}, \eng{road maintenance}.
\item Les verbes et leurs prépositions : \eng{to provide sb WITH sth}, \eng{to equip sth WITH sth}, \eng{to cater FOR}, \eng{to advocate FOR}, mais \eng{to support/assist/help} sans préposition.
\item Les bénéficiaires : \eng{the needy}, \eng{the disabled}, \eng{the underprivileged} = adjectifs substantivés, verbe au pluriel, jamais de \emph{-s}.
\end{itemize}
\end{aretenir}

Nous savons désormais nommer les acteurs de la solidarité, ce qu'ils acquièrent et ceux qu'ils aident. Dans la prochaine section, nous verrons comment ces mots se combinent naturellement entre eux : les \emph{collocations}, les expressions idiomatiques et les familles de mots.

\coursec{Collocations, expressions idiomatiques et mots de la même famille}

Dans la section précédente, nous avons découvert le vocabulaire des \eng{support groups} et des \eng{community resources and services} : banques alimentaires, dispensaires, bibliothèques, centres de jeunesse. Mais ces ressources n'existent que parce que des gens agissent ensemble : ils \emph{collectent} des fonds, ils \emph{font} des dons, ils \emph{lancent} des campagnes. Or, en anglais, chaque nom s'associe à un verbe précis, et pas à un autre : on dit \eng{make a donation} et jamais \eng{do a donation}. Ces associations obligées s'appellent des \cle{collocations}. C'est précisément ce que nous allons étudier maintenant, avec les expressions idiomatiques et les familles de mots qui permettront à votre vocabulaire de « s'étendre » naturellement.

\courssub{Observation : un texte pour découvrir les collocations}

\begin{textebox}[“Old Students Give Back to Their School” — The Buea Herald (article original)]
The alumni association of Government High School Buea has launched an appeal to raise funds for a new computer room. At a fundraising event held last Saturday, the president of the association announced that the school needed at least two million francs CFA. “Every little helps,” she said. “Even five hundred francs can make a difference. We are asking everyone to dig deep into their pockets.”

A local businessman made a generous donation of five hundred thousand francs, and several former students promised to make regular contributions until the target is reached. “We simply want to give back to the community that educated us,” explained one of them, who now works for a charitable organisation in Douala.

The campaign is not only about money. Volunteers have offered to do their bit by painting the new room and installing the furniture. The head teacher thanked everyone warmly: “This worthy cause meets the real needs of our students. When people lend a helping hand, a whole school can be transformed.”

The organisers hope to reach their target before the end of the school year. A closing gala will be held in June, with music, food and a lottery, to celebrate the generosity of all donors.
\end{textebox}

\textbf{Glossaire :} \emph{alumni} (nom pluriel, anciens élèves) ; \emph{appeal} (nom, appel) ; \emph{target} (nom, objectif à atteindre) ; \emph{lottery} (nom, tombola) ; \emph{head teacher} (nom, proviseur/directeur) ; \emph{former students} (nom, anciens élèves).

\textbf{Questions de compréhension :}
\begin{enumerate}
\item What is the alumni association raising funds for?
\item How much did the local businessman donate?
\item Name two things volunteers offered to do besides giving money.
\item When will the closing gala take place?
\end{enumerate}

Relevez dans le texte les groupes \emph{verbe + nom} : \eng{launch an appeal}, \eng{raise funds}, \eng{make a donation}, \eng{reach a target}, \eng{hold a gala}, \eng{lend a helping hand}, \eng{meet the needs of}. Aucun de ces verbes n'est interchangeable : on ne peut pas dire \eng{do a donation} ni \eng{make funds}. C'est le cœur de cette section.

\courssub{Collocations verbe + nom}

\begin{definition}[Collocation]
Une \cle{collocation} est une association habituelle et quasi obligatoire de deux mots : un verbe et son nom, ou un adjectif et son nom. En anglais, on ne choisit pas le verbe librement : on dit \eng{make a donation} mais \eng{do voluntary work}, \eng{raise funds} mais \eng{launch a campaign}. Ces paires s'apprennent comme des blocs indissociables.
\end{definition}

\begin{propriete}[Les grands verbes du fund-raising]
\begin{center}
\begin{tabularx}{\linewidth}{|X|X|X|}\hline
\rowcolor{popblueL} Collocation anglaise & Traduction française & Exemple en contexte \\ \hline
\eng{raise funds / raise money} & collecter des fonds & \eng{They raised funds for the clinic.} \\ \hline
\eng{raise awareness} & sensibiliser & \eng{The campaign raised awareness about malaria.} \\ \hline
\eng{make a donation / a contribution} & faire un don / une contribution & \eng{She made a donation of 10,000 FCFA.} \\ \hline
\eng{launch an appeal / a campaign} & lancer un appel / une campagne & \eng{The NGO launched an appeal for volunteers.} \\ \hline
\eng{hold / organise a gala / an event} & organiser un gala / un événement & \eng{We held a gala at the town hall.} \\ \hline
\eng{reach a target} & atteindre un objectif & \eng{We reached our target in three weeks.} \\ \hline
\eng{meet the needs of} & répondre aux besoins de & \eng{This project meets the needs of orphans.} \\ \hline
\eng{lend a (helping) hand} & donner un coup de main & \eng{Neighbours lent a hand after the flood.} \\ \hline
\eng{do voluntary work} & faire du bénévolat & \eng{He does voluntary work on Saturdays.} \\ \hline
\eng{support a cause} & soutenir une cause & \eng{Many artists support this cause.} \\ \hline
\end{tabularx}
\end{center}
\end{propriete}

\begin{attention}[Les verbes pièges : make, do, raise]
\begin{itemize}
\item On dit \eng{make a donation}, jamais \eng{do a donation}. En revanche : \eng{do voluntary work}, \eng{do one's bit}. Il n'y a pas de logique : il faut mémoriser chaque bloc.
\item \eng{To raise} (collecter, élever — verbe transitif, toujours suivi d'un complément) ne se confond pas avec \eng{to rise} (s'élever, augmenter — verbe intransitif, irrégulier : \eng{rise, rose, risen}) : \eng{They raised money.} / \eng{Prices rise.}
\item Attention au calque du français « organiser une collecte » : on préfère \eng{to organise a fundraising event} ou \eng{to hold a collection}.
\item \eng{To collect funds} existe, mais \eng{to raise funds} est la collocation la plus fréquente et la plus attendue à l'examen.
\end{itemize}
\end{attention}

\courssub{Collocations adjectif + nom}

L'adjectif, lui aussi, a ses partenaires privilégiés. Un don n'est pas \eng{big} mais \eng{generous} ; une cause n'est pas \eng{good} mais \eng{worthy}.

\begin{center}
\begin{tabularx}{\linewidth}{|X|X|X|}\hline
\rowcolor{popblueL} Adjectif + nom & Traduction & Remarque \\ \hline
\eng{a generous donation} & un don généreux & pas \eng{a big donation} \\ \hline
\eng{a charitable organisation} & une organisation caritative & \eng{charitable} = adjectif de \eng{charity} \\ \hline
\eng{voluntary work} & travail bénévole & forme britannique standard ; \eng{volunteer work} est américain \\ \hline
\eng{community service} & service à la communauté & aussi « travaux d'intérêt général » \\ \hline
\eng{a worthy cause} & une cause noble / qui en vaut la peine & \eng{worthy} = digne, qui mérite \\ \hline
\eng{a fundraising event} & un événement de collecte de fonds & orthographe : voir plus bas \\ \hline
\eng{a non-profit organisation} & une organisation à but non lucratif & à ne pas confondre avec \eng{NGO} \\ \hline
\eng{a non-governmental organisation (NGO)} & une organisation non gouvernementale (ONG) & \eng{NGO} = \eng{non-governmental organisation} \\ \hline
\eng{a regular contribution} & une contribution régulière & utile pour les cotisations \\ \hline
\end{tabularx}
\end{center}

\begin{attention}[Non-profit ou NGO ? Ne confondez pas]
Une \eng{non-profit organisation} est une organisation \emph{à but non lucratif} (elle ne distribue pas de bénéfices). Une \eng{NGO} (\eng{non-governmental organisation}) est une organisation \emph{non gouvernementale} (indépendante de l'État). Beaucoup d'ONG sont aussi à but non lucratif, mais les deux mots ne sont pas synonymes : \eng{NGO} ne veut pas dire « à but non lucratif ».
\end{attention}

\begin{exemplebox}[Exemples traduits]
\eng{Every little helps: even 500 FCFA can make a difference.} « Tout petit don compte : même 500 FCFA peuvent faire la différence. »

\eng{The alumni decided to give back to the community.} « Les anciens élèves ont décidé de rendre à la communauté ce qu'elle leur avait donné. »

\eng{She works for a charitable organisation that supports a worthy cause.} « Elle travaille pour une organisation caritative qui soutient une noble cause. »
\end{exemplebox}

\courssub{Expressions idiomatiques de la générosité}

\begin{definition}[Expression idiomatique]
Une \cle{expression idiomatique} est un groupe de mots dont le sens global ne se devine pas à partir des mots pris un à un. Elle ne se traduit jamais mot à mot : il faut chercher l'équivalent français.
\end{definition}

\begin{center}
\begin{tabularx}{\linewidth}{|X|X|X|}\hline
\rowcolor{popblueL} Idiom anglais & Sens réel & Équivalent français \\ \hline
\eng{to lend a helping hand} & aider concrètement & donner un coup de main \\ \hline
\eng{to give back to the community} & rendre à la communauté ce qu'on a reçu & s'acquitter envers la société \\ \hline
\eng{every little helps} & chaque petit don compte & les petits ruisseaux font les grandes rivières \\ \hline
\eng{charity begins at home} & on doit d'abord aider les siens & charité bien ordonnée commence par soi-même \\ \hline
\eng{to dig deep (into one's pockets)} & donner généreusement, faire un effort financier & se saigner (aux quatre veines) \\ \hline
\eng{to do one's bit} & apporter sa petite contribution & mettre la main à la pâte, faire sa part \\ \hline
\end{tabularx}
\end{center}

\begin{exemplebox}[Les idioms en contexte]
\eng{When the school roof collapsed, the whole neighbourhood lent a helping hand.} « Quand le toit de l'école s'est effondré, tout le quartier a donné un coup de main. »

\eng{Come on, dig deep! The orphanage needs new beds.} « Allez, faites un effort ! L'orphelinat a besoin de lits neufs. »

\eng{I cannot give much, but I want to do my bit.} « Je ne peux pas donner beaucoup, mais je veux faire ma part. »
\end{exemplebox}

\begin{attention}[« Charity begins at home » : un proverbe, pas une phrase ordinaire]
\eng{Charity begins at home} est un \emph{proverbe} figé : il signifie que l'on doit d'abord prendre soin de sa famille et de ses proches avant d'aider les étrangers. Ne le traduisez jamais mot à mot (« la charité commence à la maison ») dans une rédaction : utilisez l'équivalent français « charité bien ordonnée commence par soi-même », ou reformulez : \eng{we should first help the people around us}. De même, \eng{to dig deep} ne veut pas dire « creuser profondément » dans ce contexte !
\end{attention}

\begin{savaistu}[La culture du don dans le monde anglophone]
En Grande-Bretagne, les \eng{charity shops} (magasins caritatifs) vendent des vêtements et des objets d'occasion donnés par le public, au profit d'œuvres de charité. Aux États-Unis, les écoles et les églises organisent des \eng{bake sales} (ventes de gâteaux) et des \eng{jumble sales} (ventes de charité) pour financer leurs projets. Au Cameroun anglophone, les \eng{alumni associations} comme les associations d'anciens élèves de Buea jouent exactement ce rôle : elles collectent des fonds pour les salles de classe, les laboratoires et les bourses. Le mot \eng{fundraising} est donc bien plus qu'un mot de vocabulaire : c'est une pratique sociale quotidienne.
\end{savaistu}

\courssub{Les familles de mots : un verbe, toute une famille}

Observez : \eng{donate} (verbe), \eng{donation} (action), \eng{donor} (personne). L'anglais construit des familles entières à partir d'une racine, avec des suffixes réguliers. Connaître ces familles, c'est multiplier votre vocabulaire par trois sans effort.

\begin{propriete}[Formation des noms d'action et des noms de personne]
Pour former le \emph{nom d'action} à partir du verbe, on ajoute souvent \cle{-tion} (\eng{donate → donation}, \eng{contribute → contribution}) ou \cle{-ation} (\eng{organise → organisation}). Le nom de \emph{personne} prend \cle{-er} ou \cle{-or} : \eng{organiser}, \eng{donor}, \eng{contributor}. Attention à la chute du \eng{-e} final : \eng{donate → donation} (et non \eng{donatetion}). Cas particulier : \eng{volunteer} est à la fois le verbe et le nom de personne (suffixe \eng{-eer}, comme \eng{engineer}).
\end{propriete}

\begin{center}
\begin{tabular}{|l|l|l|}\hline
\rowcolor{popblueL} Verbe & Nom (action) & Nom (personne) / adjectif \\ \hline
\eng{to donate} & \eng{donation} & \eng{donor} \\ \hline
\eng{to contribute} & \eng{contribution} & \eng{contributor} \\ \hline
\eng{to volunteer} & \eng{volunteering} & \eng{volunteer} ; \eng{voluntary} (adj.), \eng{voluntarily} (adv.) \\ \hline
\eng{to organise} & \eng{organisation} & \eng{organiser} \\ \hline
\eng{to assist} & \eng{assistance} & \eng{assistant} \\ \hline
\eng{to provide} & \eng{provision} & \eng{provider} \\ \hline
\eng{---} & \eng{charity} & \eng{charitable} (adj.) \\ \hline
\end{tabular}
\end{center}

\begin{popfigure}
\begin{tikzpicture}[
  grow=down,
  level distance=1.5cm,
  level 1/.style={sibling distance=4.2cm},
  every node/.style={draw, rounded corners, align=center, font=\small, inner sep=4pt},
  edge from parent/.style={draw=popmuted, -{Stealth[length=2mm]}}
]
\node[fill=popblueL, align=center]{\textbf{volunteer}\\ \footnotesize verbe et nom : se porter volontaire / bénévole}
  child { node[fill=popgreenL, align=center]{\textbf{voluntary}\\ \footnotesize adjectif : bénévole} }
  child { node[fill=poporangeL, align=center]{\textbf{voluntarily}\\ \footnotesize adverbe : bénévolement} }
  child { node[fill=poppurpleL, align=center]{\textbf{volunteering}\\ \footnotesize nom : le bénévolat} };
\end{tikzpicture}
\legende{Arbre de la famille de \eng{volunteer} : un seul mot-racine donne un verbe, un nom de personne, un adjectif et un adverbe.}
\end{popfigure}

\begin{exemplebox}[Une famille, trois phrases]
\eng{Many donors gave money last year.} « De nombreux donateurs ont donné de l'argent l'année dernière. » (personne)

\eng{Their donations built the new well.} « Leurs dons ont permis de construire le nouveau puits. » (action)

\eng{They donate regularly to the village fund.} « Ils donnent régulièrement à la caisse du village. » (verbe)
\end{exemplebox}

\courssub{Prononciation et orthographe : les pièges à déjouer}

\begin{propriete}[Prononciation des mots clés (notation française)]
\begin{itemize}
\item \eng{charity} : « tchar-i-ti » (\eng{ch} = « tch » comme dans \eng{match})
\item \eng{volunteer} : « vo-leun-tier » (accent sur la dernière syllabe !)
\item \eng{donation} : « do-né-cheun » (le \eng{-tion} se prononce « cheun »)
\item \eng{organisation} : « or-ga-naï-zé-cheun »
\item \eng{contribution} : « kon-tri-biou-cheun »
\end{itemize}
Règle générale : la terminaison \eng{-tion} se prononce toujours « cheun », jamais « ti-on ».
\end{propriete}

\begin{attention}[Orthographe : fundraising, pas « found-raising »]
\begin{itemize}
\item On écrit \eng{fundraising} (un seul mot, le plus courant) ou \eng{fund-raising} (avec trait d'union). Jamais \eng{found-raising} : \eng{fund} (fonds) et \eng{found} (fonder ; aussi prétérit de \eng{find}) sont deux mots différents.
\item \eng{organisation} est la graphie britannique traditionnelle ; \eng{organization} (avec \eng{z}) est la graphie américaine, aussi acceptée en Grande-Bretagne. Les deux existent, mais ne mélangez pas les deux graphies dans une même copie. Dans cette leçon, nous utilisons la graphie britannique en \eng{-isation}.
\item \eng{voluntary} s'écrit avec \eng{-ary} (comme \eng{library}), pas \eng{-ery}.
\end{itemize}
\end{attention}

\courssub{Méthode et entraînement}

\begin{methode}[Comment mémoriser les collocations durablement]
\begin{enumerate}
\item Apprenez toujours le verbe \emph{avec} son nom habituel : mémorisez le bloc \eng{make a donation}, jamais le mot \eng{donation} tout seul.
\item Classez vos blocs dans un carnet en deux colonnes : verbes d'action (\eng{raise, launch, hold, reach}) d'un côté, noms (\eng{funds, appeal, gala, target}) de l'autre, et reliez-les.
\item Rédigez une phrase \emph{personnelle} pour chaque paire, liée à votre vie : \eng{My class raised funds for our end-of-year party.}
\item Relisez vos phrases à voix haute en respectant la prononciation (« do-né-cheun », « vo-leun-tier »).
\item Testez-vous une semaine plus tard : cachez la colonne anglaise et retrouvez le bloc à partir du français.
\end{enumerate}
\end{methode}

\begin{exoresolu}[Complétez avec la bonne collocation]
Énoncé — Complétez chaque phrase avec le verbe qui convient : \eng{raise}, \eng{make}, \eng{launch}, \eng{reach}, \eng{lend}.

1. The students decided to \ldots\ a campaign to help the hospital.
2. We hope to \ldots\ our target of one million francs.
3. Could you \ldots\ me a hand with the posters?
4. The company will \ldots\ a donation to the orphanage.
5. They want to \ldots\ awareness about clean water.
\tcblower
Solution détaillée :

1. \eng{launch} — on « lance » une campagne : \eng{launch a campaign}.
2. \eng{reach} — un objectif s'« atteint » : \eng{reach a target}.
3. \eng{lend} — l'idiom est \eng{lend a (helping) hand} (« donner un coup de main »).
4. \eng{make} — la collocation obligatoire est \eng{make a donation} (jamais \eng{do}).
5. \eng{raise} — \eng{raise awareness} = « sensibiliser ».
\end{exoresolu}

\begin{exercice}[À vous de jouer]
1. Traduisez en anglais : « Les anciens élèves ont lancé un appel pour collecter des fonds, et chacun a fait sa part. »
2. Trouvez le nom d'action et le nom de personne de : \eng{contribute}, \eng{assist}, \eng{provide}.
3. Corrigez les trois erreurs : \eng{They did a big donation to the found-raising event.}
\end{exercice}

\begin{corrige}[Exercice « À vous de jouer »]
1. \eng{The alumni launched an appeal to raise funds, and everyone did their bit.}
2. \eng{contribute → contribution / contributor} ; \eng{assist → assistance / assistant} ; \eng{provide → provision / provider}.
3. \eng{They made a generous donation to the fundraising event.} (collocation \eng{make a donation} et non \eng{do} ; adjectif \eng{generous} et non \eng{big} ; orthographe \eng{fundraising} et non \eng{found-raising}.)
\end{corrige}

\begin{aretenir}[L'essentiel de la section]
\begin{itemize}
\item Une collocation est un bloc indissociable : \eng{make a donation}, \eng{raise funds}, \eng{launch a campaign}, \eng{reach a target}, \eng{hold a gala}.
\item Les adjectifs ont leurs partenaires : \eng{generous donation}, \eng{worthy cause}, \eng{voluntary work}, \eng{charitable organisation}.
\item Les idioms (\eng{lend a helping hand}, \eng{dig deep}, \eng{do one's bit}, \eng{every little helps}) ne se traduisent jamais mot à mot.
\item Familles de mots : verbe + \eng{-tion/-ation} = action ; + \eng{-er/-or} = personne ; \eng{volunteer} est à la fois verbe et nom de personne.
\item Prononciation : \eng{-tion} = « cheun » ; orthographe : \eng{fundraising}, jamais \eng{found-raising}.
\item \eng{NGO} = \eng{non-governmental organisation} (ONG), à ne pas confondre avec \eng{non-profit organisation} (à but non lucratif).
\end{itemize}
\end{aretenir}

\coursec{Faux amis et pièges des francophones}

Dans la section précédente, nous avons enrichi notre vocabulaire grâce aux collocations (\eng{to raise funds}, \eng{to make a donation}), aux expressions idiomatiques (\eng{to dig deep into one's pockets}) et aux familles de mots (\eng{to donate — a donation — a donor}). Mais attention : plus un mot anglais ressemble à un mot français, plus il est dangereux. Le domaine de la collecte de fonds et de l'entraide communautaire regorge de \cle{faux amis} (\eng{false friends}) : des mots qui se ressemblent dans les deux langues mais ne veulent pas dire la même chose. Cette section vous apprend à les repérer, à les corriger et à ne plus jamais tomber dans leurs pièges — ni à l'écrit, ni à l'oral, ni au Bac.

\courssub{Observer avant de généraliser : deux phrases, deux erreurs}

Lisez attentivement ces deux phrases entendues dans une classe de Terminale à Yaoundé :

\begin{exemplebox}[Deux phrases à corriger]
\eng{*I assisted to the fund-raising gala last Saturday.} (incorrect)\par
\eng{*The benefactors of the project received school fees.} (incorrect)
\end{exemplebox}

L'élève qui a écrit la première phrase voulait dire « J'ai assisté au gala de collecte de fonds » ; celui qui a écrit la seconde voulait dire « Les bénéficiaires du projet ont reçu les frais de scolarité ». Dans les deux cas, le mot anglais choisi existe bien, mais il a un \emph{autre} sens que le mot français auquel il ressemble. Voici les versions correctes :

\begin{exemplebox}[Corrections]
\eng{I attended the fund-raising gala last Saturday.} « J'ai assisté au gala de collecte de fonds samedi dernier. »\par
\eng{The beneficiaries of the project received school fees.} « Les bénéficiaires du projet ont reçu les frais de scolarité. »
\end{exemplebox}

Retenez l'idée centrale : en anglais comme en français, un mot a sa propre vie. La ressemblance de forme ne garantit jamais l'identité de sens.

\courssub{Les sept faux amis à connaître par cœur}

\begin{attention}[Erreur fréquente au Bac : to assist]
\eng{*I assisted to the fund-raising gala.} est \textbf{faux} pour deux raisons : \eng{to assist} signifie \emph{aider} (et non « assister à »), et il se construit \emph{sans} préposition \eng{to}. Pour dire « assister à » (être présent), utilisez \eng{to attend} ou \eng{to be present at}.\par
Dites donc : \eng{I attended the fund-raising gala.} « J'ai assisté au gala. » ou \eng{I helped to organize the gala.} « J'ai aidé à organiser le gala. »
\end{attention}

\begin{propriete}[to assist = aider ; to attend = assister à]
\begin{itemize}
\item \cle{to assist} (quelqu'un) = \emph{aider}, \emph{prêter assistance}. Prononciation : « é-ssis-te ». Construction : \eng{to assist somebody (in/with something)}, sans préposition \eng{to}. \eng{Volunteers assisted the victims of the flood.} « Des volontaires ont aidé les victimes de l'inondation. »
\item \cle{to attend} = \emph{assister à}, \emph{être présent à} (une réunion, une école, une cérémonie). Prononciation : « é-tende ». \eng{She attends a secondary school in Buea.} « Elle fréquente un lycée à Buea. »
\item Noms associés : \eng{assistance} = \emph{aide} (\eng{financial assistance} « aide financière ») ; \eng{an assistant} = un assistant, un aide.
\end{itemize}
\end{propriete}

\begin{propriete}[a donation = un don ; dotation = endowment / grant]
\cle{a donation} (prononciation : « do-né-icheune ») est un \emph{don}, un geste libre et volontaire : \eng{The church made a donation of blankets to the orphanage.} « L'église a fait don de couvertures à l'orphelinat. » Le français « dotation » (somme attribuée par une institution, un budget) se dit \eng{an endowment} (prononciation : « ène-dou-èmente ») ou \eng{a grant} : \eng{The school received a government grant.} « L'école a reçu une dotation (subvention) de l'État. »
\end{propriete}

\begin{propriete}[to support = soutenir ; supporter (tolérer) = to tolerate / to bear / to stand]
\cle{to support} a group, a family, a project = \emph{soutenir} (financièrement, moralement) : \eng{We support the local support group with monthly donations.} « Nous soutenons le groupe d'entraide local par des dons mensuels. » En revanche, « supporter » au sens de \emph{tolérer, endurer} se dit \eng{to tolerate}, \eng{to bear} ou \eng{to stand} : \eng{I can't stand this noise.} « Je ne supporte pas ce bruit. » Attention aussi : \eng{a supporter} = un partisan, un sympathisant (d'une cause, d'un club) — jamais « celui qui tolère ».
\end{propriete}

\begin{propriete}[a contribution = un apport volontaire ; cotisation = membership fee / dues / subscription]
\cle{a contribution} (prononciation : « kon-tri-biou-cheune ») est un \emph{apport volontaire} (d'argent, de temps, d'idées) : \eng{Everyone made a small contribution to the harvest.} « Chacun a apporté une petite contribution à la collecte. » La « cotisation » — somme \emph{obligatoire et régulière} versée par les membres d'une association — se dit \eng{a membership fee}, \eng{dues} (toujours au pluriel, prononciation : « diouze ») ou \eng{a subscription} : \eng{Members pay their dues every month.} « Les membres paient leurs cotisations chaque mois. »
\end{propriete}

\begin{aretenir}[Donation ou cotisation ?]
Une \cle{donation} est \textbf{libre et ponctuelle} : on donne ce qu'on veut, quand on veut. Les \cle{membership fees / dues} sont \textbf{obligatoires et réguliers} : on les paie pour rester membre. Comparez : \eng{She made a donation of 5,000 francs to the fund-raising.} « Elle a fait un don de 5 000 francs. » / \eng{The annual subscription to the alumni association is 2,000 francs.} « La cotisation annuelle à l'association des anciens élèves est de 2 000 francs. »
\end{aretenir}

\begin{propriete}[to collect : collecter ou collectionner ? Le contexte décide]
\cle{to collect} a deux grands emplois : \eng{to collect money / funds / donations} = \emph{collecter, récolter} (une quête) : \eng{Students collected money for the new borehole.} « Les élèves ont collecté de l'argent pour le nouveau forage. » ; \eng{to collect stamps / coins} = \emph{collectionner} (un hobby) : \eng{My uncle collects old coins.} « Mon oncle collectionne les pièces anciennes. » Le nom \eng{a collection} garde la même ambiguïté : \eng{a collection for the sick} « une quête pour les malades » ; \eng{a stamp collection} « une collection de timbres ».
\end{propriete}

\begin{propriete}[a benefactor = un bienfaiteur ; bénéficiaire = beneficiary]
Piège majeur : les deux mots existent en anglais avec des sens \emph{opposés}. \cle{a benefactor} (prononciation : « bé-ne-fècteure ») = celui qui \textbf{donne}, le bienfaiteur ; \cle{a beneficiary} (prononciation : « bé-ne-fi-chi-ére ») = celui qui \textbf{reçoit}, le bénéficiaire. Ne les inversez jamais.
\end{propriete}

\begin{exemplebox}[Benefactor ou beneficiary ?]
\eng{The benefactor of the project is a local businessman.} « Le bienfaiteur du projet est un homme d'affaires local. »\par
\eng{The beneficiaries received school fees.} « Les bénéficiaires ont reçu les frais de scolarité. »
\end{exemplebox}

\begin{propriete}[actually = en fait ; actuellement = currently / at present]
\cle{actually} (prononciation : « èk-choû-e-li ») signifie \emph{en fait, en réalité} (et sert souvent à corriger une idée fausse) : \eng{People think the gala was expensive, but it actually cost nothing.} « Les gens pensent que le gala a coûté cher, mais en fait il n'a rien coûté. » « Actuellement » (en ce moment) se dit \eng{currently} (prononciation : « keur-ren-te-li »), \eng{at present} ou \eng{at the moment} : \eng{We are currently collecting funds for the library.} « Nous collectons actuellement des fonds pour la bibliothèque. »
\end{propriete}

\courssub{Tableau récapitulatif des faux amis}

\begin{center}
\begin{tabularx}{\linewidth}{|X|X|X|}
\hline
\rowcolor{popblueL} Mot anglais & Sens réel en anglais & Piège français \\
\hline
to assist & aider (sans \eng{to}) & « assister à » = to attend \\
\hline
a donation & un don (libre, ponctuel) & « dotation » = endowment / grant \\
\hline
to support & soutenir & « supporter » (tolérer) = to tolerate / to bear / to stand \\
\hline
a contribution & un apport volontaire & « cotisation » = membership fee / dues \\
\hline
to collect money & collecter de l'argent & « collectionner » = to collect stamps (hobby) \\
\hline
a benefactor & un bienfaiteur (qui donne) & « bénéficiaire » = a beneficiary \\
\hline
actually & en fait, en réalité & « actuellement » = currently / at present \\
\hline
\end{tabularx}
\end{center}

\begin{popfigure}
\begin{tikzpicture}[
  col/.style={font=\bfseries\small, text=white, align=center},
  cell/.style={font=\small, align=left, text width=3.4cm},
  cellb/.style={font=\small, align=left, text width=3.6cm},
  cellc/.style={font=\small, align=left, text width=3.8cm},
  row/.style={minimum height=0.9cm}
]
\node[col, fill=popblue, minimum width=3.8cm] (h1) at (0,0) {English word};
\node[col, fill=popgreen, minimum width=4.0cm] (h2) at (4.0,0) {Real meaning};
\node[col, fill=poporange, minimum width=4.2cm] (h3) at (8.2,0) {French trap};
\node[cell, fill=popblueL, minimum width=3.8cm, minimum height=0.9cm] at (0,-0.9) {to assist};
\node[cellb, fill=popgreenL, minimum width=4.0cm, minimum height=0.9cm] at (4.0,-0.9) {to help};
\node[cellc, fill=poporangeL, minimum width=4.2cm, minimum height=0.9cm] at (8.2,-0.9) {« assister à » = to attend};
\node[cell, fill=popblueL, minimum width=3.8cm, minimum height=0.9cm] at (0,-1.85) {a benefactor};
\node[cellb, fill=popgreenL, minimum width=4.0cm, minimum height=0.9cm] at (4.0,-1.85) {the one who gives};
\node[cellc, fill=poporangeL, minimum width=4.2cm, minimum height=0.9cm] at (8.2,-1.85) {« bénéficiaire » = a beneficiary};
\node[cell, fill=popblueL, minimum width=3.8cm, minimum height=0.9cm] at (0,-2.8) {actually};
\node[cellb, fill=popgreenL, minimum width=4.0cm, minimum height=0.9cm] at (4.0,-2.8) {in fact};
\node[cellc, fill=poporangeL, minimum width=4.2cm, minimum height=0.9cm] at (8.2,-2.8) {« actuellement » = currently};
\node[cell, fill=popblueL, minimum width=3.8cm, minimum height=0.9cm] at (0,-3.75) {a contribution};
\node[cellb, fill=popgreenL, minimum width=4.0cm, minimum height=0.9cm] at (4.0,-3.75) {a voluntary gift};
\node[cellc, fill=poporangeL, minimum width=4.2cm, minimum height=0.9cm] at (8.2,-3.75) {« cotisation » = dues};
\end{tikzpicture}
\legende{Tableau à trois colonnes : mot anglais, sens réel, piège français.}
\end{popfigure}

\courssub{Le lexique en situation : un texte piège}

\begin{textebox}[A fund-raising day at Government High School, Bamenda]
Last term, our school organised a fund-raising day to buy books for the library. Many people attended the event, and several volunteers assisted the organising committee: they collected money at the entrance and sold tickets for the raffle. A local businessman, the main benefactor of the project, made a generous donation of 200,000 francs. “I actually studied in this school twenty years ago,” he said with a smile, “so I am happy to support it today.” The members of the Parents' Association also paid their annual dues during the event, and every class made a small contribution. At the end of the day, the headmaster announced the results: the school had raised enough money to buy three hundred books. The real beneficiaries, of course, are the students, who will now enjoy a better library. The headmaster thanked everyone and reminded us that solidarity is currently the pride of our community.
\end{textebox}

\textbf{Glossaire.} \emph{raffle} (n.) : tombola ; \emph{dues} (n. plur.) : cotisations ; \emph{to raise money} (v.) : récolter de l'argent ; \emph{headmaster} (n.) : proviseur ; \emph{pride} (n.) : fierté ; \emph{organising committee} (n.) : comité d'organisation ; \emph{borehole} (n.) : forage.

\textbf{Questions de compréhension.}
\begin{enumerate}
\item Who attended the fund-raising day, and what did the volunteers do?
\item Who is the benefactor of the project, and what did he give?
\item What did the members of the Parents' Association pay?
\item Who are the real beneficiaries of the project?
\item Find in the text two words that are false friends for a French speaker and explain their real meaning.
\end{enumerate}

\courssub{Méthode : comment éviter les calques du français}

\begin{methode}[Le réflexe anti-faux-amis en quatre étapes]
\begin{enumerate}
\item \textbf{Repérez} le mot que vous voulez écrire : ressemble-t-il à un mot français (\eng{assist}, \eng{donation}, \eng{support}, \eng{actually}...) ?
\item \textbf{Interrogez-vous} : « Est-ce le sens anglais ou le sens français que j'ai en tête ? » Traduisez mentalement votre phrase en français clair.
\item \textbf{Vérifiez} dans votre liste de faux amis (ou un dictionnaire) : le mot anglais porte-t-il bien ce sens ?
\item \textbf{Remplacez} si nécessaire par l'équivalent correct : \eng{to attend}, \eng{currently}, \eng{dues}, \eng{beneficiary}...
\end{enumerate}
Exemple : vous voulez écrire « J'assiste actuellement aux réunions du groupe d'entraide. » Étape 1 : \eng{assist} et \eng{actually} sont suspects. Étape 2 : sens voulu = « être présent » et « en ce moment ». Étape 3 : \eng{to assist} = aider ; \eng{actually} = en fait. Étape 4 : \eng{I currently attend the support group meetings.}
\end{methode}

\courssub{Entraînement}

\begin{exoresolu}[Corriger les faux amis]
Énoncé — Corrigez les phrases suivantes, toutes fautives à cause d'un faux ami.\par
1. \eng{*I assisted to the meeting of the support group.} (« J'ai assisté à la réunion. »)\par
2. \eng{*The benefactors received food and clothes.} (« Les bénéficiaires ont reçu... »)\par
3. \eng{*We are actually collecting funds for the borehole.} (« Nous collectons actuellement... »)\par
4. \eng{*Members must pay a donation every month.} (« Les membres doivent payer une cotisation. »)\par
5. \eng{*I cannot support this heat.} (« Je ne supporte pas cette chaleur. »)
\tcblower
Solution détaillée.\par
1. \eng{I attended the meeting of the support group.} — « assister à » = \eng{to attend} ; \eng{to assist} = aider et ne prend pas \eng{to}.\par
2. \eng{The beneficiaries received food and clothes.} — celui qui reçoit est le \eng{beneficiary} ; le \eng{benefactor} est celui qui donne.\par
3. \eng{We are currently collecting funds for the borehole.} — « actuellement » = \eng{currently} ; \eng{actually} = en fait.\par
4. \eng{Members must pay their dues every month.} (ou \eng{a membership fee}) — une cotisation obligatoire n'est pas une \eng{donation}, qui est libre et ponctuelle.\par
5. \eng{I cannot stand this heat.} (ou \eng{I cannot bear this heat.}) — « supporter » (tolérer) = \eng{to stand / to bear / to tolerate} ; \eng{to support} = soutenir.
\end{exoresolu}

\begin{exercice}[À vous de jouer]
Traduisez en anglais, en faisant attention aux faux amis :
\begin{enumerate}
\item Les bienfaiteurs du projet ont assisté à la cérémonie de remise des dons.
\item Les membres paient leurs cotisations chaque mois, mais les dons sont libres.
\item En fait, les bénéficiaires ont déjà reçu les fournitures scolaires.
\item Des volontaires aident le comité à collecter de l'argent pour le forage.
\item Nous soutenons actuellement trois groupes d'entraide à Douala.
\end{enumerate}
\end{exercice}

\begin{corrige}[Exercice : traduction et faux amis]
1. \eng{The benefactors of the project attended the donation ceremony.} — \eng{benefactors} = ceux qui donnent ; « assister à » = \eng{to attend}.\par
2. \eng{Members pay their dues every month, but donations are voluntary.} — « cotisations » = \eng{dues} ; « dons » = \eng{donations}.\par
3. \eng{Actually, the beneficiaries have already received the school supplies.} — « en fait » = \eng{actually} ; « bénéficiaires » = \eng{beneficiaries}.\par
4. \eng{Volunteers assist the committee in collecting money for the borehole.} — « aider » = \eng{to assist} (sans \eng{to}) ; « collecter » = \eng{to collect}.\par
5. \eng{We currently support three support groups in Douala.} — « actuellement » = \eng{currently} ; « soutenir » = \eng{to support}.
\end{corrige}

\begin{aretenir}[L'essentiel de la section]
Sept faux amis à ne plus jamais confondre : \eng{to assist} (aider) / \eng{to attend} (assister à) ; \eng{a donation} (un don) / \eng{a grant} (une dotation) ; \eng{to support} (soutenir) / \eng{to tolerate} (supporter) ; \eng{a contribution} (apport libre) / \eng{dues} (cotisations) ; \eng{to collect money} (collecter) / \eng{to collect stamps} (collectionner) ; \eng{a benefactor} (bienfaiteur) / \eng{a beneficiary} (bénéficiaire) ; \eng{actually} (en fait) / \eng{currently} (actuellement). Avant d'écrire un mot qui ressemble au français, posez-vous toujours la question : « sens anglais ou sens français ? »
\end{aretenir}

\coursec{Emploi du lexique en contexte et production guidée}

Dans la section précédente, nous avons appris à débusquer les \cle{faux amis} et les calques du français qui guettent le candidat au Probatoire et au Baccalauréat : on ne dit pas \eng{to make a collect} mais \eng{to hold a collection}, on ne dit pas \eng{to assist an event} mais \eng{to attend an event}. Savoir \emph{éviter} les erreurs ne suffit cependant pas : il faut maintenant savoir \emph{produire}. C'est l'objet de cette dernière section : réemployer tout le lexique de la leçon dans des situations de communication réelles — un dialogue, un article de journal, une annonce, une lettre d'appel à contributions et une interview — comme on vous le demandera à l'examen et dans la vie scolaire.

\courssub{Dialogue modèle : organiser une collecte pour la bibliothèque}

Observez d'abord ce dialogue entre deux élèves de Tle A du lycée de Nkongsamba. Repérez les formules de demande d'aide, les questions sur les sommes récoltées et les expressions d'encouragement.

\begin{textebox}[Model dialogue — Raising funds for the school library]
\textbf{A:} Hello, Mirabel! Our support group is raising funds to buy new books and shelves for the school library. Could you help us?

\textbf{B:} That's a worthy cause! How much have you raised so far?

\textbf{A:} We have collected about 150,000 francs, but our target is 500,000 francs. We are still far from it.

\textbf{B:} How can I help?

\textbf{A:} You could make a donation, however small, or volunteer at our charity concert next Saturday. We are also appealing to well-wishers and former students of the school.

\textbf{B:} I'll do both. I can donate 5,000 francs and sell tickets at the entrance. By the way, have you asked the Parent-Teacher Association for support?

\textbf{A:} Yes, they have pledged 100,000 francs, and a local businessman has promised to donate in kind: twenty boxes of chalk and some furniture.

\textbf{B:} Wonderful! Every little helps. I'm sure you will reach your target before the end of the term.

\textbf{A:} Thank you so much for your generosity. We really appreciate it.
\end{textebox}

\textbf{Glossaire.} \emph{worthy cause} (n.) : cause noble ; \emph{so far} (adv.) : jusqu'à présent ; \emph{target} (n.) : objectif chiffré ; \emph{well-wisher} (n.) : bienfaiteur, personne de bonne volonté ; \emph{to pledge} (v.) : promettre (un don) ; \emph{in kind} (adv.) : en nature ; \emph{every little helps} (expr.) : chaque petit geste compte.

\begin{aretenir}[Formules-clés du dialogue]
\begin{itemize}
\item Demander de l'aide : \eng{Could you help us raise funds for...?} « Pourriez-vous nous aider à récolter des fonds pour... ? »
\item Demander le montant récolté : \eng{How much have we raised so far?} « Combien avons-nous récolté jusqu'à présent ? » (present perfect + \eng{so far})
\item Lancer un appel : \eng{We are appealing to well-wishers...} « Nous lançons un appel aux bienfaiteurs... »
\item Proposer son aide : \eng{How can I help?} / \eng{You could make a donation or volunteer.} « Vous pourriez faire un don ou vous porter volontaire. »
\item Encourager : \eng{Every little helps!} « Chaque petit geste compte ! »
\end{itemize}
\end{aretenir}

\begin{exemplebox}[Exemples traduits]
\eng{We are appealing to all well-wishers to help us reach our target of one million francs.} « Nous lançons un appel à tous les bienfaiteurs pour nous aider à atteindre notre objectif d'un million de francs. »

\eng{The support group has raised 300,000 francs so far, thanks to generous donors.} « Le groupe de soutien a récolté 300 000 francs jusqu'à présent, grâce à des donateurs généreux. »

\eng{Could you contribute in cash or in kind?} « Pourriez-vous contribuer en espèces ou en nature ? »
\end{exemplebox}

\courssub{Texte à trous : un article de journal sur une action caritative}

Voici un article original, tel qu'on pourrait en lire dans la presse anglophone du Cameroun. Complétez-le avec le lexique de la leçon : c'est exactement le type d'exercice proposé en compréhension de l'écrit.

\begin{exoresolu}[Gap-fill — A charity drive in Bamenda]
\textbf{Énoncé.} Complete the article with these words: \emph{donation, fund-raising, volunteers, target, well-wishers, in kind, appeal, support group, proceeds, worthy}.

Last Saturday, the Women's \textbf{(1)} \ldots\ of Bamenda organised a \textbf{(2)} \ldots\ campaign to equip the district hospital. The group launched an \textbf{(3)} \ldots\ to all \textbf{(4)} \ldots\ in the North-West Region. Their \textbf{(5)} \ldots\ was two million francs. About fifty \textbf{(6)} \ldots\ sold raffle tickets and collected money at the commercial avenue. A local company made a generous \textbf{(7)} \ldots\ of 500,000 francs, while market women contributed \textbf{(8)} \ldots\ , offering food and drinks for the event. The \textbf{(9)} \ldots\ from the charity concert will be used to buy beds and mattresses. “It is a \textbf{(10)} \ldots\ cause,” said the chairlady, “and every little helps.”

\tcblower
\textbf{Solution détaillée.} (1) \emph{support group} : collocation fixe, « groupe de soutien ». (2) \emph{fund-raising} : adjectif composé devant \emph{campaign}. (3) \emph{appeal} : \eng{to launch an appeal to somebody}, « lancer un appel à ». (4) \emph{well-wishers} : les bienfaiteurs à qui l'on s'adresse. (5) \emph{target} : l'objectif chiffré (deux millions). (6) \emph{volunteers} : ce sont eux qui vendent les billets ; attention au pluriel avec \emph{fifty}. (7) \emph{donation} : \eng{to make a donation of...}, collocation obligatoire (jamais \eng{do a donation}). (8) \emph{in kind} : en nature (nourriture, boissons), par opposition à \eng{in cash}. (9) \emph{proceeds} : les recettes, toujours au pluriel dans ce sens. (10) \emph{worthy} : \eng{a worthy cause}, « une noble cause ».
\end{exoresolu}

\courssub{Transformation de phrases : du calque à la collocation correcte}

Exercice redoutable mais formateur : on vous donne une phrase contenant un calque du français ; vous la réécrivez avec la collocation anglaise correcte.

\begin{center}
\begin{tabularx}{\linewidth}{|X|X|}
\hline
\rowcolor{popblueL} \textbf{Calque du français (à corriger)} & \textbf{Collocation correcte} \\
\hline
\eng{Please, we beg you to give us money.} & \eng{We kindly appeal to you for support.} \\
\hline
\eng{They did a donation of 10,000 francs.} & \eng{They made a donation of 10,000 francs.} \\
\hline
\eng{We want to gain money for the health centre.} & \eng{We want to raise money for the health centre.} \\
\hline
\eng{The association organises a collect of clothes.} & \eng{The association is holding a collection of clothes.} \\
\hline
\eng{Many persons assisted the charity concert.} & \eng{Many people attended the charity concert.} \\
\hline
\eng{We thank you for your generosity in advance.} & \eng{Your generous contribution will be highly appreciated.} \\
\hline
\end{tabularx}
\end{center}

\begin{attention}[Dans une annonce formelle, bannissez les calques]
Ne traduisez jamais « S'il vous plaît, nous vous supplions » par \eng{Please, we beg you} : c'est un calque maladroit, jugé trop familier, voire incorrect dans un document officiel. Préférez les formules consacrées : \eng{We kindly appeal to...} « Nous lançons un appel à... », \eng{Your generous contribution will be highly appreciated.} « Votre généreuse contribution sera très appréciée. », \eng{Any donation, however small, will be welcome.} « Tout don, même modeste, sera le bienvenu. » De même, \eng{to gain money} signifie « gagner de l'argent » (par son travail) ; pour une collecte, on dit \eng{to raise money/funds}. Enfin, dans la langue courante, on préfère \eng{people} à \eng{persons}, réservé aux textes juridiques ou administratifs.
\end{attention}

\begin{exercice}[Transform the sentences]
Rewrite each sentence, replacing the French calque with the correct collocation.
\begin{enumerate}
\item \eng{Our club wants to gain funds for new jerseys.}
\item \eng{The chief did a donation of two bags of cement.}
\item \eng{Please, we beg you, come to our fund-raising gala.}
\item \eng{Over 200 persons assisted the launching ceremony.}
\end{enumerate}
\end{exercice}

\begin{corrige}[Exercice « Transform the sentences »]
\begin{enumerate}
\item \eng{Our club wants to raise funds for new jerseys.} (\eng{raise funds}, pas \eng{gain funds}.)
\item \eng{The chief made a donation of two bags of cement.} (\eng{make a donation}.)
\item \eng{We kindly appeal to you to attend our fund-raising gala.} ou \eng{Your presence at our fund-raising gala will be highly appreciated.}
\item \eng{Over 200 people attended the launching ceremony.} (\eng{attend} = assister à ; \eng{people}, pas \eng{persons} dans ce contexte.)
\end{enumerate}
\end{corrige}

\courssub{Rédaction guidée : l'appel à contributions (appeal letter)}

L'\cle{appeal letter} (lettre d'appel à contributions) est un document semi-formel très fréquent dans la vie associative camerounaise et parfois demandé en \emph{guided writing}. Sa structure est rigide : respectez-la.

\begin{methode}[Rédiger un appel à contributions en cinq étapes]
\begin{enumerate}
\item \textbf{Présenter le groupe et le projet} : \eng{The Science Club of Government High School Buea is raising funds to equip the school laboratory.}
\item \textbf{Expliquer le besoin} (resources/services) : \eng{Our laboratory lacks microscopes and chemicals, and students cannot do practical work.}
\item \textbf{Fixer un objectif chiffré} : \eng{Our target is one million francs. So far, we have raised 400,000 francs.}
\item \textbf{Dire comment contribuer} (cash or in kind) : \eng{You may contribute in cash at the bursar's office or in kind (equipment, books).}
\item \textbf{Remercier les donateurs à l'avance} : \eng{Your generous contribution will be highly appreciated. Thank you for supporting education.}
\end{enumerate}
\end{methode}

\begin{popfigure}
\begin{center}
\begin{tikzpicture}[
node distance=0.45cm,
box/.style={draw=popblue, fill=popblueL, rounded corners=2pt, text width=8.6cm, align=center, font=\small, inner sep=5pt},
fleche/.style={-{Stealth[length=2.5mm]}, thick, poporange}
]
\node[box, fill=popgoldL, draw=popgold] (titre) {\textbf{Title / Heading} — \eng{Appeal for Support}};
\node[box, below=of titre] (pb) {\textbf{1. Problem} — present the group and the need};
\node[box, below=of pb] (obj) {\textbf{2. Target} — state the amount aimed at};
\node[box, below=of obj] (besoin) {\textbf{3. What we need} — resources / services (cash or in kind)};
\node[box, below=of besoin] (don) {\textbf{4. How to give} — where, when, to whom};
\node[box, below=of don, fill=popgreenL, draw=popgreen] (merci) {\textbf{5. Thanks} — thank donors in advance};
\draw[fleche] (titre) -- (pb);
\draw[fleche] (pb) -- (obj);
\draw[fleche] (obj) -- (besoin);
\draw[fleche] (besoin) -- (don);
\draw[fleche] (don) -- (merci);
\end{tikzpicture}
\end{center}
\legende{Structure d'un appel à contributions : du titre aux remerciements.}
\end{popfigure}

\begin{exemplebox}[Modèle d'annonce (80-100 mots)]
\textbf{APPEAL FOR SUPPORT}

\eng{The Readers' Club of GBHS Bamenda is raising funds to equip the school library with new books and shelves. Many of our books are old and torn, and students have nothing to read. Our target is 500,000 francs; so far, we have raised 150,000 francs. We are appealing to all well-wishers, parents and former students to support this worthy cause. You may contribute in cash at the bursar's office or in kind (books, shelves, chairs). A charity concert will be held on 15th March in the school hall. Your generous contribution will be highly appreciated. Every little helps!}
\end{exemplebox}

\begin{exercice}[Guided writing]
Your support group wants to buy a computer and a printer for your school. Write an appeal letter of 80-100 words to the community, following the five-step structure. Use at least eight target words from the lesson (e.g. \emph{raise funds, target, appeal, donation, in kind, well-wishers, worthy cause, appreciate}).
\end{exercice}

\begin{corrige}[Exercice « Guided writing » — éléments de correction]
Barème indicatif : respect de la structure en cinq étapes (titre, problème, objectif chiffré, modalités du don, remerciements) ; au moins huit mots-cibles correctement employés ; aucune formule calquée (\eng{we beg you} est éliminatoire) ; registre semi-formel ; longueur respectée. Exemple de début attendu : \eng{The Young Farmers' Club is raising funds to buy a computer and a printer for the school office. Our target is 350,000 francs...}
\end{corrige}

\courssub{Jeu de rôle : interview d'un responsable d'ONG}

\begin{experience}[Role play — Interviewing an NGO leader]
Travail en binômes. L'élève A est journaliste du journal scolaire ; l'élève B est le président d'une ONG qui vient de construire un forage (\emph{borehole}) dans un village.
\begin{enumerate}
\item Le journaliste prépare cinq questions avec le lexique de la leçon : \eng{How did you raise the funds? Who are your main donors? What was your target? Did people contribute in kind? Do you have any volunteers?}
\item Le responsable répond en réemployant au moins six mots-cibles : \emph{fund-raising, donation, pledge, proceeds, support group, in cash, in kind, grateful}.
\item La classe écoute et coche dans une grille chaque mot-cible entendu.
\item Inversion des rôles, puis compte rendu oral de deux minutes devant la classe.
\end{enumerate}
\end{experience}

\courssub{Auto-évaluation : grille de reprise du lexique}

Avant de rendre votre production (annonce, lettre ou compte rendu), vérifiez que vous avez réemployé au moins \textbf{huit mots-cibles}. Cochez la colonne de droite.

\begin{center}
\begin{tabularx}{\linewidth}{|X|l|l|}
\hline
\rowcolor{popblueL} \textbf{Mot-cible} & \textbf{Traduction} & \textbf{Utilisé ?} \\
\hline
\eng{to raise funds / money} & récolter des fonds & $\square$ \\
\hline
\eng{target} & objectif chiffré & $\square$ \\
\hline
\eng{to appeal to well-wishers} & lancer un appel aux bienfaiteurs & $\square$ \\
\hline
\eng{to make a donation} & faire un don & $\square$ \\
\hline
\eng{in cash / in kind} & en espèces / en nature & $\square$ \\
\hline
\eng{support group} & groupe de soutien & $\square$ \\
\hline
\eng{volunteer} & bénévole / se porter volontaire & $\square$ \\
\hline
\eng{proceeds} & recettes & $\square$ \\
\hline
\eng{worthy cause} & noble cause & $\square$ \\
\hline
\eng{to pledge} & promettre (un don) & $\square$ \\
\hline
\eng{highly appreciated} & très apprécié(e) & $\square$ \\
\hline
\eng{every little helps} & chaque petit geste compte & $\square$ \\
\hline
\end{tabularx}
\end{center}

\begin{aretenir}[Ce qu'il faut retenir de la leçon]
Pour communiquer efficacement sur les activités de collecte de fonds : utilisez les collocations exactes (\eng{raise funds}, \eng{make a donation}, \eng{hold a collection}, \eng{reach a target}) ; structurez tout appel à contributions en cinq étapes (groupe et projet, besoin, objectif chiffré, modalités du don, remerciements) ; employez les formules formelles de politesse (\eng{We kindly appeal to...}, \eng{Your generous contribution will be highly appreciated.}) et jamais les calques du français ; enfin, réemployez au moins huit mots-cibles dans chaque production, à l'oral comme à l'écrit. Ce lexique est un classique des épreuves de \emph{reading} et de \emph{guided writing} au Probatoire et au Baccalauréat.
\end{aretenir}

\newpage

\coursec{Activité d'intégration}

\courssub{Objectifs de l'activité}

À la fin de cette activité d'intégration, l'élève sera capable de :

\begin{itemize}
\item réemployer à l'oral et à l'écrit le lexique des collectes de fonds (\cle{raise funds}, \cle{donation}, \cle{appeal}, \cle{proceeds}, \cle{target amount}, \cle{well-wishers}, \cle{in kind}, \cle{volunteers}) dans une situation de communication authentique ;
\item produire un court texte persuasif (affiche d'appel à contributions) de 80 à 100 mots, correctement organisé ;
\item conduire un échange oral de sollicitation d'un donateur avec les formules de politesse appropriées ;
\item éviter les faux amis étudiés dans la leçon (\eng{to assist} $\neq$ « assister à », \eng{actually} $\neq$ « actuellement », \eng{a benefactor} $\neq$ « un bénéficiaire », \eng{to achieve} $\neq$ « acheter ») ;
\item travailler en groupe, répartir les rôles et évaluer la production des autres selon une grille de critères explicites.
\end{itemize}

\courssub{Situation}

Vous êtes élèves de Terminale A au lycée bilingue de votre localité. Le village de Mbankomo, à une vingtaine de kilomètres de Yaoundé, ne dispose pas de point d'eau potable : les habitants, surtout les femmes et les enfants, parcourent chaque jour plusieurs kilomètres jusqu'à une source non aménagée. Le comité de développement du village a lancé l'opération « Operation Borehole » pour financer le forage d'un puits équipé d'une pompe. Le coût total est estimé à 2 500 000 FCFA. Le chef du village a écrit aux écoles de la région pour demander leur aide. Votre classe a décidé de répondre à cet appel en créant des \eng{support groups} chargés chacun de concevoir une campagne de collecte de fonds.

Voici la lettre du comité de développement, qui sert de document déclencheur.

\begin{textebox}[Letter from the Mbankomo Development Committee]
Dear students,

Our village of Mbankomo has no safe drinking water. Every morning, women and children walk about three kilometres to reach an unprotected stream. During the dry season, the stream almost dries up, and waterborne diseases are common among our children.

We have decided to launch “Operation Borehole”. Our target amount is 2,500,000 FCFA, which will cover the drilling of a borehole, a hand pump and a concrete platform. A local company has already promised a donation of 500,000 FCFA, and several well-wishers abroad have pledged support. We are also grateful for contributions in kind: cement, pipes, or even free labour from skilled volunteers.

We are appealing to schools, churches, associations and business people to help us raise the remaining funds. Every donation, however small, will bring us closer to our goal. All proceeds will be managed transparently by a committee of five members, and a financial report will be read publicly at the end of the campaign.

We count on your generosity and your dynamism.

Yours faithfully,

The Chairman, Mbankomo Development Committee
\end{textebox}

\textbf{Glossaire :} \emph{drinking water} (eau potable) ; \emph{stream} (ruisseau) ; \emph{to dry up} (s'assécher) ; \emph{waterborne diseases} (maladies transmises par l'eau) ; \emph{drilling} (forage) ; \emph{to pledge} (promettre un don) ; \emph{however small} (si petit soit-il) ; \emph{transparently} (en toute transparence).

\begin{attention}[Lecture attentive : deux pièges dans la lettre]
\begin{itemize}
\item \eng{to appeal to somebody} signifie « lancer un appel à quelqu'un, solliciter » : \eng{We are appealing to schools\ldots} « Nous lançons un appel aux écoles\ldots ». Ne pas confondre avec \eng{to appeal to} au sens de « plaire à » (\eng{This music appeals to young people.}).
\item \eng{proceeds} (toujours au pluriel) signifie « les recettes, le produit d'une collecte », et non « les procédures ». Prononciation : « pro-ssidz ».
\end{itemize}
\end{attention}

\courssub{Consignes}

Travail en groupes de quatre. L'activité se déroule en cinq tâches numérotées, de la compréhension à la production.

\begin{enumerate}
\item \textbf{Compréhension du document (10 min).} Lisez la lettre en silence, puis répondez oralement dans votre groupe : What problem does the village face? What is the target amount? Who has already contributed, and how? What does “contributions in kind” mean? Who will manage the proceeds? Soulignez dans la lettre au moins huit mots ou expressions du champ lexical de la collecte de fonds.

\item \textbf{Création du support group (5 min).} Choisissez un nom anglais pour votre groupe (par exemple \eng{Hope Builders}, \eng{Water for All}) et fixez votre propre objectif : votre groupe s'engage à récolter une partie de la somme (par exemple 250 000 FCFA) ou à financer un équipement précis (la pompe, la plateforme en béton). Désignez un \eng{chairperson}, un \eng{treasurer}, un \eng{secretary} et un \eng{spokesperson}.

\item \textbf{Production écrite : l'affiche d'appel à contributions (25 min).} Rédigez une affiche/annonce de 80 à 100 mots destinée aux élèves, aux parents et aux commerçants du quartier. Votre texte doit contenir au moins huit mots-cibles de la leçon : \eng{raise funds}, \eng{donation}, \eng{appeal}, \eng{volunteers}, \eng{proceeds}, \eng{in kind}, \eng{well-wishers}, \eng{target}, \eng{support group}, \eng{generosity}. Il doit comporter un titre accrocheur, le problème, l'objectif chiffré, les modalités de don et une phrase de remerciement. Attention aux faux amis : n'écrivez jamais \eng{assist the ceremony} pour « assister à la cérémonie » (\eng{attend}), ni \eng{actually} pour « actuellement » (\eng{currently, at present}).

\item \textbf{Production orale : le jeu de rôle (15 min par groupe).} Préparez un mini-dialogue de 8 à 10 répliques : le \eng{spokesperson} et le \eng{treasurer} de votre groupe rendent visite à un donateur potentiel (un commerçant du marché, joué par un autre membre) pour solliciter un don. Utilisez les formules de sollicitation polie : \eng{Would you consider making a donation?}, \eng{Any contribution, however small, would be greatly appreciated.}, \eng{We would be grateful for your support.} Le donateur pose au moins deux questions (sur l'emploi des fonds, sur le montant) avant de se décider. Jouez la scène devant la classe.

\item \textbf{Vote et évaluation (10 min).} Après chaque prestation, la classe note la campagne sur une grille de 20 points : lexique (5 pts), correction grammaticale (5 pts), organisation de l'affiche (5 pts), créativité et conviction (5 pts). La classe vote ensuite pour la campagne la plus convaincante. Les affiches sont ramassées et notées par l'enseignant avec la même grille.
\end{enumerate}

\courssub{Ressources à mobiliser}

\begin{aretenir}[Boîte à outils de la campagne]
\begin{itemize}
\item \textbf{Lexique :} \eng{to raise funds / to launch an appeal / to make a donation / to donate / to pledge / proceeds / target amount / well-wishers / volunteers / contributions in kind / a support group / a fundraising campaign / generosity / to reach one's goal}.
\item \textbf{Collocations :} \eng{to launch a campaign}, \eng{to reach a target}, \eng{to make an appeal for funds}, \eng{to collect donations}, \eng{to contribute generously}, \eng{safe drinking water}.
\item \textbf{Structures de sollicitation :} \eng{Would you mind contributing\ldots?}, \eng{We are appealing to your generosity\ldots}, \eng{Your donation will help us to\ldots}, \eng{We would be grateful if you could\ldots}
\item \textbf{Grammaire utile :} le futur avec \eng{will} pour les promesses (\eng{Your gift will change lives}) ; les modaux \eng{could / would} pour la politesse ; l'impératif pour les slogans (\eng{Give water, give life!}) ; le présent simple pour les faits (\eng{Mbankomo has no safe water}).
\item \textbf{Faux amis à éviter :} \eng{to assist} = aider (\eng{attend} = assister à) ; \eng{actually} = en fait (\eng{currently} = actuellement) ; \eng{a benefactor} = un bienfaiteur (\eng{a beneficiary} = un bénéficiaire) ; \eng{to achieve} = accomplir, atteindre (\eng{to buy} = acheter) ; \eng{deception} = tromperie, duperie (\eng{disappointment} = déception).
\end{itemize}
\end{aretenir}

\begin{attention}[Le piège de « déception »]
Le français « déception » ne se traduit jamais par \eng{deception}. \eng{Deception} signifie « tromperie, duperie » : \eng{He obtained the money by deception.} « Il a obtenu l'argent par tromperie. » Pour « déception », on dit \eng{disappointment} : \eng{To our great disappointment, the target was not reached.} « À notre grande déception, l'objectif n'a pas été atteint. »
\end{attention}

\begin{methode}[Réussir l'affiche d'appel à contributions]
\begin{enumerate}
\item \textbf{Accrocher} avec un titre impératif ou une question : \eng{Give water, give life!} / \eng{Will you help Mbankomo?}
\item \textbf{Énoncer le problème} en une ou deux phrases au présent simple : \eng{Mbankomo has no safe drinking water.}
\item \textbf{Chiffrer l'objectif} : \eng{Our target is 250,000 FCFA.}
\item \textbf{Expliquer comment donner} (argent, dons en nature, bénévolat) : \eng{Contributions in kind are also welcome.}
\item \textbf{Rassurer} sur la gestion des fonds : \eng{All proceeds will go directly to the village committee.}
\item \textbf{Remercier} et conclure par un slogan : \eng{We thank you for your generosity.}
\end{enumerate}
\end{methode}

\courssub{Proposition de production et corrigé commenté}

\begin{exoresolu}[Modèle d'affiche et de dialogue — Operation Borehole]
Énoncé : produire l'affiche d'appel à contributions (80-100 mots) et le début du dialogue de sollicitation pour le groupe fictif « Water for All ».
\tcblower
\textbf{Affiche modèle :}

\eng{WATER FOR ALL — HELP MBANKOMO DRINK SAFE WATER!}

\eng{The village of Mbankomo has no drinking water. Our support group, “Water for All”, is launching an appeal to raise funds for a borehole. Our target is 250,000 FCFA. We invite students, parents, traders and all well-wishers to make a donation, however small. Contributions in kind — cement, pipes or free labour — are also welcome. Our volunteers will collect your gifts every Friday at the school gate. All proceeds will go directly to the village committee. Give water, give life! We thank you for your generosity.}

(98 mots — dix mots-cibles employés : \eng{support group}, \eng{appeal}, \eng{raise funds}, \eng{target}, \eng{well-wishers}, \eng{donation}, \eng{in kind}, \eng{volunteers}, \eng{proceeds}, \eng{generosity}.)

\textbf{Dialogue modèle (extrait) :}

\eng{— Good morning, Madam. We are members of “Water for All”, a support group from the bilingual high school.\\
— Good morning. What can I do for you?\\
— We are raising funds to drill a borehole in Mbankomo. Would you consider making a donation?\\
— How will the money be used?\\
— All proceeds will go to the village committee, and a public financial report will be made. Any contribution, however small, would be greatly appreciated.\\
— In that case, I will give 10,000 FCFA and some bags of cement in kind.\\
— Thank you very much for your generosity, Madam. You are a true well-wisher!}

\textbf{Commentaire des choix de langue :} le titre de l'affiche utilise l'impératif (\eng{Help\ldots Give water, give life!}), typique des slogans anglais. Le futur \eng{will collect}, \eng{will go} exprime l'engagement ferme. La sollicitation orale privilégie \eng{Would you consider\ldots?} et \eng{would be greatly appreciated}, plus polis qu'un impératif direct. Le mot \eng{well-wisher} est réemployé avec un sourire en conclusion. Aucun faux ami : on dit \eng{safe drinking water} (et non « potable water »), \eng{however small} pour « si petit soit-il ».
\end{exoresolu}

\begin{exercice}[Pour s'entraîner — réemploi du lexique]
Complétez les phrases suivantes avec le mot ou l'expression qui convient : \eng{proceeds}, \eng{pledged}, \eng{in kind}, \eng{target}, \eng{appeal}, \eng{benefactor}, \eng{volunteers}, \eng{raise}.

\begin{enumerate}
\item The students launched an \ldots\ to the whole neighbourhood.
\item An anonymous \ldots\ offered to pay for the hand pump.
\item Our \ldots\ is 250,000 FCFA; we have already collected half of it.
\item Many traders contributed \ldots : rice, soap and exercise books.
\item The \ldots\ of the concert will be handed over to the village committee.
\item Ten \ldots\ will sell tickets at the school gate.
\item The church has \ldots\ 100,000 FCFA to the project.
\item We hope to \ldots\ enough money before the end of the term.
\end{enumerate}
\end{exercice}

\begin{corrige}[Exercice « Pour s'entraîner »]
\begin{enumerate}
\item \eng{appeal} — « Les élèves ont lancé un appel à tout le quartier. »
\item \eng{benefactor} — « Un bienfaiteur anonyme a proposé de payer la pompe manuelle. » (et non \eng{beneficiary} = bénéficiaire)
\item \eng{target} — « Notre objectif est de 250 000 FCFA ; nous avons déjà récolté la moitié. »
\item \eng{in kind} — « De nombreux commerçants ont contribué en nature : riz, savon et cahiers. »
\item \eng{proceeds} — « Les recettes du concert seront remises au comité du village. »
\item \eng{volunteers} — « Dix bénévoles vendront des tickets à la grille de l'école. »
\item \eng{pledged} — « L'église a promis 100 000 FCFA au projet. » (verbe régulier : \eng{pledge, pledged, pledged})
\item \eng{raise} — « Nous espérons récolter assez d'argent avant la fin du trimestre. »
\end{enumerate}
\end{corrige}

\courssub{Bilan de l'activité}

Cette tâche d'intégration a permis de mobiliser l'ensemble des savoir-faire de la leçon dans une chaîne complète de communication : comprendre un document authentique (la lettre du comité), s'approprier le lexique des collectes de fonds et des groupes de soutien, le réemployer dans une production écrite persuasive et dans un échange oral de sollicitation. Le contexte — le financement d'un forage dans un village camerounais — donne un sens concret et civique au vocabulaire : les élèves ne mémorisent pas une liste, ils s'en servent pour convaincre. L'évaluation par grille (lexique, grammaire, organisation, créativité) et le vote final développent enfin l'esprit critique et la coopération. Pour aller plus loin, la classe peut réellement organiser une mini-collecte au profit d'un projet scolaire, en rédigeant un compte rendu en anglais : \eng{Our class raised\ldots The proceeds were handed over to\ldots}

\newpage

\coursec{Méthodes et exercices résolus}

\courssub{Méthode 1 : Identifier et mémoriser le vocabulaire du thème (fund-raising et support groups)}

\begin{methode}[Comment apprendre et reconnaître le lexique de la collecte de fonds]
Pour maîtriser le vocabulaire d'un champ lexical au Bac, il ne suffit pas de traduire mot à mot : il faut organiser, prononcer et réemployer.
\begin{enumerate}
\item \textbf{Repérer le thème du texte.} Dès les premières lignes, cherchez les mots-clés : \eng{fund-raising, charity, donation, community}. Ils annoncent le champ lexical et vous aident à deviner le sens des mots inconnus.
\item \textbf{Classer le vocabulaire en familles.} Faites trois colonnes dans votre cahier : \emph{people} (\eng{donor, volunteer, beneficiary, fundraiser, sponsor}), \emph{actions} (\eng{to raise funds, to donate, to sponsor, to organize a gala, to contribute}), \emph{things/events} (\eng{a charity gala, a raffle, a donation, a borehole, a scholarship, a charity walk}).
\item \textbf{Apprendre les mots avec leur collocution.} Ne retenez jamais un mot seul : retenez \eng{to raise funds} (et non \eng{to lift funds}), \eng{to make a donation}, \eng{to support a cause}, \eng{to organize an event}.
\item \textbf{Deviner le sens par le contexte.} Si vous lisez \eng{The villagers contributed money, cement and labour to build the health centre}, vous comprenez que \eng{contributed} signifie « ont contribué, donné » grâce aux mots \eng{money, cement, labour}.
\item \textbf{Vérifier la prononciation et l'orthographe.} Notez la prononciation en lettres françaises : \eng{fund-raising} « fond-ré-zing », \eng{charity} « tcha-ri-ti », \eng{borehole} « bo-rol », \eng{volunteer} « vo-lon-tir », \eng{donation} « do-né-chonne », \eng{sponsor} « spon-seur ».
\item \textbf{Réemployer immédiatement.} Après chaque liste, écrivez trois phrases personnelles sur votre école ou votre quartier : \eng{Our class organized a fund-raising event to buy books for the school library.} « Notre classe a organisé une collecte de fonds pour acheter des livres à la bibliothèque scolaire. »
\end{enumerate}
\end{methode}

\begin{exoresolu}[Compréhension de texte et repérage lexical]
\textbf{Énoncé.} Read the text and answer the questions.

\emph{Last Saturday, the Old Students' Association of Government High School Bafut organized a fund-raising gala to renovate the school library. More than two hundred guests attended the event. Local businesses sponsored the evening and donated prizes for a raffle. Thanks to the generosity of donors and volunteers, the association raised over one million francs CFA. The money will be used to buy new shelves, computers and textbooks for the students.}

\begin{enumerate}
\item What was the purpose of the fund-raising gala?
\item Who sponsored the event?
\item Find in the text: (a) a word meaning « person who gives money »; (b) a word meaning « person who works without being paid »; (c) an expression meaning « collected money ».
\item How will the money be used?
\end{enumerate}
\tcblower
\textbf{Solution détaillée.}
\begin{enumerate}
\item \textbf{Question 1.} Le mot \eng{purpose} demande le \emph{but}. Dans le texte : \eng{to renovate the school library} (l'infinitif de but introduit par \eng{to}). Réponse rédigée : \eng{The purpose of the fund-raising gala was to renovate the school library.}
\item \textbf{Question 2.} On cherche le sujet du verbe \eng{sponsored} : \eng{Local businesses sponsored the evening}. Réponse : \eng{Local businesses sponsored the event.}
\item \textbf{Question 3.} (a) \eng{donor} = « donateur » (famille de \eng{to donate} + suffixe \eng{-or} de l'agent) ; (b) \eng{volunteer} = « bénévole » ; (c) \eng{raised over one million francs} : l'expression \eng{to raise money/funds} signifie « collecter des fonds » — attention, ce n'est pas « lever » au sens physique.
\item \textbf{Question 4.} Structure passive du texte : \eng{The money will be used to buy...} Réponse : \eng{The money will be used to buy new shelves, computers and textbooks for the students.}
\end{enumerate}
\textbf{Conclusion méthodologique.} Au Bac, les réponses de compréhension doivent être des phrases complètes, construites à partir des mots du texte, jamais de simples mots isolés recopiés.
\end{exoresolu}

\courssub{Méthode 2 : Réemployer le lexique à l'oral et à l'écrit}

\begin{methode}[Produire un dialogue ou un paragraphe avec le vocabulaire du thème]
\begin{enumerate}
\item \textbf{Préparer une banque de structures.} Avant de parler ou d'écrire, listez cinq structures prêtes à l'emploi :
\begin{itemize}
\item \eng{We are organizing a fund-raising event to + verbe} (but) ;
\item \eng{We need volunteers to help us + verbe} ;
\item \eng{Would you like to contribute to our project?} ;
\item \eng{The money raised will be used to + verbe} ;
\item \eng{Thanks to your generosity, we managed to + verbe}.
\end{itemize}
\item \textbf{Respecter la chronologie d'un projet.} Un paragraphe de production suit l'ordre logique : \emph{problem} $\rightarrow$ \emph{project} $\rightarrow$ \emph{fund-raising activities} $\rightarrow$ \emph{result}. Utilisez les connecteurs \eng{first, then, moreover, finally, thanks to}.
\item \textbf{Varier les verbes.} Ne répétez pas \eng{to give} partout : alternez \eng{to donate, to contribute, to sponsor, to support, to fund}.
\item \textbf{À l'oral, utiliser les formules de politesse.} Pour solliciter : \eng{Could you help us...?}, \eng{We would be grateful if you could...}. Pour remercier : \eng{We really appreciate your support.}
\item \textbf{Se relire.} Vérifiez : sujet + verbe conjugué, prépositions (\eng{to contribute \textbf{to}}, \eng{to rely \textbf{on}}, \eng{responsible \textbf{for}}), et orthographe (\eng{fund-raising} avec un trait d'union).
\end{enumerate}
\end{methode}

\begin{exoresolu}[Production écrite guidée]
\textbf{Énoncé.} Your class wants to raise funds to buy a water pump for your village. Write a short paragraph (60--80 words) presenting the project. Use at least four words or expressions from the lesson: \eng{fund-raising, to raise funds, donation, volunteer, community, to contribute}.

\tcblower
\textbf{Solution détaillée.}

\emph{Plan.} Problème (pas d'eau potable) $\rightarrow$ projet (acheter une pompe) $\rightarrow$ actions de collecte $\rightarrow$ appel à contribution.

\emph{Paragraphe modèle :}

\eng{Our village has no clean drinking water, so the students of our class have decided to organize a fund-raising campaign to buy a water pump. We plan to raise funds through a charity football match and a raffle. Several volunteers have already promised to help us, and local traders have made generous donations. We invite every member of the community to contribute to this important project. Together, we can bring clean water to our village.}

« Notre village n'a pas d'eau potable, c'est pourquoi les élèves de notre classe ont décidé d'organiser une campagne de collecte de fonds pour acheter une pompe à eau. Nous prévoyons de collecter des fonds grâce à un match de football de charité et une tombola. Plusieurs bénévoles ont déjà promis de nous aider, et des commerçants locaux ont fait de généreux dons. Nous invitons chaque membre de la communauté à contribuer à ce projet important. Ensemble, nous pouvons apporter l'eau potable à notre village. »

\emph{Justifications.}
\begin{itemize}
\item \eng{to organize a fund-raising campaign} : collocation correcte ; \eng{fund-raising} est un adjectif composé avec trait d'union placé devant le nom.
\item \eng{to raise funds through...} : la préposition \eng{through} exprime le moyen (« par le biais de »).
\item \eng{have already promised} : present perfect, car l'action passée a un résultat visible maintenant.
\item \eng{have made generous donations} : collocation \eng{to make a donation} (jamais \eng{to do a donation}).
\item \eng{to contribute to this project} : préposition \eng{to} obligatoire après \eng{contribute}.
\end{itemize}
\textbf{Conclusion.} Un bon paragraphe de production au Bac = un plan clair + quatre ou cinq mots du thème correctement employés + des verbes conjugués sans faute.
\end{exoresolu}

\courssub{Méthode 3 : Éviter les faux amis et les calques du français}

\begin{methode}[Déjouer les pièges de traduction dans ce champ lexical]
\begin{enumerate}
\item \textbf{Identifier le faux ami.} Quand un mot anglais ressemble à un mot français, méfiez-vous : comparez toujours les deux sens avant d'écrire.
\item \textbf{Mémoriser les paires critiques du thème.}
\begin{center}
\begin{tabularx}{\linewidth}{|X|X|X|}
\hline
\rowcolor{popblueL} Français & English correct & Faux ami à éviter \\
\hline
collecter des fonds & to raise funds & \eng{to found} (= fonder) \\
\hline
une collecte de fonds & a fund-raising event & \eng{a collection of funds} (calque) \\
\hline
faire un don & to make a donation & \eng{to do a donation} \\
\hline
une œuvre de charité & a charity & \eng{a charity work} (calque) ; \eng{charity} seul suffit \\
\hline
bénévole (n.) & a volunteer & \eng{a voluntary} (adjectif seulement) \\
\hline
soutenir (financièrement) & to support / to sponsor & \eng{to support} ne prend jamais \eng{for} : \eng{to support a cause} \\
\hline
une salle des fêtes & a community hall & \eng{a party room} (calque) \\
\hline
\end{tabularx}
\end{center}
\item \textbf{Attention aux prépositions.} \eng{to contribute \textbf{to} a project}, \eng{to rely \textbf{on} donations}, \eng{to be responsible \textbf{for} the event}, \eng{to take part \textbf{in} a gala}.
\item \textbf{Attention à l'ordre des mots.} L'adjectif précède le nom : \eng{a generous donor} (et non \eng{a donor generous}). Les adverbes de fréquence se placent avant le verbe principal : \eng{They often organize raffles.}
\item \textbf{Vérifier au brouillon.} Soulignez chaque mot « suspect » (ressemblant au français) et demandez-vous : « Est-ce le mot que j'ai vu dans le cours ? »
\end{enumerate}
\end{methode}

\begin{exoresolu}[Correction de phrases fautives]
\textbf{Énoncé.} Each sentence contains one mistake (false friend, wrong preposition or word order). Correct it.

\begin{enumerate}
\item The association did a donation of books to the school.
\item Many voluntaries helped during the fund-raising gala.
\item We contributed for the construction of the borehole.
\item They organized a gala very successful last month.
\item The community depends of donations to pay the nurses.
\end{enumerate}
\tcblower
\textbf{Solution détaillée.}
\begin{enumerate}
\item \eng{did a donation} $\rightarrow$ \eng{\textbf{made} a donation}. Collocation obligatoire : \eng{to make a donation} ; \eng{to do} ne se combine pas avec \eng{donation}. Phrase correcte : \eng{The association made a donation of books to the school.}
\item \eng{voluntaries} $\rightarrow$ \eng{\textbf{volunteers}}. Faux ami : \eng{voluntary} est un adjectif (« volontaire, facultatif ») ; le nom de la personne est \eng{volunteer}. Phrase correcte : \eng{Many volunteers helped during the fund-raising gala.}
\item \eng{contributed for} $\rightarrow$ \eng{contributed \textbf{to}}. Calque du français « contribuer à » : en anglais, la préposition est \eng{to}. Phrase correcte : \eng{We contributed to the construction of the borehole.}
\item \eng{a gala very successful} $\rightarrow$ \eng{a \textbf{very successful gala}}. En anglais, l'adjectif (même accompagné d'un adverbe) se place \emph{avant} le nom, contrairement au français « un gala très réussi ». Phrase correcte : \eng{They organized a very successful gala last month.}
\item \eng{depends of} $\rightarrow$ \eng{depends \textbf{on}}. Calque de « dépendre de ». Phrase correcte : \eng{The community depends on donations to pay the nurses.}
\end{enumerate}
\textbf{Conclusion.} Trois réflexes sauvent des points : collocation (\eng{make a donation}), préposition (\eng{contribute to, depend on}), ordre des mots (adjectif avant le nom).
\end{exoresolu}

\courssub{Méthode 4 : Exploiter les mots de la même famille}

\begin{methode}[Former et employer les dérivés d'un même mot]
\begin{enumerate}
\item \textbf{Construire le tableau de famille.} Pour chaque mot du thème, cherchez le verbe, le nom d'action, le nom de personne et l'adjectif.
\begin{center}
\begin{tabular}{|l|l|l|l|}
\hline
\rowcolor{popblueL} Verbe & Nom (action) & Nom (personne) & Adjectif \\
\hline
to donate & donation & donor & --- \\
\hline
to volunteer & volunteering & volunteer & voluntary \\
\hline
to organize & organization & organizer & organized \\
\hline
to contribute & contribution & contributor & --- \\
\hline
to fund & funding / fund & funder & funded \\
\hline
to benefit & benefit & beneficiary & beneficial \\
\hline
\end{tabular}
\end{center}
\item \textbf{Identifier la nature grammaticale demandée.} Avant de transformer un mot, demandez-vous : « Ai-je besoin d'un verbe, d'un nom ou d'un adjectif dans cette phrase ? » Regardez l'entourage : après un article (\eng{a, the}) il faut un nom ; après un sujet, un verbe ; devant un nom, un adjectif.
\item \textbf{Attention à l'orthographe des dérivés.} \eng{to organize} $\rightarrow$ \eng{organization} (la graphie britannique \eng{organise/organisation} est aussi acceptée, mais soyez constant dans toute votre copie) ; \eng{benefit} $\rightarrow$ \eng{beneficiary} (« bé-né-fi-chia-ri »).
\item \textbf{Réemployer dans une phrase.} Pour chaque dérivé, écrivez une phrase : \eng{The beneficiaries of the project are the pupils of the primary school.} « Les bénéficiaires du projet sont les élèves de l'école primaire. »
\end{enumerate}
\end{methode}

\begin{exoresolu}[Formation de mots (word building)]
\textbf{Énoncé.} Complete each sentence with the correct form of the word in brackets.

\begin{enumerate}
\item The \ldots\ of the new classrooms will start next month. (construct)
\item Many \ldots\ gave money to buy drugs for the health centre. (donate)
\item The project is entirely \ldots\ by an international charity. (fund)
\item The \ldots\ of the gala thanked all the guests warmly. (organize)
\item Clean water is very \ldots\ to the whole community. (benefit)
\end{enumerate}
\tcblower
\textbf{Solution détaillée.}
\begin{enumerate}
\item \eng{\textbf{construction}}. Après l'article \eng{the}, il faut un nom : verbe \eng{to construct} $\rightarrow$ nom \eng{construction} (suffixe \eng{-tion}). Phrase : \eng{The construction of the new classrooms will start next month.}
\item \eng{\textbf{donors}}. Sujet de \eng{gave}, précédé de \eng{many} : il faut un nom de personne au pluriel : \eng{donor} + \eng{-s}. Phrase : \eng{Many donors gave money to buy drugs for the health centre.}
\item \eng{\textbf{funded}}. Structure passive : auxiliaire \eng{is} + participe passé \eng{funded} + agent introduit par \eng{by}. Phrase : \eng{The project is entirely funded by an international charity.}
\item \eng{\textbf{organizer}} (ou \eng{organizers}). Après \eng{the}, nom de personne : celui qui organise. Phrase : \eng{The organizer of the gala thanked all the guests warmly.}
\item \eng{\textbf{beneficial}}. Après le verbe \eng{is} et l'adverbe \eng{very}, il faut un adjectif : \eng{benefit} $\rightarrow$ \eng{beneficial} (« bénéfique »), suivi de la préposition \eng{to}. Phrase : \eng{Clean water is very beneficial to the whole community.}
\end{enumerate}
\textbf{Conclusion.} La méthode infaillible : observer l'entourage du mot manquant (article, auxiliaire, adverbe) pour déterminer la nature grammaticale, puis appliquer le bon suffixe.
\end{exoresolu}

\courssub{Méthode 5 : Comprendre un enregistrement ou un dialogue sur le thème (listening)}

\begin{methode}[Réussir la compréhension orale sur les collectes de fonds]
\begin{enumerate}
\item \textbf{Avant l'écoute, anticiper.} Lisez les questions et prédisez le vocabulaire que vous allez entendre : \eng{gala, raffle, sponsor, target, amount, borehole}.
\item \textbf{Première écoute : saisir le sens général.} Répondez mentalement à trois questions : \emph{Who? What? Why?} Ne vous bloquez pas sur un mot inconnu.
\item \textbf{Deuxième écoute : noter les détails.} Repérez les chiffres (montants, dates, nombres de participants) et les noms de personnes ou de lieux. Les chiffres sont presque toujours interrogés au Bac.
\item \textbf{Attention aux pièges sonores.} Distinguez \eng{thirteen} / \eng{thirty}, \eng{fifteen} / \eng{fifty} ; notez que \eng{can't} se prononce souvent « kènt » à l'américaine.
\item \textbf{Rédiger la réponse en phrase complète.} Recyclez les mots de la question : \eng{How much did they raise?} $\rightarrow$ \eng{They raised five hundred thousand francs.}
\end{enumerate}
\end{methode}

\begin{exoresolu}[Compréhension orale — script et questions]
\textbf{Énoncé.} Listen to the recording (script below) and answer the questions.

\emph{Script :} \eng{Good morning. This is Radio Bamenda. The Women's Support Group of Nkwen is pleased to announce its annual charity walk next Saturday. Participants will walk from the market square to the community hall. Our target is to raise two million francs to equip the health centre with new beds. Registration is free, but every participant is encouraged to find at least two sponsors. Refreshments will be offered to all walkers.}

\begin{enumerate}
\item Who is organizing the charity walk?
\item What is the target amount, and what is it for?
\item What must every participant try to do?
\end{enumerate}
\tcblower
\textbf{Solution détaillée.}
\begin{enumerate}
\item La première phrase du script donne le sujet : \eng{The Women's Support Group of Nkwen}. Réponse rédigée : \eng{The Women's Support Group of Nkwen is organizing the charity walk.}
\item Deux informations à combiner : le montant (\eng{two million francs}) et le but (introduit par \eng{to} + infinitif : \eng{to equip the health centre with new beds}). Réponse : \eng{Their target is to raise two million francs to equip the health centre with new beds.} Vocabulaire : \eng{target} = « objectif » ; \eng{to equip} = « équiper ».
\item Le script dit : \eng{every participant is encouraged to find at least two sponsors}. Réponse : \eng{Every participant is encouraged to find at least two sponsors.} Piège : \eng{Registration is free} ne signifie pas qu'il n'y a rien à faire — la consigne porte sur les \eng{sponsors}.
\end{enumerate}
\textbf{Conclusion.} En listening, chaque réponse se construit avec les mots du script ; notez les chiffres au brouillon dès la première écoute.
\end{exoresolu}

\begin{attention}[Les 5 erreurs qui coûtent des points au Bac]
\begin{enumerate}
\item \textbf{Faux ami : \eng{to raise} / \eng{to rise}.} \eng{To raise funds} (collecter) est transitif : il a toujours un complément. \eng{To rise} (s'élever) est intransitif. Écrire \eng{They rose funds} est une faute grave. Prononciation : \eng{raise} « réz », \eng{rise} « raïz ».
\item \textbf{Calque : \eng{to do a donation}.} On dit \eng{to \textbf{make} a donation} ou \eng{to \textbf{give} a donation}. De même : \eng{to make a contribution}, jamais \eng{to do a contribution}.
\item \textbf{Prépositions fautives.} Retenez : \eng{to contribute \textbf{to}}, \eng{to depend \textbf{on}}, \eng{to take part \textbf{in}}, \eng{to be responsible \textbf{for}}, \eng{to be interested \textbf{in} helping}. Une préposition calquée du français coûte un point à chaque occurrence.
\item \textbf{Ordre des mots : adjectif avant le nom.} Écrivez \eng{a successful fund-raising event}, \eng{generous donors}, jamais \eng{an event successful} ni \eng{donors generous}. L'adverbe de fréquence se place avant le verbe : \eng{They \textbf{often} organize charity walks.}
\item \textbf{Orthographe et forme des mots du thème.} \eng{fund-raising} (trait d'union), \eng{volunteer} (nom de personne ; \eng{voluntary} est l'adjectif), \eng{beneficiary} (et non \eng{beneficiaire}), \eng{organization} (ou graphie britannique \eng{organisation} — soyez constant), \eng{community} (deux \eng{m} : c-o-m-m-u-n-i-t-y).
\end{enumerate}
\end{attention}

\newpage

\coursec{Exercices}

\courssub{Je vérifie mes connaissances}

\begin{exercice}[QCM : la bonne collocation]
Pour chaque phrase, entoure la lettre de la seule réponse correcte.

\begin{enumerate}
\item Our association decided to \ldots\ a public appeal for the new health centre.\\
a) do \quad b) launch \quad c) rise \quad d) make
\item Many well-wishers \ldots\ a donation to the orphanage last Christmas.\\
a) made \quad b) did \quad c) launched \quad d) raised
\item The students helped to \ldots\ funds for the school library.\\
a) rise \quad b) lift \quad c) raise \quad d) grow
\item The \ldots\ of the charity concert will be used to equip the village borehole.\\
a) benefits \quad b) profits \quad c) proceeds \quad d) wages
\item A \ldots\ is a person who gives money or goods to help others.\\
a) beneficiary \quad b) donor \quad c) beggar \quad d) customer
\item The NGO distributed food to the \ldots\ families of the neighbourhood.\\
a) needy \quad b) necessary \quad c) needed \quad d) needing
\end{enumerate}
\end{exercice}

\begin{exercice}[Vrai ou faux ? Justifie ta réponse]
Indique si chaque affirmation est \textbf{true} ou \textbf{false}, puis justifie en une phrase en anglais.

\begin{enumerate}
\item A \eng{fundraiser} is an event organized to collect money for a cause.
\item The \eng{beneficiary} is the person who gives the money.
\item \eng{In kind} donations are donations of goods or services, not money.
\item A \eng{support group} is a business company that sells community services.
\item \eng{To volunteer} means to work without being paid, by free choice.
\end{enumerate}
\end{exercice}

\begin{exercice}[Vocabulaire : trouve le mot]
Complète chaque définition avec un mot de la leçon commençant par la lettre indiquée.

\begin{enumerate}
\item A deep hole dug to provide water to a village : b\ldots
\item A public request for money or help : an a\ldots
\item Money or goods given to help a person or an organization : a d\ldots
\item The sum of money an association hopes to collect : the t\ldots
\item A person who works for free to help others : a v\ldots
\item A person who gives money to a charity : a d\ldots
\end{enumerate}
\end{exercice}

\begin{exercice}[Conjugaison et formes à compléter]
Mets le verbe entre parenthèses à la forme qui convient (temps, voix ou forme nominale).

\begin{enumerate}
\item Last year, our club \ldots\ (raise) enough money to buy ten desks.
\item The new classrooms \ldots\ (equip) with computers next term.
\item So far, the support group \ldots\ (collect) two million francs.
\item The villagers thanked the \ldots\ (donate) warmly during the ceremony.
\item While the appeal \ldots\ (launch), many volunteers were distributing leaflets.
\item Her \ldots\ (contribute) to the project was highly appreciated.
\end{enumerate}
\end{exercice}

\courssub{Je m'entraîne}

\begin{exercice}[Traduction]
Traduis les phrases suivantes en anglais, en utilisant le vocabulaire de la leçon.

\begin{enumerate}
\item Notre association a lancé un appel aux dons pour construire un forage dans le village.
\item Les bénévoles ont collecté des fonds pour les familles nécessiteuses.
\item Les bienfaiteurs peuvent faire un don en argent ou en nature.
\item Les recettes du concert de charité serviront à équiper le centre de santé.
\item Le groupe de soutien a remercié tous les donateurs pour leur généreuse contribution.
\end{enumerate}
\end{exercice}

\begin{exercice}[Faux amis : corrige les calques]
Chaque phrase contient un calque du français. Corrige l'erreur et réécris la phrase correctement.

\begin{enumerate}
\item I assisted to the fund-raising gala last Saturday.
\item The beneficiaires thanked their benefactors for the gifts.
\item The association made a collect of clothes for the orphans.
\item We must actualize the list of donors before the meeting.
\item The money was destined to the construction of a borehole.
\item The volunteers demanded a contribution from each parent.
\item Our school realized a big profit during the charity fair.
\item The mayor assisted the ceremony in person.
\end{enumerate}
\end{exercice}

\begin{attention}[Les faux amis à retenir]
\begin{center}
\begin{tabularx}{\linewidth}{|X|X|X|}
\hline
\rowcolor{popblueL} Français & Piège & Forme correcte \\
\hline
assister à (une cérémonie) & \eng{assist} & \eng{attend} \\
\hline
actualiser & \eng{actualize} & \eng{update} \\
\hline
demander (poliment) & \eng{demand} & \eng{ask for} \\
\hline
réaliser (un profit) & \eng{realize} & \eng{make / earn} \\
\hline
destiné à & \eng{destined to} & \eng{intended for} \\
\hline
une collecte & \eng{a collect} & \eng{a collection / a fund-raising drive} \\
\hline
\end{tabularx}
\end{center}
\end{attention}

\begin{textebox}[A village gets clean water — texte à trous]
The women of Mankon have launched an \ldots\ to well-wishers at home and abroad. Their \ldots\ is to raise five million francs before the end of the year. Thanks to a generous \ldots\ from a local businessman and the hard work of twenty \ldots\ , the village has already collected half of the money. The \ldots\ of last month's charity concert will pay for the pipes. Once the \ldots\ is completed, it will be \ldots\ with a modern hand pump. Every \ldots\ , however small, will help the \ldots\ families who still walk long distances to fetch water. The local \ldots\ for widows has promised to maintain the pump.
\end{textebox}

\begin{exercice}[Texte à trous type Bac]
Complète l'article ci-dessus avec dix mots de la liste ci-dessous. Chaque mot ne peut être utilisé qu'une seule fois.

\begin{center}
\begin{tabular}{|l|l|l|l|l|}
\hline
\rowcolor{popblueL} donation & volunteers & proceeds & appeal & borehole \\
\hline
\rowcolor{popblueL} needy & contribution & target & equipped & support group \\
\hline
\end{tabular}
\end{center}
\end{exercice}

\begin{exercice}[Appariements : expressions idiomatiques]
Associe chaque expression anglaise (1--10) à son équivalent français (a--j).

\begin{center}
\begin{tabularx}{\linewidth}{|l X|}
\hline
\rowcolor{popblueL} \multicolumn{2}{|l|}{\textbf{English expressions}} \\
\hline
1 & to lend a helping hand \\
2 & every little helps \\
3 & to give back to the community \\
4 & to dig deep into one's pockets \\
5 & charity begins at home \\
6 & to do one's bit \\
7 & a drop in the ocean \\
8 & to pull together \\
9 & to be in need \\
10 & actions speak louder than words \\
\hline
\end{tabularx}
\end{center}

\begin{center}
\begin{tabularx}{\linewidth}{|l X|}
\hline
\rowcolor{popgreenL} \multicolumn{2}{|l|}{\textbf{Équivalents français}} \\
\hline
a & être dans le besoin \\
b & donner un coup de main \\
c & les actes valent mieux que les paroles \\
d & se donner à fond (financièrement) \\
e & les petits ruisseaux font les grandes rivières \\
f & une goutte d'eau dans l'océan \\
g & rendre à la communauté ce qu'elle nous a donné \\
h & la charité bien ordonnée commence par soi-même \\
i & apporter sa pierre à l'édifice \\
j & unir ses efforts \\
\hline
\end{tabularx}
\end{center}
\end{exercice}

\courssub{Je me prépare au Bac}

\begin{textebox}[Harvest thanksgiving in Muyuka]
Every year in November, the people of Muyuka, a small town in the South West Region, celebrate their harvest thanksgiving. The event is more than a religious ceremony: it is the biggest fund-raising activity of the year. Weeks before the great day, the village development committee launches an appeal to all sons and daughters of the town, including those living in Douala, Yaoundé and even abroad.

On thanksgiving Sunday, after the church service, families bring gifts in kind: bags of rice, bunches of plantains, goats and even building materials. These gifts are auctioned in front of the church, and the proceeds are carefully recorded by the treasurer. Last year, the committee set a target of three million francs to renovate the community health centre. To everyone's surprise, the auction raised four million. The extra money was used to buy benches for the primary school.

The beneficiaries of these projects are not only the needy. The whole community profits from better services. As the traditional ruler declared during the ceremony, “When we pull together, no project is too big for Muyuka.” Volunteers have already promised to do their bit again next year.
\end{textebox}

\begin{exercice}[Compréhension écrite type Bac]
Lis le texte ci-dessus, puis réponds aux questions.

\textbf{A. Vrai ou faux ?} Justifie chaque réponse par une citation du texte.
\begin{enumerate}
\item The harvest thanksgiving is only a religious ceremony.
\item People living outside Muyuka are not concerned by the appeal.
\item The committee collected more money than expected.
\end{enumerate}

\textbf{B. QCM.} Choisis la bonne réponse.
\begin{enumerate}
\item The gifts are \ldots\ after the church service.\\
a) burned \quad b) auctioned \quad c) hidden \quad d) eaten
\item The extra money was used to \ldots\\
a) pay the treasurer \quad b) build a church \quad c) buy school benches \quad d) repair the roads
\end{enumerate}

\textbf{C. Réponses courtes.} Réponds en phrases complètes, avec tes propres mots.
\begin{enumerate}
\item What kinds of gifts do families bring on thanksgiving Sunday?
\item According to the traditional ruler, what makes big projects possible?
\end{enumerate}

\textbf{D. Langue.}
\begin{enumerate}
\item Trouve dans le texte : (a) un synonyme de \eng{aim}, (b) un synonyme de \eng{money collected from sales}, (c) le contraire de \eng{rich people}.
\item Mets à la voix passive : \eng{The committee recorded the proceeds carefully.}
\end{enumerate}
\end{exercice}

\begin{textebox}[Texte à traduire]
Chaque année, notre école organise une journée de collecte de fonds pour aider les élèves nécessiteux. L'année dernière, les bénévoles ont lancé un appel aux parents et aux bienfaiteurs de la ville. Grâce à leurs généreuses contributions, nous avons dépassé notre objectif : les recettes de la vente de charité ont permis d'acheter des livres et d'équiper la bibliothèque. Le proviseur a remercié tous les donateurs et a promis que l'argent serait utilisé uniquement pour les projets de la communauté scolaire.
\end{textebox}

\begin{exercice}[Traduction type Bac]
Traduis le passage ci-dessus en anglais.
\end{exercice}

\begin{exercice}[Rédaction type Bac : lettre d'appel aux dons]
\textbf{Sujet.} Your school needs a new library. As president of the students' support group, write an appeal letter to well-wishers asking for donations in cash or in kind. Explain the project, state your target, and say how the proceeds will be used.

\textbf{Consignes :}
\begin{itemize}
\item Write between \textbf{100 and 120 words}.
\item Use the layout of a formal letter (address, date, salutation, closing formula).
\item Use at least \textbf{six words or expressions} from the lesson (appeal, donation, in cash or in kind, target, proceeds, well-wishers, contribution, equipped, needy, lend a helping hand\ldots).
\item Sign as \eng{Manka'a Njoh}, Government High School Buea.
\end{itemize}
\end{exercice}

\begin{methode}[Rédiger une lettre d'appel aux dons]
\begin{enumerate}
\item \textbf{En-tête} : ton adresse en haut à droite, puis la date (\eng{Government High School, P.O. Box 63, Buea, 15th March 2025}).
\item \textbf{Salutation} : \eng{Dear Sir or Madam,} si tu ne connais pas le destinataire.
\item \textbf{Introduction} : présente-toi et annonce l'objet de la lettre (\eng{I am writing to launch an appeal for\ldots}).
\item \textbf{Développement} : explique le projet, indique le \eng{target} (objectif financier), précise que les dons peuvent être \eng{in cash or in kind}, et dis comment les \eng{proceeds} seront utilisés.
\item \textbf{Conclusion} : remercie par avance (\eng{Every contribution, however small, will be highly appreciated.}).
\item \textbf{Formule de politesse} : \eng{Yours faithfully,} (avec \eng{Dear Sir or Madam}), puis ta signature.
\end{enumerate}
\end{methode}



\newpage

\coursec{Corrigés des exercices}

\begin{corrige}[Exercice 1 — QCM : la bonne collocation]
\begin{enumerate}
\item \textbf{b) launch} — la collocation correcte est \eng{to \cle{launch} an \cle{appeal}} (lancer un appel). \eng{Do} et \eng{make} ne se combinent pas avec \eng{appeal} ; \eng{rise} (intransitif) signifie « s'élever ».
\item \textbf{a) made} — on dit \eng{to \cle{make} a \cle{donation}} (faire un don). \eng{To raise} s'emploie avec \eng{funds} ou \eng{money}, pas avec \eng{a donation}.
\item \textbf{c) raise} — \eng{to \cle{raise} \cle{funds}} (collecter des fonds). Attention : \eng{rise} (rose, risen) est intransitif et ne prend jamais de complément d'objet.
\item \textbf{c) proceeds} — \eng{the \cle{proceeds} of an event} = les recettes d'une manifestation. \eng{Profits} évoque un bénéfice commercial ; \eng{benefits} = avantages ; \eng{wages} = salaires.
\item \textbf{b) donor} — le \eng{\cle{donor}} (donateur) donne ; le \eng{\cle{beneficiary}} (bénéficiaire) reçoit.
\item \textbf{a) needy} — \eng{\cle{needy} families} = les familles nécessiteuses. \eng{Necessary} = nécessaire (faux ami) ; \eng{needed} et \eng{needing} sont des formes verbales inadaptées ici.
\end{enumerate}
\textbf{Barème indicatif :} 1 point par réponse correcte, soit 6 points.
\end{corrige}

\begin{corrige}[Exercice 2 — Vrai ou faux ?]
\begin{enumerate}
\item \textbf{True.} \eng{A \cle{fundraiser} is an event organized to raise money for a good cause.}
\item \textbf{False.} \eng{The \cle{beneficiary} is the person who receives the money or the help; the \cle{donor} is the one who gives it.}
\item \textbf{True.} \eng{\cle{In-kind} donations consist of goods or services, for example bags of rice or free \cle{labour}, not money.}
\item \textbf{False.} \eng{A \cle{support group} is a group of people who help one another or support a cause; it is not a business company.}
\item \textbf{True.} \eng{A \cle{volunteer} works by free choice and without payment.}
\end{enumerate}
\textbf{Barème indicatif :} 1 point par réponse (0,5 pour vrai/faux + 0,5 pour la justification en anglais correct), soit 5 points.
\end{corrige}

\begin{corrige}[Exercice 3 — Vocabulaire : trouve le mot]
\begin{enumerate}
\item \eng{\cle{borehole}} (un forage) — prononciation : « bor-ohl ».
\item \eng{\cle{appeal}} (un appel) — collocation : \eng{to launch an appeal}.
\item \eng{\cle{donation}} (un don) — collocation : \eng{to make a donation}.
\item \eng{\cle{target}} (l'objectif) — collocations : \eng{to set a target}, \eng{to reach a target}, \eng{to exceed a target}.
\item \eng{\cle{volunteer}} (un bénévole) — attention à l'orthographe : \textbf{-eer} final.
\item \eng{\cle{donor}} (un donateur) — attention à l'orthographe : \textbf{-or} final, comme \eng{benefactor}, \eng{actor}.
\end{enumerate}
\textbf{Barème indicatif :} 1 point par mot correct (orthographe exigée), soit 6 points.
\end{corrige}

\begin{corrige}[Exercice 4 — Conjugaison et formes]
\begin{enumerate}
\item \eng{raised} — prétérit simple, marqué par \eng{last year} ; verbe régulier : \emph{raise, raised, raised}.
\item \eng{will be equipped} — futur passif (\eng{will be} + participe passé) : les salles « seront équipées » ; indice : \eng{next term}.
\item \eng{has collected} — present perfect, imposé par \eng{so far} (jusqu'à présent) ; sujet singulier : \eng{has}.
\item \eng{donors} — nom de personne au pluriel (les donateurs) ; dérivation : \eng{donate} $\rightarrow$ \eng{donor}.
\item \eng{was being launched} — prétérit continu passif : action en arrière-plan introduite par \eng{while} ; \eng{was being} + participe passé.
\item \eng{contribution} — nom après l'adjectif possessif \eng{her} ; dérivation : \eng{contribute} $\rightarrow$ \eng{contribution}.
\end{enumerate}
\textbf{Barème indicatif :} 1 point par forme correcte, soit 6 points.
\end{corrige}

\begin{corrige}[Exercice 5 — Traduction]
\begin{enumerate}
\item \eng{Our association launched an appeal for donations to build a borehole in the village.}
\item \eng{The volunteers raised funds for the needy families.}
\item \eng{Well-wishers can make a donation in cash or in kind.}
\item \eng{The proceeds of the charity concert will be used to equip the health centre.}
\item \eng{The support group thanked all the donors for their generous contribution.}
\end{enumerate}
\textbf{Points de vigilance :}
\begin{itemize}
\item Prétérit pour les actions passées datées (phrases 1, 2, 5).
\item Futur avec \eng{will} dans la phrase 4 (passif \eng{will be used}).
\item Ne pas traduire « nécessiteuses » par \eng{necessary} (faux ami) $\rightarrow$ \eng{needy}.
\item « Bienfaiteurs » se dit \eng{well-wishers} ou \eng{benefactors}, jamais \eng{well-doers}.
\item \eng{Proceeds} est toujours pluriel dans ce sens (le verbe est au pluriel : \eng{were}, \eng{are}).
\end{itemize}
\textbf{Barème indicatif :} 2 points par phrase (1 pour le lexique, 1 pour la grammaire), soit 10 points.
\end{corrige}

\begin{corrige}[Exercice 6 — Faux amis : corrige les calques]
\begin{enumerate}
\item \eng{I \cle{attended} the fund-raising gala last Saturday.} — \eng{to attend} = assister à (être présent) ; \eng{to assist} = aider.
\item \eng{The \cle{beneficiaries} thanked their benefactors for the gifts.} — orthographe anglaise : \eng{beneficiaries} (avec \emph{-i-} et \emph{-ies} au pluriel).
\item \eng{The association organized a \cle{collection} of clothes for the orphans.} — « une collecte » se dit \eng{a collection} ou \eng{a fund-raising drive} ; \eng{a collect} n'existe pas comme nom.
\item \eng{We must \cle{update} the list of donors before the meeting.} — \eng{to update} = actualiser ; \eng{actual} (adj.) signifie « réel ».
\item \eng{The money was \cle{intended for} the construction of a borehole.} — \eng{intended for} = destiné à ; \eng{destined to} évoque le destin/fatalité.
\item \eng{The volunteers \cle{asked} each parent for a contribution.} — \eng{to demand} est trop autoritaire (exiger) ; \eng{to ask for} = demander poliment.
\item \eng{Our school \cle{made} a big profit during the charity fair.} — collocation : \eng{to make a profit} ; \eng{to realise} (UK) / \eng{realize} (US) = se rendre compte.
\item \eng{The mayor \cle{attended} the ceremony in person.} — même correction qu'à la phrase 1 : \eng{to attend}.
\end{enumerate}
\textbf{Barème indicatif :} 1 point par phrase corrigée, soit 8 points.
\end{corrige}

\begin{corrige}[Exercice 7 — Texte à trous type Bac]
Dans l'ordre du texte :
\begin{enumerate}
\item \eng{appeal} — \eng{launched an appeal} (lancer un appel).
\item \eng{target} — l'objectif : cinq millions de francs.
\item \eng{donation} — \eng{a generous donation from a local businessman}.
\item \eng{volunteers} — vingt bénévoles qui travaillent dur.
\item \eng{proceeds} — les recettes du concert de charité.
\item \eng{borehole} — le forage, une fois achevé (\eng{completed}).
\item \eng{equipped} — \eng{it will be equipped with a modern hand pump} (passif futur).
\item \eng{contribution} — \eng{every contribution, however small}.
\item \eng{needy} — les familles nécessiteuses.
\item \eng{support group} — le groupe de soutien aux veuves.
\end{enumerate}
\textbf{Barème indicatif :} 1 point par mot bien placé, soit 10 points. Chaque mot n'est utilisé qu'une fois : vérifiez qu'aucun mot n'apparaît deux fois dans votre copie.
\end{corrige}

\begin{corrige}[Exercice 8 — Appariements : expressions idiomatiques]
1--b ; 2--e ; 3--g ; 4--d ; 5--h ; 6--i ; 7--f ; 8--j ; 9--a ; 10--c.

\begin{center}
\begin{tabularx}{\linewidth}{|X|X|}
\hline
\rowcolor{popblueL} Expression anglaise & Équivalent français \\
\hline
\eng{to \cle{lend a helping hand}} & donner un coup de main \\
\hline
\eng{\cle{every little helps}} & les petits ruisseaux font les grandes rivières \\
\hline
\eng{to \cle{give back to the community}} & rendre à la communauté ce qu'elle nous a donné \\
\hline
\eng{to \cle{dig deep into one's pockets}} & se donner à fond (financièrement) \\
\hline
\eng{\cle{charity begins at home}} & la charité bien ordonnée commence par soi-même \\
\hline
\eng{to \cle{do one's bit}} & apporter sa pierre à l'édifice \\
\hline
\eng{a \cle{drop in the ocean}} & une goutte d'eau dans l'océan \\
\hline
\eng{to \cle{pull together}} & unir ses efforts \\
\hline
\eng{to \cle{be in need}} & être dans le besoin \\
\hline
\eng{\cle{actions speak louder than words}} & les actes valent mieux que les paroles \\
\hline
\end{tabularx}
\end{center}
\textbf{Barème indicatif :} 1 point par appariement correct, soit 10 points.
\end{corrige}

\begin{corrige}[Exercice 9 — Compréhension écrite type Bac]
\textbf{A. Vrai ou faux ?}
\begin{enumerate}
\item \textbf{False.} “The event is more than a religious ceremony: it is the biggest fund-raising activity of the year.”
\item \textbf{False.} The appeal is launched “to all sons and daughters of the town, including those living in Douala, Yaoundé and even abroad.”
\item \textbf{True.} “The committee set a target of three million francs… the auction raised four million.”
\end{enumerate}

\textbf{B. QCM.}
\begin{enumerate}
\item \textbf{b) auctioned} — “These gifts are auctioned in front of the church.”
\item \textbf{c) buy school benches} — “The extra money was used to buy benches for the primary school.”
\end{enumerate}

\textbf{C. Réponses courtes.}
\begin{enumerate}
\item \eng{They bring gifts in kind, such as bags of rice, bunches of plantains, goats and building materials.}
\item \eng{According to him, big projects are possible when people pull together, that is, when they unite their efforts.}
\end{enumerate}

\textbf{D. Langue.}
\begin{enumerate}
\item (a) \eng{target} ; (b) \eng{proceeds} ; (c) \eng{the needy}.
\item \eng{The proceeds were carefully recorded by the committee.} — passif au prétérit : \eng{were} + participe passé ; attention, \eng{proceeds} est pluriel, d'où \eng{were}.
\end{enumerate}

\textbf{Barème indicatif (sur 20) :} A : 3 $\times$ 2 points (réponse + citation) = 6 ; B : 2 $\times$ 2 points = 4 ; C : 2 $\times$ 2 points = 4 ; D : 3 $\times$ 1 point + 3 points = 6.

\textbf{Points de vigilance :} citez le texte entre guillemets anglais “ … ” pour les justifications ; répondez en phrases complètes en C ; ne recopiez pas le texte mot à mot en C (« with your own words »).
\end{corrige}

\begin{corrige}[Exercice 10 — Traduction type Bac]
\textbf{Proposition de traduction :}

\eng{Every year, our school organises a fund-raising day to help needy students. Last year, the volunteers launched an appeal to the parents and to the well-wishers of the town. Thanks to their generous contributions, we exceeded our target: the proceeds of the charity sale made it possible to buy books and to equip the library. The principal thanked all the donors and promised that the money would be used only for the projects of the school community.}

\textbf{Points de vigilance :}
\begin{itemize}
\item \eng{needy students} (élèves nécessiteux) — jamais \eng{necessary students} (faux ami).
\item \eng{to launch an appeal to somebody} — préposition \eng{to}.
\item \eng{to exceed our target} (dépasser notre objectif) — on accepte aussi \eng{to go beyond our target}.
\item \eng{the proceeds of the charity sale} — \eng{proceeds} est toujours pluriel dans ce sens.
\item \eng{the principal} (le proviseur) — attention au faux ami : \eng{principal} = directeur d'école ; \eng{principle} = principe.
\item Discours rapporté au passé : \eng{promised that the money would be used} (concordance des temps : \eng{will} $\rightarrow$ \eng{would}).
\item Orthographe UK : \eng{organises} (et non \eng{organizes}).
\end{itemize}

\textbf{Barème indicatif (sur 10) :} 5 points pour le lexique du thème (fund-raising day, appeal, well-wishers, contributions, target, proceeds, donors…), 3 points pour la grammaire (temps, passif, concordance), 2 points pour l'exactitude globale du sens.
\end{corrige}

\begin{corrige}[Exercice 11 — Rédaction type Bac : lettre d'appel aux dons]
\textbf{Modèle de rédaction :}

\begin{quote}
Government High School\\
P.O. Box 63, Buea\\
15th March 2025

Dear Sir or Madam,

I am Manka'a Njoh, president of the students' support group of Government High School Buea. I am writing to launch an appeal to all well-wishers for our new school library.

Our target is to raise two million francs before the end of the school year. Your donation, in cash or in kind — books, shelves or computers — will lend a helping hand to hundreds of needy students. The proceeds of this appeal will be used to buy books and to equip the reading room.

Every contribution, however small, will be highly appreciated. Remember: every little helps.

Yours faithfully,\\
Manka'a Njoh
\end{quote}

Ce modèle compte environ 110 mots et réemploie huit expressions de la leçon : \eng{support group}, \eng{appeal}, \eng{well-wishers}, \eng{target}, \eng{donation}, \eng{in cash or in kind}, \eng{proceeds}, \eng{contribution}.

\textbf{Barème indicatif (sur 20) :}
\begin{itemize}
\item \textbf{Respect de la tâche} (lettre d'appel, projet expliqué, objectif chiffré, emploi des recettes annoncé) : 5 points.
\item \textbf{Mise en page de la lettre formelle} (adresse, date, salutation, formule finale, signature) : 4 points.
\item \textbf{Réemploi du lexique de la leçon} (au moins six expressions) : 5 points.
\item \textbf{Correction grammaticale et orthographique} : 4 points.
\item \textbf{Respect de la longueur} (100--120 mots) : 2 points.
\end{itemize}

\textbf{Points de vigilance :}
\begin{itemize}
\item \eng{Yours faithfully} s'accorde avec \eng{Dear Sir or Madam} (et \eng{Yours sincerely} avec un destinataire nommé).
\item Ne pas écrire « I am writing you » mais \eng{I am writing to you}.
\item Éviter le calque « to sensitize well-wishers » au sens de « lancer un appel » ; utiliser \eng{to launch an appeal to}.
\item Vérifier l'orthographe de \eng{volunteers}, \eng{beneficiaries}, \eng{contribution}.
\item Utiliser l'orthographe britannique (\eng{organises}, \eng{realise}, \eng{labour}).
\end{itemize}
\end{corrige}

\newpage

\coursec{Fiche bilan}

\begin{popfigure}
\begin{tikzpicture}[mindmap, grow cyclic, every node/.style={concept, font=\small},
root concept/.append style={concept color=popblue, font=\bfseries\small, text=white},
level 1/.append style={level distance=3.6cm, sibling angle=51},
level 1 concept/.append style={font=\scriptsize}]
\node[root concept]{Fund-raising \& Community Support}
child[concept color=poporange]{ node[align=center]{Fund-raising activities\\ \scriptsize charity, donation, raffle, gala} }
child[concept color=popgreen]{ node[align=center]{Support groups\\ \scriptsize volunteers, NGO, self-help} }
child[concept color=poppurple]{ node[align=center]{Community resources\\ \scriptsize borehole, health centre, hall} }
child[concept color=popgold]{ node[align=center]{Collocations\\ \scriptsize raise funds, make a donation} }
child[concept color=poppink]{ node[align=center]{Word families\\ \scriptsize donate, donation, donor} }
child[concept color=popmuted]{ node[align=center]{False friends\\ \scriptsize assist, eventually, actually} }
child[concept color=popblue!60]{ node[align=center]{Idioms\\ \scriptsize give a helping hand, chip in} };
\end{tikzpicture}
\legende{Carte mentale de la leçon : le lexique de la collecte de fonds et du soutien communautaire.}
\end{popfigure}

\begin{bilan}
\textbf{1. Lexique clé : fund-raising activities}
\begin{itemize}
\item \eng{fund-raising} (n.) : collecte de fonds ; \eng{a fund-raiser} : un organisateur / une activité de collecte
\item \eng{to raise funds / money} : collecter des fonds ; \eng{a charity} : une œuvre caritative
\item \eng{a donation} : un don ; \eng{a donor} : un donateur ; \eng{to donate} : faire un don
\item \eng{a raffle} : une tombola ; \eng{a charity gala} : un gala de bienfaisance ; \eng{a sponsored walk} : une marche sponsorisée
\item \eng{a contribution} : une contribution ; \eng{to contribute} : contribuer ; \eng{proceeds} (n. pl.) : les recettes
\item \eng{a fundraising target} : un objectif de collecte ; \eng{to reach a target} : atteindre un objectif
\end{itemize}
\textbf{2. Lexique clé : support groups and community resources}
\begin{itemize}
\item \eng{a support group} : un groupe de soutien ; \eng{a self-help group} : un groupe d'entraide
\item \eng{a volunteer} : un bénévole ; \eng{voluntary work} : le bénévolat ; \eng{an NGO} (non-governmental organisation) : une ONG
\item \eng{a community} : une communauté ; \eng{community development} : le développement communautaire
\item \eng{resources} : les ressources ; \eng{services} : les services ; \eng{facilities} : les infrastructures
\item \eng{a borehole} : un forage ; \eng{a health centre} : un centre de santé ; \eng{a community hall} : une salle des fêtes communautaire
\item \eng{to equip a school} : équiper une école ; \eng{to provide clean water} : fournir de l'eau potable
\end{itemize}
\textbf{3. Collocations et expressions à connaître par cœur}
\begin{center}
\begin{tabularx}{\linewidth}{|X|X|}\hline
\rowcolor{popblueL} \textbf{English} & \textbf{Français} \\ \hline
to raise funds for a project & collecter des fonds pour un projet \\ \hline
to make a donation to a charity & faire un don à une œuvre caritative \\ \hline
to organise a fund-raising event & organiser une manifestation de collecte \\ \hline
to give a helping hand & donner un coup de main \\ \hline
to chip in & cotiser, participer financièrement \\ \hline
to rely on donations & compter sur les dons \\ \hline
to volunteer one's time & offrir bénévolement son temps \\ \hline
to benefit the whole community & profiter à toute la communauté \\ \hline
\end{tabularx}
\end{center}
\textbf{4. Mots de la même famille}
\begin{itemize}
\item \eng{to donate} (v.) $\rightarrow$ \eng{a donation} (n.) $\rightarrow$ \eng{a donor} (personne)
\item \eng{to contribute} (v.) $\rightarrow$ \eng{a contribution} (n.) $\rightarrow$ \eng{a contributor} (personne)
\item \eng{to volunteer} (v. et n.) $\rightarrow$ \eng{voluntary} (adj.) $\rightarrow$ \eng{voluntarily} (adv.)
\item \eng{to help} (v. et n.) $\rightarrow$ \eng{helpful} (adj.) $\rightarrow$ \eng{helpless} (adj. contraire)
\item \eng{to fund} (v.) $\rightarrow$ \eng{a fund / funds} (n.) $\rightarrow$ \eng{fund-raising} (n. et adj.)
\end{itemize}
\textbf{5. Faux amis et pièges}
\begin{itemize}
\item \eng{to assist} = aider (formel), pas « assister à » : \emph{to attend a meeting} = assister à une réunion.
\item \eng{actually} = en fait, réellement ; « actuellement » = \emph{currently, at present}.
\item \eng{eventually} = finalement ; « éventuellement » = \emph{possibly, if necessary}.
\item \eng{a benefit} = un avantage ; « un bénéfice » (argent) = \emph{a profit}.
\item On dit \eng{to raise money} (collecter), jamais \emph{to rise money} ; \emph{rise} est intransitif.
\end{itemize}
\textbf{6. Méthode express}
\begin{enumerate}
\item Repérer le thème du texte (qui collecte ? pour quoi ? pour qui ?).
\item Souligner les collocations verbe + nom (\eng{raise funds}, \eng{make a donation}).
\item Réemployer 5 à 8 mots du champ lexical dans sa propre production orale ou écrite.
\item Vérifier les faux amis avant de rendre la copie.
\end{enumerate}
\end{bilan}

\begin{aretenir}[Checklist avant le Bac]
\begin{itemize}
\item[$\square$] Je connais le vocabulaire de la collecte de fonds : \eng{fund-raising, donation, donor, raffle, proceeds}.
\item[$\square$] Je connais le vocabulaire des groupes de soutien : \eng{support group, volunteer, NGO, self-help group}.
\item[$\square$] Je sais nommer les ressources communautaires : \eng{borehole, health centre, community hall, facilities}.
\item[$\square$] J'emploie correctement les collocations : \eng{to raise funds, to make a donation, to reach a target}.
\item[$\square$] Je distingue \eng{raise} (transitif) et \eng{rise} (intransitif).
\item[$\square$] Je connais les familles de mots : \eng{donate / donation / donor}, \eng{volunteer / voluntary}.
\item[$\square$] J'évite les faux amis : \eng{assist, actually, eventually, benefit}.
\item[$\square$] Je sais utiliser les expressions idiomatiques : \eng{to give a helping hand, to chip in}.
\item[$\square$] Je prononce correctement \eng{charity} (« tcha-ri-ti ») et \eng{volunteer} (« vo-lon-tier »).
\item[$\square$] Je sais rédiger un court paragraphe décrivant une activité de collecte de fonds dans ma communauté.
\end{itemize}
\end{aretenir}

\findecours

\end{document}
